% concatenated from Shapes.tex
\documentclass[12pt]{amsart}
\usepackage{latexsym}
\usepackage[utf8]{inputenc}
\usepackage{amsmath}
\usepackage{amsthm}
\usepackage{amsfonts}
\usepackage{amssymb}
\usepackage{epsfig}
\usepackage{nicefrac}
\usepackage[all]{xy}
\usepackage{graphicx}
\usepackage{graphics} 
\graphicspath{ {./images/} }
\usepackage{subcaption}
\usepackage{enumerate}
\DeclareGraphicsRule{.tif}{png}{.png}{`convert #1 `basename #1 .tif`.png}
\DeclareMathOperator{\sign}{sign}
\usepackage[T1]{fontenc}
\usepackage{subcaption}
\usepackage[font=small,skip=0pt]{caption}
\usepackage{tikz}
\usepackage{subfig}
 
%\newcommand*\mycirc[1]{%
%   \begin{tikzpicture}
%     \node[draw,circle,inner sep=1pt] {#1};
%   \end{tikzpicture}}
%
% \usepackage{graphicx}
%\usepackage{caption}
%\usepackage{mwe}    
%\renewcommand\thefigure{Main-\arabic{figure}}
%\DeclareCaptionLabelFormat{simplab}{Figure~#2}
%\DeclareCaptionFormat{nocap}{}
%\DeclareCaptionSubType*{figure}
%\renewcommand\thesubfigure{\arabic{subfigure}}
%\captionsetup[subfigure]{%
%            font={normalsize},
%            labelfont={bf},
%            labelformat=simplab,
%            labelsep=colon,
%            singlelinecheck=trued
%            position=bottom
%            }%
%


%\renewcommand{\thefootnote}{{\arabic{footnote})}}

\newtheorem{theorem}{Theorem}[section]
\newtheorem{corollary}[theorem]{Corollary}
\newtheorem{proposition}[theorem]{Proposition}
\newtheorem{lemma}[theorem]{Lemma}
\newtheorem{claim}[theorem]{Claim}
\newtheorem{question}[theorem]{Question}
\newtheorem{conjecture}[theorem]{Conjecture}


\theoremstyle{definition}
\newtheorem{assumption}[theorem]{Assumption}
\newtheorem{remark}[theorem]{Remark}
\newtheorem{example}[theorem]{Example}
\newtheorem{condition}[theorem]{Condition}
\newtheorem{definition}[theorem]{Definition}








%\newcommand{\eqdef}{\;{:=}\;}
%\newcommand\PSL{\operatorname{PSL}}
%\newcommand\coker{\operatorname{coker}}
%\newcommand\supt{\operatorname{supt}}
%\newcommand\const{\operatorname{const}}
%\newcommand\Index{\operatorname{Index}}
%\newcommand\Symp{\operatorname{Symp}}
%\newcommand\Ham{\operatorname{Ham}}
%\newcommand\Def{{\overset {\rm {def}}{\ =\ }}}

\def\bigmid{\ \rule[-3.5mm]{0.1mm}{9mm}\ }

\newcommand{\DD}{{\mathbb D}}
\newcommand{\CC}{{\mathbb C}}
\newcommand{\RR}{{\mathbb R}}
\newcommand{\ZZ}{{\mathbb Z}}
\newcommand{\NN}{{\mathbb N}}
\newcommand{\TT}{{\mathbb T}}



\begin{document}


\title[Constraints on Lagrangians, and the shape invariant]{New constraints on Lagrangian embeddings, and the shape invariant}
\date{\today}
\author{Richard Hind}
\address{RH: Department of Mathematics\\
University of Notre Dame\\
255 Hurley\\
Notre Dame, IN 46556, USA.}
\author{Ely Kerman}
\address{EK: Department of Mathematics\\
University of Illinois at Urbana-Champaign\\
1409 West Green Street\\
Urbana, IL 61801, USA.}
\thanks{Both authors are supported by grants from the Simons Foundation}


\begin{abstract}
For a large class of toric domains in $\RR^4$,  we determine which product Lagrangian tori can be mapped into the domain by a Hamiltonian diffeomorphism. In other words, we compute the Hamiltonian shape invariant of these toric domains, as defined in \cite{hizh}. The argument relies on new intersection results for product Lagrangian tori in symplectic polydisks. For Hamiltonian diffeomorphisms which map certain Lagrangian product tori back into the polydisk, we establish intersections between the images and a one-parameter family of product Lagrangian tori that includes (is based at) the original torus. For symplectic polydisks with area ratios less than two, we strengthen this to establish intersections between the Hamiltonian images and the original Lagrangian torus. As a soft complement to these intersection results we also present an embedding construction which demonstrates that this intersection rigidity vanishes when the one-parameter family of product Lagrangian tori is replaced by a natural packing by Lagrangian tori.

\end{abstract}

\maketitle

\section{Introduction}

Let $L$ be a Lagrangian submanifold of $\RR^{2n}$ equipped with its standard symplectic structure. Given an open subset $U \subset \RR^{2n}$ it is natural to ask whether there exists a Hamiltonian diffeomorphism of $\RR^{2n}$ that maps $L$ into $U$. One can then ask whether such maps must also satisfy special constraints. For example, if $L$ is contained in $U$, then  the 
set  $$\{ \phi(L) \mid \phi \in \mathrm{Ham}(\RR^{2n}), \phi(L) \subset U\}$$ is nonempty and one can ask if any pair of elements must intersect. Even in simple settings, these pairwise intersections may not be guaranteed. They may also only be a small part of a richer intersection story.

Consider, for example, the standard open disk $D(b) \subset \RR^2$ of area $b$ and the standard (Lagrangian) circle $L(s) $ enclosing area $s<b$. Suppose that $b > 2s$. Then  $L(s)$ and its Hamiltonian images in $D(b)$ need not intersect. On the other hand, the two intersection statements below always hold and follow immediately from simple area considerations.

\begin{theorem}\label{thm:2a} If $\phi$ is a Hamiltonian diffeomorphism of $\RR^2$ such that $\phi(L(s)) \subset D(b)$, then $\phi(L(s))$ must intersect the family $$\bigcup_{t\in [s,b-s]}L(t).$$ 
%Moreover, if $b\leq 2s$, then $\phi(L(s))$ must intersect $L(s)$.
\end{theorem}

%The role of the family  may also be played by a finite collection of Lagrangians.

\begin{theorem}\label{thm:2b}
    Let $k \geq 2$ be the unique integer such that $ks<b \leq (k+1)s$ and suppose that  $$\bigsqcup_{j=1}^k\Lambda_j$$ is a disjoint union of embedded closed curves $\Lambda_j$ in $D(b)$ that each bound a domain of area  $s$.  If $\phi$ is a Hamiltonian diffeomorphism of $\RR^2$ that maps $L(s)$ into $D(b)$, then $\phi(L(s))$ must intersect at least one of the $\Lambda_j$.
\end{theorem}
\begin{figure}[!h]
    \centering
    \includegraphics[width=0.85\linewidth]{2in2pie-1.pdf}
    \vspace{-6.5cm}
    \caption{Any Hamiltonian image of $L(s)$ in $D(b)$ must intersect both $\bigcup_{t\in [s,b-s]}L(t)$ and $\bigsqcup_{j=1}^k\Lambda_j$.}
    \label{fig:2in2}
\end{figure}


\bigskip








In this paper we consider if and how the rigidity statements in Theorems \ref{thm:2a} and \ref{thm:2b} might extend to  four dimensions. Identifying $\RR^4$ with $\CC^2$, we consider product Lagrangian tori,  
\begin{equation*} 
L(r,s) =\{ (z_1,z_2) \in \CC^2 \mid \pi |z_1|^2 = r, \, \pi |z_2|^2 = s \},
%X_{\{ (r,s) \}} = \mu^{-1}(r,s).
\end{equation*}
% the simple one-parameter families of tori,
% \begin{equation} \label{prod-torus}
% L(r,[a,b)) =  \bigcup_{t \in [a,b)} L(r, t) = S(r) \times \bigcup_{t \in [a,b)}S(t),
% %X_{\{ (r,s) \}} = \mu^{-1}(r,s).
% \end{equation}
and symplectic polydisks,
\begin{equation*} \label{poly}
P(a,b) = \{(z_1,z_2) \in \CC^2 \mid \pi |z_1|^2 <a, \, \pi |z_2|^2 <b \}.
\end{equation*}

\subsubsection{Generalizing Theorem \ref{thm:2a}}\label{secgen}
Our main result is the following extension of Theorem \ref{thm:2a} to dimension four.

\begin{theorem}\label{lagint}
Suppose that  $1/2\leq r<1$ and $1 \leq s <b/2$. If $\phi$ is a Hamiltonian diffeomorphism of $\RR^4$ that maps $L(r,s)$ into $P(1,b)$, then $\phi(L(r,s))$ intersects the codimension one submanifold
\begin{equation} 
\bigcup_{t \in [s,b-s]} L(r, t).
\end{equation}
%Moreover, if $b<2s$, then $\phi(L(r,s))$ intersects $L(r,s)$.
\end{theorem}


The first three constraints on $r$ and $s$ in Theorem \ref{lagint} are each necessary. If $r$ is at least $1$, then Theorem 3 of \cite{ho} implies that there is no Hamiltonian diffeomorphism of $\RR^4$ that maps $L(r,s)$ into any polydisk of the form $P(1,b)$. If $r$ is less than $1/2$, then there are Hamiltonian diffeomorphisms of $D(1)$ that displace $D(r)$ from itself, and any such map can be used to construct a Hamiltonian diffeomorphism of $\RR^4$ that maps $L(r,s)$ into $P(1,b) \smallsetminus \bigcup_{t \in [s,b-s]} L(r, t).$ If $s$ is less than $1$, and greater than $r$, then $L(s,r)$, which is Hamiltonian isotopic to $L(r,s)$, is contained in $P(1,b)$ and is disjoint from the family $\bigcup_{t \in [s,b-s]} L(r, t)$. 

% The theorem itself implies that the last constraint, $s<b$, is also necessary. In particular, Theorem \ref{lagint} implies that if $s \geq b$ (and $1/2 \le r<1$), then no Hamiltonian diffeomorphism can map $L(r,s)$ into $P(1,b)$. 



\medskip

% Pushing the proof of Theorem \ref{lagint} a bit further, we establish the following refinement.

% \begin{theorem}\label{lagint2}
% Suppose that $1/2\leq r<1$ and $1 \leq s <b$.
% % Suppose that $r$ and $s$ satisfy
% % \begin{equation}\label{rs} \frac{1}{2} \leq r <1 \leq s.\end{equation}
% If $\phi$ is a Hamiltonian diffeomorphism of $\RR^4$ that maps $L(r,s)$ into $P(1,b)$, then $\phi(L(r,s))$ intersects 
% \begin{equation} \label{prod-torus}
% \bigcup_{t \in [s,b-s]} L(r, t).
% \end{equation}
% % If $b > 2s$, then every element of $\mathcal{L}(L(r,s), P(1,b))$ must intersect the set 
% % \begin{equation} \label{prod-torus}
% % L(r,[s,b)) =  \bigcup_{t \in [s,b-s)} L(r, t).
% %\end{equation}
% \end{theorem}

In the limit $b\to 2s$, Theorem \ref{lagint} yields the following Lagrangian intersection result which appears to be new. 

\begin{corollary}\label{newlag}
Suppose that $1/2\leq r<1$ and $1 \leq s <b \le 2s$. If $\phi$ is a Hamiltonian diffeomorphism of $\RR^4$ that maps $L(r,s)$ into $P(1,b)$, then $$L(r,s) \cap \phi(L(r,s)) \neq \emptyset.$$
\end{corollary}


As described below in Section \ref{sec:2lag}, one can use the methods developed here to prove a more general result. We show, in Corollary \ref{newlagv2}, that any two Hamiltonian images of $L(r,s)$ in $P(1,b)$ must intersect.


\begin{remark}
It would be interesting to know if these intersections can be detected by some refinement of Lagrangian Floer theory. In the same spirit, it would also be interesting to know whether there is a topological lower bound for the number of intersection points here.
\end{remark}



\begin{remark}
Versions of Corollary \ref{newlag} fail in dimension greater than $4$. More precisely, it follows from \cite[Theorem A]{chek} that when $\alpha>1$ is irrational, the product torus $L=L(1, \alpha, 2)$ is Hamiltonian isotopic, in $\RR^6$, to infinitely many distinct product tori lying arbitrarily close to $L$. This is discussed in section 2 of the article \cite{atallah}, as well as in \cite{brsh}. 
\end{remark}

\begin{remark} If one compactifies $P(1,b)$ to $S^2 \times S^2$ with a nonmonotone symplectic form, then the analogue of Corollary \ref{newlag} fails almost completely. Corollary 1.10 of \cite{brki} implies that there is a set of split Lagrangian tori in $S^2 \times S^2$, of full measure, such that for every $L$ in the set every Weinstein neighborhood of $L$ contains a disjoint Hamiltonian image of $L$. For more discussion on the $S^2 \times S^2$ case see Remark \ref{brendelschmitz}.
    
\end{remark}


\medskip

In the limit $b\to \infty$, Theorem \ref{lagint} yields  the following Lagrangian refinement of Gromov's nonsqueezing theorem from \cite{gr}.
\begin{corollary}\label{binf}
   Suppose that $r<1 \le s$. If there is a Hamiltonian diffeomorphism of $\RR^4$ that maps $L(r,s)$ into 
   $$(D(1) \times \RR^2) \smallsetminus \bigcup_{t \in [s,\infty)} L(r, t),$$ then $r <1/2$.
\end{corollary}



\bigskip


\subsubsection{Generalizing Theorem \ref{thm:2b}} The setting of Theorem \ref{thm:2b} can be extended to dimension four in the following natural way. Given a collection of disjoint embedded closed curves  
$$\{\Lambda_1, \dots, \Lambda_{k}\} \subset D(b)$$
each bounding area $s$, as in Theorem \ref{thm:2b}, consider the $k$ disjoint Lagrangian tori
$$L_j = L(r) \times \Lambda_j$$ in $P(2r,b).$ Given a Hamiltonian diffeomorphism $\phi$ of $\RR^4$ with $\phi(L_1) \subset P(2r,b)$, one can ask whether $\phi(L_1)$ must intersect at least one of the $L_j$. 
By continuity, when asking this question, one can assume that  $$b < (k+1)s.$$ 


In this setting, the work of Polterovich and Shelukhin from \cite{ps} yields a local rigidity theorem in the spirit of Theorem \ref{thm:2b}. The relationship between $r$ and the positive number
 \begin{align}\label{eq:g}
    \mathfrak{o} = (k+1)s-b,
\end{align}
plays a crucial role.
\begin{theorem}(Theorem C, \cite{ps})\label{thm:ps}
If $r < \mathfrak{o}/(k-1),$ then there is no Hamiltonian diffeomorphism supported in $P(2r,b)$ that displaces $L(r,s)$ from all of the $L_j.$
\end{theorem}

In the present work, we show, by construction, that even this local rigidity disappears when $r$ is sufficiently large.

\begin{theorem}\label{thm:hm}
If $k$ is even, then for every $r> 6\mathfrak{o}/k $ there is a Hamiltonian diffeomorphism of $P(2r,b)$ that displaces $L(r,s)$ from all of the  $L_j.$
\end{theorem}

This construction also extends to the limit $b \to \infty$.

\begin{corollary}\label{cor:disc}
    Let $\Lambda \colon S^1 \times \NN \to \RR^2$ be a proper embedding such that each $\Lambda_j = \Lambda(S^1 \times \{j\})$ bounds a disk of area $s$. Set $L_j = L(1/2) \times \Lambda_j$. There exists a Hamiltonian diffeomorphism $\phi$ of the cylinder  $D(1) \times \RR^2 $ such that $\phi(L(1/2,s)) \cap L_j=\emptyset$ for all $j \in \NN$.
\end{corollary}

The following related questions remain unresolved.

\begin{question}
    For which values of $r$ between $\mathfrak{o}/(k-1)$  and $6\mathfrak{o}/k$, if any,  does the rigidity statement of Theorem \ref{thm:ps} hold?
\end{question}

\begin{question}
    If one allows all Hamiltonian diffeomorphisms of $\RR^4$, not just those supported in $P(2r,b)$, does the rigidity statement  of Theorem \ref{thm:ps} hold for any positive value of $r$?
\end{question}

\begin{remark}\label{brendelschmitz}
    If one compactifies $P(2r, b)$ to $S^2 \times S^2$, then the gap between rigidity and nonrigidity shrinks to a point. Theorem C of \cite{ps} asserts that for  $r < \mathfrak{o}/(k-1)$ there is no Hamiltonian diffeomorphism of $S^2 \times S^2$ that displaces $L(r,s)$ from all of the $L_j.$ On the other hand, by Corollary 1.16 of \cite{brki}, for almost every $s<b$ satisfying $r > \mathfrak{o}/(k-1)$,
    %, then for almost all $L(r,s)$ (in the copy of $S^2 \times S^2$ obtained by compactifying $P(2r,b)$)
    there is a infinite sequence of Hamiltonian diffeomorphisms $\phi_j$ of $S^2 \times S^2$ such that the images $\phi_j(L(r,s))$ are pairwise disjoint.

    There is more flexibility if we also consider Lagrangian tori $L(x,y) \subset P(2r,b)$ with $x \neq r$. Assuming our compactification is not monotone, that is, $2r \neq b$, recent work of Schmitz, \cite[Theorem B]{schmitz} says that, given $\varepsilon>0$, there exists a (single) Hamiltonian diffeomorphism $\psi$ of $S^2 \times S^2$ such that, for almost every $(x,y) \in \{0<x<2r, \, 0< y<b, \, |x-r| > \varepsilon \}$, the images $\psi^j(L(x,y))$ are pairwise disjoint.
\end{remark}



% %\begin{remark}
%  %If $b >2s$, then there are Hamiltonian diffeomorphisms supported in $P(1,b)$ that displace $L(r,s)$ from  $\cup_{t \in[s,b-s)} L(r,t).$
% %\end{remark}
% \subsection{Related results and questions}

% \subsubsection{Restricting to $\mathrm{Ham}(P(1,b))$.}

% % \bigskip

% % \noindent {\bf Restricting to $\mathrm{Ham}(P(1,b))$.}
% The primary novel feature of Theorem \ref{lagint} and Theorem \ref{thm:2s} is that they hold for all Hamiltonian diffeomorphisms  of $\RR^4$ that map $L(r,s)$ into $P(1,b)$, not just Hamiltonian diffeomorphisms of (supported in) $P(1,b)$. For this smaller class of maps some of these intersection results follow from previous work. 

% The work of Entov and Polterovich from \cite{ep} implies the following weaker version of Theorem \ref{thm:2s}.  
% \begin{theorem}(\cite{ep})
% Suppose $r$ and $s$ satisfy \eqref{rs}. If $\varphi$ is a Hamiltonian diffeomorphism  of $P(1,2s)$, then $\varphi(L(r,s))$ must intersect $L(r,s)$.
% \end{theorem}


% The work of Polterovich and Shelukhin in \cite{ps}, which involves intersections among finite collections of Lagrangian tori (see below), implies the following weaker version of Theorem \ref{lagint}.

% \begin{theorem}(\cite{ps})
% Suppose $r$ and $s$ satisfy \eqref{rs}. If $b$ is sufficiently large\footnote{Here $b$ must be chosen to be large enough for Theorem $C$ of \cite{ps} to hold for $a=1/2b$.}, then for every Hamiltonian diffeomorphism $\varphi$ of $P(1,b)$ the image $\varphi(L(r,s))$ intersects the family $\bigcup_{t \in [s,b)} L(r, t).$
% \end{theorem}




% \begin{proof}

% Arguing by contradiction, we assume that for all large $b$ there is a Hamiltonian diffeomorphism $\phi$ of $P(1,b)$ such that $\phi(L(r,s))$ is disjoint from $L(r, [s,b))$. Choose $b>0$ large enough so that Theorem $C$ of \cite{ps} holds for $a=1/2b$. Consider the map  $\phi_b$ defined by $$\phi_b(p) =\frac{1}{\sqrt{b}} \, \phi\left(\sqrt{b} p\right).$$ This is Hamiltonian diffeomorphism 
% of $ P(1/b,1)$, and $\phi_b(L(r/b,s/b))$ is disjoint from 
% the family of Lagrangian tori $L(r/b,[s/b,1)).$

% Now compactify $P(1/b,1)$ to $S^2(2a) \times S^2(1)$. Since $\phi_b$ is a Hamiltonian diffeomorphism 
% of $ P(1/b,1)$ it has compact support in $ P(1/b,1)$ and extends to a unique Hamiltonian diffeomorphism of $S^2(2a) \times S^2(1)$. We then have $\phi_b(L(r/b,s/b))$ is disjoint from
% $$E(2a) \times \hat{B}((b-s)/b),$$
% where $E(2a)$ is the equator in $S^2(2a)$ and $\hat{B}((b-s)/b)$ is the complement of the standard open disc of area $s/b$ in $S^2(1)$. As we now show, the existence of such a Hamiltonian diffeomorphism of $S^2(2a) \times S^2(1)$ is forbidden by the results of \cite{ps}.

% To see this, let $K$ be the unique integer satisfying \[K< b/s \le K+1.\] There are $K$ disjoint discs of area $s/b$ in $S^2(1)$ such that one of them is the boundary of $\hat{B}((b-s)/b)$ and the others are contained in the interior of $\hat{B}((b-s)/b)$. Hence there are $K$ disjoint Lagrangian tori in $S^2(2a) \times S^2(1)$, including $L(r/b,s/b)$, that are contained in $$E(2a) \times \hat{B}((b-s)/b) \subset S^2(2a) \times S^2(1)$$. By Theorem C of \cite{ps}, for any  Hamiltonian diffeomorphism $\psi$ of $S^2(2a) \times S^2(1)$ the image $\psi(L(r/b,s/b)$ must intersect one of these $K$ tori and hence must intersect  $E(2a) \times \hat{B}((b-s)/b)$. This contradicts the existence of the map $\phi_b$, as claimed.
% \end{proof}

% \subsubsection{Finite intersections}

% It is natural to ask if the one-parameter family of Lagrangian tori in Theorem \ref{lagint} can be replaced by a finite (or discrete) subcollection of Lagrangian tori. Here we describe the simplest version of this question and show that the answer is complicated even when one restricts the problem to Hamiltonian diffeomorphisms of $P(1,b)$. We discuss two partial answers in this setting: a positive answer in the form of a rigidity theorem  from \cite{ps}, and a second, negative, answer which follows from a construction presented here. The complete answer to this question is still unknown and merits further study.



% To pose the relevant question, we begin with an elementary model in dimension two. For the open disc $D(b)$ and any $s <b$ there is a unique positive integer $k$ such that \begin{equation}\label{bs}
%   ks<b\leq (k+1)s .
% \end{equation}
% Let  $\{\Lambda_j\}$, for $j=1, \dots k$, be a collection of disjoint embedded closed curves in $D(b)$ that each bound area $s$. The 
% %We may also assume that $\Lambda_1 = S(s).$ 
%  following two dimensional rigidity statement follows immediately from \eqref{bs}.

% \medskip
% \noindent{\emph {For any Hamiltonian diffeomorphism $\phi$ of $\RR^2$ that maps $S(s)$ into $D(b)$, the image $\phi(S(s))$ must intersect one of the $\Lambda_j$.}}
% \medskip

% Now consider a stabilized version of this picture. Given the closed curves $\Lambda_j \subset D(b)$ as above, consider the $k$ disjoint Lagrangian tori
% $$L_j = S(r) \times \Lambda_j$$ in the polydisk  $D(2r) \times D(b)= P(2r, b).$ Note that for every Hamiltonian diffeomorphism  $\phi$ of $\RR^2$ with $\phi(S(r)) \subset D(2r)$ we have $\phi(S(r)) \cap S(r) \neq \emptyset$.  The following question arises.


% \begin{question}
% If $\phi$ is  Hamiltonian diffeomorphism of $\RR^4$ that maps $L(r,s)$ into $P(2r,b)$, must the image $\phi(L(r,s))$ intersect one of the $L_j$?
% \end{question}

% When one restricts this question to Hamiltonian diffeomorphisms with support in $P(2r,b)$, then the answer depends on the size of $r$ relative to the overlap \begin{align}\label{eq:g}
%     \mathfrak{o} = (k+1)s-b.
% \end{align}

% The work of Polterovich and Shelukhin in \cite{ps} yields the following rigidity theorem.
% \begin{theorem}(Theorem C, \cite{ps})
% If $r < \mathfrak{o}/(k+1),$ then there is no Hamiltonian diffeomorphism supported in $P(2r,b)$ that displaces $L(r,s)$ from all of the $L_j.$
% \end{theorem}

% On the other hand, we show, by example, that this rigidity disappears when $r$ is sufficiently large.

% \begin{theorem}\label{thm:hm}
% If $k$ is even, then for every $r> 6\mathfrak{o}/k $ there is a Hamiltonian diffeomorphism of $P(2r,b)$ that displaces $L(r,s)$ from all of the  $L_j.$
% \end{theorem}

% The discrete version of Corollary \ref{binf} also fails in this setting.

% \begin{corollary}\label{cor:disc}
%     Let $\Lambda \colon S^1 \times \NN \to \RR^2$ be a proper embedding such that each $\Lambda_j = \Lambda(S^1 \times \{j\})$ bounds a disc of area $s$. Set $L_j = S(\frac12) \times \Lambda_j$.There exists a Hamiltonian diffeomorphism $\phi$ of $Z(1)$ such that $\phi(L(\frac12,s)) \cap L_j=\emptyset$ for all $j \in \NN$.
% \end{corollary}




\subsection{The Hamiltonian shape invariant}
The original motivation for Theorem \ref{lagint} concerns the Hamiltonian shape invariant  defined  by Hind and Zhang in \cite{hizh}. The Hamiltonian shape of a subset $X$ of $\RR^4$ is the set 
\begin{align*} 
{\rm Sh}(X) & : = \left\{(r, s) \, \big| \, r \le s, \,  L(r, s) \hookrightarrow X \right\} \subset \RR_{>0}^2,
\end{align*}
where the notation $L(r, s) \hookrightarrow X$ implies that there exists a Hamiltonian diffeomorphism $\Phi$ of $\RR^4$ with $\Phi(L(r,s)) \subset X$. This set-valued symplectic invariant of $X$ has the following properties:\\
\begin{enumerate}
    \item ${\rm Sh}(cX) = c^2 {\rm Sh}(X), \text{ for all } c>0$,\\
    \item $X \hookrightarrow Y \implies {\rm Sh}(X)\subset {\rm Sh}(Y)$,\\
    \item ${\rm Sh}(\cup_i X_i) \supset \cup_i {\rm Sh}(X_i)$,\\
    \item ${\rm Sh}(\cap_i X_i) \subset \cap_i {\rm Sh}(X_i).$\\
\end{enumerate}
It can thus be viewed as a set-valued symplectic capacity. Like all capacities, the Hamiltonian shape is highly nontrivial to compute. A natural starting class of examples to consider are toric domains. 

Recall that a 
% Theorem \ref{lagint} and Theorem \ref{bs} will allow us to compute the Hamiltonian shape of broad classes of toric domains.
{\em toric domain} in $\RR^4$ is a subset of the form $X_{\Omega} = \mu^{-1} (\Omega)$ where $\Omega$ is a subset of $\RR^2_{\ge 0}$
and $\mu: \RR^4=\CC^2 \to \RR^2_{\ge 0}$ is the standard moment map defined by $$\mu(z_1, z_2) = (\pi|z_1|^2, \pi |z_2|^2).$$ 
% For example, the symplectic polydisk $P(a,b)$ is the toric domain $X_{p(a,b)}$ where $p(a,b)= [0,a)\times [0,b)$, and the symplectic ellipsoid  $E(a,b)$ is the toric domain $X_{e(a,b)}$ where $$e(a,b) = \{(r,s) \in \RR^2_{\ge 0} \mid s< b-rb/a \}$$ 
For each point $(r,s)$ in the interior of $\Omega$, $\rm{int}(\Omega)$, the inverse image $\mu^{-1}(r,s)$ is the Lagrangian torus $L(r,s) \subset X_{\Omega}$. Hence, the set $\rm{int}(\Omega) \cap \{ r \leq s\}$ is contained in ${\rm Sh} (X_{\Omega})$. This direct connection between $\Omega$ and ${\rm Sh}(X_\Omega)$ inspires the hope that there might be a simple 
formula, or algorithm, for computing the Hamiltonian shape of $X_{\Omega}$ directly from $\Omega$.

This has been achieved for some important standard toric domains. In \cite{ho}, Hind and Opshstein compute the Hamiltonian shapes of balls and symplectic polydisks. In \cite{hizh}, Hind and Zhang compute the Hamiltonian shape of symplectic ellipsoids 
\begin{equation*} 
E(a,b) = \mu^{-1}(\{(r,s) \in \RR^2_{\ge 0} \mid x/a+y/b <1 \})
\end{equation*}
with $\frac{b}{a} \in \NN_{\geq 2}.$
Theorem \ref{lagint} and Corollary \ref{newlag} allow us to compute the Hamiltonian shape of two general families of toric domains.


\subsubsection{Information loss} Before stating our results, we recall two essential observations from  \cite{ho} concerning the ways in which information about $\Omega$ may become hidden in ${\rm Sh}(X_{\Omega})$. The first source of information loss in the passage from $\Omega$ to ${\rm Sh}(X_{\Omega})$, already accounted for in the definition of the shape, is due to a basic symmetry. Since the tori $L(r,s)$ and $L(s,r)$ are Hamiltonian isotopic, if one of them symplectically embeds in a domain, then so does the other.
Let $\sigma:(r,s) \mapsto (s,r)$ be the reflection in the diagonal in $\RR^2_{>0}$.
It follows that 
${\rm Sh}(X_{\Omega})$ contains the set $\Omega^+ \cap \{r \leq s\}$ where $$ \Omega^+= \Omega \cup \sigma \Omega$$
%\{ (r,s) \mid \{(r,s),\, (s,r)\} \cap \Omega \neq \emptyset \}.$$ Some properties of $\Omega$ can be lost in the transition to $\Omega^+$.  
This transition from $\Omega$ to $\Omega^+$ can obscure features of $\Omega$. In particular, for distinct subsets $\Omega_1$ and $\Omega_2$ we can have $\Omega_1^+ = \Omega_2^+$, see Figure \ref{fig:blind1}.

\begin{figure}
    \centering
    \includegraphics[width=0.8\linewidth]{Omega+.pdf}
    \vspace{-7.5cm}
    \caption{For these distinct regions $\Omega_1^+ = \Omega_2^+ =\Omega_3^+ $.}
    \label{fig:blind1}
\end{figure}


\bigskip

A second, more subtle, reason why some features  of $\Omega$ are not visible in ${\rm Sh}(X_{\Omega})$ is a consequence of the following Lagrangian embedding theorem of Hind and Opshtein.

\begin{theorem}(\cite{ho} \cite[Theorem 5.2]{hizh2})
For every $x>2$ and every $\epsilon >0$ there is a Hamiltonian diffeomorphism $\Phi$ of $\RR^4$ such that 
$\Phi(L(1,x))$ is contained in the toric domain $X_{(1+\epsilon)Q}$, where $Q$ is the convex hull of the set of points $\{(0,0), (2,0), (0,2), (1,2)\}$.
\end{theorem}

\begin{remark}
    The set $X_Q$ is the closure of $P(2,2) \cap E(2,4).$
\end{remark}




% The following natural question has been much studied.

% \begin{question} \label{qfibre}
% For which interior points $(r,s) \in \Omega$ is the fibre torus $\mu^{-1}(r,s)$ nondisplaceable by Hamiltonian diffeomorphisms of $X_{\Omega}.$
% \end{question}

% Theorem \ref{lagint} yields a set of examples relevant to the following natural generalization of Question \ref{qfibre}.

% \begin{question}
% For which subsets $U \subset \Omega$ is  $\mu^{-1}(U)$ nondisplaceable by Hamiltonian diffeomorphisms of $\RR^4$ that preserve $X_{\Omega}.$
% \end{question}

% For $\Omega= [0,1]\times \RR_{\ge 0}$ the corresponding toric domain is the symplectic cylinder $Z(1) := D(1) \times \RR^2$. Given $(r,s) \in [0,1]\times \RR_{\ge 0}$ let $$U_{r,s} = \{(r,t) \mid t\geq s\}.$$ 
% In this setting, Theorem \ref{lagint} implies the following.


% \begin{corollary}
%     If  $\frac{1}{2} \leq r <1$ and $s \ge 1$, then the subset $\mu^{-1}(U_{r,s}) \subset Z(1)$ can not be displaced by any Hamiltonian diffeomorphism of the cylinder $Z(1).$
% \end{corollary}

% Note that each subset $\mu^{-1}(U_{r,s}) \subset Z(1)$ has zero volume and is comprised entirely of displaceable fibres. 

An immediate implication of this result is that 
${\rm Sh}(X_{\Omega})$ contains the unbounded set $$\{0\leq r < m(\Omega)\} \cap \{r\leq s\} $$ where the width of the infinite strip is given by $$m(\Omega) = \sup \{ c \mid cQ \subset \Omega \, \, \mathrm{or} \, \, cQ \subset \sigma \Omega \}.$$ Hence, no part of $\Omega$ that lies in the infinite strip $\{0\leq r < m(\Omega)\}$ is visible in ${\rm Sh}(X_{\Omega})$.

\bigskip

In all examples where the shape of $X_{\Omega}$ is known, there is nothing else to the story, and one has
\begin{equation}\label{eq:shape}
    {\rm Sh}(X_{\Omega}) = \Bigl(\Omega^+  \cup \{0\leq r < m(\Omega)\} \Bigr)\cap \{r\leq s\}.
\end{equation} This equality is  established for balls and symplectic polydisks in \cite{ho}, and for symplectic ellipsoids $E(a,b)$ with $\frac{b}{a} \in \NN_{\geq 2}$ in \cite{hizh}.
Here we show that \eqref{eq:shape} holds much more generally. 


\subsubsection{New shape computations}

We will focus on subsets $\Omega \subset \RR^2_{\ge 0}$ that contain the unit square $[0,1)\times [0,1)$ and are contained in the unit strip $[0,1) \times \RR_{\ge 0}$. In this case, we have $\Omega^+ =\Omega$, $m(\Omega) =1/2$ and hence 
\begin{equation}\label{starter}
    {\rm Sh}(X_{\Omega}) \supset \Bigl(\Omega \cup \{0\leq r < 1/2\} \Bigr)\cap \{r\leq s\}.
\end{equation}
We will show that these sets are equal for two broad classes of subsets.

Theorem \ref{lagint} implies the following.

\begin{theorem}\label{main} Let $f: [0,1] \to [1, \infty]$ be any function. For the corresponding graphical subset $$\Omega(f)=\{ (r,s) \mid 0\leq r < 1, 0\leq s < f(r)\}$$ we have 
\begin{align*}
\label{fshape}
   {\rm Sh}(X_{\Omega(f)}) = \Bigl(\Omega(f)  \cup \{0\leq r < 1/2\} \Bigr)\cap \{r\leq s\}. 
\end{align*}
% In fact, \eqref{eq:shape} holds for $X_{\Omega{(f) \cup \Delta}}$ for any
% $$\Delta \subset \{ (r,s) \mid 1\leq r < f(s), 1 \leq s < f(r)\}.$$

\end{theorem}


\begin{proof}
 For $0< r\leq s$, it follows from Proposition 2.3 of \cite{chsc} that the displacement energy of $L(r,s)$ is equal to $r$. The displacement energy of the symplectic cylinder $D(1) \times \RR^2 = \mu^{-1}([0,1) \times \RR_{\ge 0})$ is equal to $1$. Since the displacement energy and shape are monotone and $X_{\Omega(f)} \subset D(1) \times \RR^2$,  if $(r,s)$ is in ${\rm Sh}(X_{\Omega(f)})$ we must have $r\leq 1$. It suffices to show that if $(r,s)$ is in ${\rm Sh}(X_{\Omega(f)})$ and $1/2 \leq r \leq 1$ then $s < f(r).$ Equivalently, it suffices to show that for any $(r,s)$ with $r \in [1/2,1]$ and $s \ge f(r)$ there is no Hamiltonian diffeomorphism of $\RR^4$ that maps $L(r,s)$ into $X_{\Omega(f)}$. Choosing any $b > 2f(r)$, this is an immediate implication of Theorem \ref{lagint}.
\end{proof}

\begin{figure}[!h]
    \centering
    \includegraphics[width=0.8\linewidth]{Shape1-2.pdf}
    \vspace{-7.5cm}
    \caption{Equality \eqref{eq:shape} holds for $X_{\Omega(f)}$.}
    \label{fig:shape1}
\end{figure}



\begin{remark} Consider two functions $f,g: [0,1] \to [1, \infty]$. Theorem \ref{main} implies that if the toric domain $X_{\Omega(f)}$ and $X_{\Omega(g)}$ are symplectomorphic, then the restrictions of the functions $f$ and $g$ to $[1/2,1]$ must be identical. A {\em rigidity of profile} result of a similar spirit has recently been asserted by Hutchings, \cite{hutch-talk}, for generic strictly convex toric domains in $\RR^4$. 
% Let $\overline{C}^0[1/2,\,1]$ be the vector space of continuous functions $g$ on $[1/2,\, 1]$ such that $g(1/2)=0$ and $\int_{\frac12}^1 g=0.$ Let $\overline{B}^0[1/2,\,1]$ be the closed unit ball in $\overline{C}^0[1/2,\,1]$ with respect to the $sup$ norm. Consider the map $g \mapsto X_{\Omega(f)}$ where 
% \[
%   f(\tau) =
%   \begin{cases}
%     2, & \text{for } 0 \leq \tau \leq 1/2 \\
%     g(\tau)+2, & \text{for } 1/2 \leq \tau \leq 1.
%   \end{cases}
% \]
% By Theorem \ref{main}, this map is an injection from $\overline{B}^0[1/2,\,1]$ into the set of domains in $\RR^4$ with volume equal to $2$, modulo symplectomorphism. 
\end{remark}


Arguing similarly, Corollary \ref{newlag} yields the following. 

\begin{theorem}\label{main2} Let $\Omega \subset \RR^2_{\geq 0}$ be any set that contains the square 
$[0,1)\times[0,1)$ and is contained in the rectangle $[0,1) \times [0,2).$
Then 
\begin{equation*}
\label{}
   {\rm Sh}(X_{\Omega}) = \Bigl(\Omega  \cup \{0\leq r < 1/2\} \Bigr)\cap \{r\leq s\}. 
\end{equation*}
\end{theorem}


The sets whose Hamiltonian shape is given by  Theorem \ref{main2} are not assumed to have any additional structure beyond the two inclusions asserted, see Figure \ref{fig:shape2}.

\begin{figure}[!h]
    \centering
    \includegraphics[width=0.8\linewidth]{Shape2-1.pdf}
    \vspace{-7.5cm}
    \caption{Equality \eqref{eq:shape} holds when $[0,1)\times[0,1)\subset \Omega \subset [0,1) \times [0,2).$}
    \label{fig:shape2}
\end{figure}



\subsubsection{On the shape of ellipsoids} For symplectic ellipsoids $E(a,b)$ with $b>2a$ and $\frac{b}{a} \notin 2\NN$, it is still not known whether the inclusion
$${\rm Sh}(E(a,b)) \subset \Bigl(\mu(E(a,b))  \subset \{0\leq r < a/2\} \Bigr)\cap \{r\leq s\}$$
is an equality. Theorem \ref{main} can be used to refine this problem. 

\begin{corollary}
For $b\geq 2a >0$ let $\mathcal{T}(a,b)$ be the triangular domain defined by the inequalities $r\le s$, $s \geq b -br/a$ and $s <a$. Then 
$${\rm Sh}(E(a,b))=\Bigl(\mu(E(a,b))  \cup \{0\leq r < a/2\} \Bigr)\cap \{r\leq s\}$$
iff no torus $L(r,s)$ with $(r,s) \in \mathcal{T}(a,b)$ can be mapped into $E(a,b)$ by a Hamiltonian diffeomorphism of $\RR^4$ .
\end{corollary}

\begin{proof}
    By monotonoicity, the set ${\rm Sh}(E(a,b))$ is contained in ${\rm Sh}(E(a,b) \cup P(a,a))$ which, by Theorem \ref{main}, is equal to $$\left(\Bigl(\mu(E(a,b))  \cup \{0\leq r < a/2\} \Bigr)\cap \{r\leq s\}\right) \cup \mathcal{T}(a,b). $$
\end{proof}


% \begin{remark}If $\frac{b}{a}$ is in $\NN_{\geq 2}$, then,  by \cite{hizh}, $${\rm Sh}(E(a,b)) = \Bigl(\mu(E(a,b))  \cup \{0\leq r < a/2\} \Bigr)\cap \{r\leq s\}.$$ It is not known whether $\mathcal{T}(a,b)$, ever contributes to the shape of $E(a,b)$.
% \end{remark}


\begin{figure}[!h]
    \centering
    \includegraphics[width=1\linewidth]{ShapeE-1.pdf}
    \vspace{-10.5cm}
    \caption{It is not known when, if ever, a Lagrangian torus $L(r,s)$ with $(r,s)$ in the red region, $\mathcal{T}(a,b)$, can be symplectically embedded into $E(a,b)$.}
    \label{fig:ellipse}
\end{figure}

{\bf Acknowledgements.} The authors would like to thank Felix Schlenk for detailed comments on a draft, and also Jun Zhang for his interest and encouragement.

\section{The proof of Theorem \ref{lagint} modulo two propositions}

In this section we present the proof of Theorem \ref{lagint} assuming two propositions whose proofs are presented in Section \ref{foliation} and Section \ref{dbuildings}. More precisely, we first present a proof of a slightly weaker result, Theorem \ref{center}, below. In Section \ref{adjust}, we then show that the proof of Theorem \ref{center} can be extended to cover Theorem \ref{lagint}.
\begin{theorem}\label{center}
Suppose that  $1/2\leq r<1$ and $1 \leq s <b$. If $\phi$ is a Hamiltonian diffeomorphism of $\RR^4$ that maps $L(r,s)$ into $P(1,b)$, then $\phi(L(r,s))$ intersects the codimension one submanifold
\begin{equation*} 
\bigcup_{t \in [s,b)} L(r, t).
\end{equation*}
%Moreover, if $b<2s$, then $\phi(L(r,s))$ intersects $L(r,s)$.
\end{theorem}


\begin{remark}
    The choice to present Theorem \ref{lagint} as an extension of Theorem \ref{center} is made for the sake of clarity. The proof of Theorem \ref{center} contains all the key technical ideas and arguments, and the extension to Theorem \ref{lagint} is easily described in these terms. 
\end{remark}

\subsection{Proof of Theorem \ref{center}}

To prove Theorem \ref{center} we argue by contradiction. The assumption to be contradicted is the following.

\begin{assumption}\label{assump}
There are real numbers $r$, $s$ satisfying 
\[ \frac{1}{2} \leq r <1 \leq s,\]
a number $b>s$, and a Hamiltonian diffeomorphism $\phi$ of $\RR^4$ that maps $L(r,s)$ into $P(1,b)$, such that  $\phi(L(r,s))$ is disjoint from the family of Lagrangian tori $\bigcup_{t \in [s,b)} L(r, t)$.
\end{assumption}

\bigskip

To simplify our notation we set $L_1=L(r,s)$, $L_2= \phi(L(r,s))$ and $$L[t] = L(r, s-1+t)$$ for $t \in [1,b-s+1)$. With this, Assumption \ref{assump} asserts that $L_2$ is  disjoint from each of the $L[t]$. As in \cite{hiker}, the general philosophy behind the proof is that the assertion: {\em $L_1$ and $L_2$ are disjoint Lagrangian submanifolds in $(M,\omega)$} implies an identical assertion for any symplectic manifold that can be obtained from $(M,\omega)$ by applying standard symplectic procedures in the complement of $L_1 \cup L_2$; like compactification, symplectic blow-up, symplectic blow-down and inflation. Such changes of the ambient symplectic manifold will be referred to as {\em scene changes}. To contradict Assumption \ref{assump}, we will show that if $L_2$ is disjoint from the $L[t]$, then we can construct a scene in which a classical Lagrangian intersection result of Gromov implies that $L_1$ and $L_2$ must intersect.  
 % In general, the set of possible scenes is richer whenever $(M,\omega)$, or any symplectic manifolds obtained from it, has many rigid symplectic curves in the compliment of $L_1 \cup L_2$. 


% To reach the desired contradiction, it will be necessary to view the assumed pair of disjoint Lagrangian tori $L_1$ and $L$ in different settings. The following result facilitates this.


% \begin{lemma}\label{axisd} There exists a Hamiltonian isotopy of $P(1,b)$ which fixes the family $F(r,s,b)$ and  displaces $L$ from $\{ z_1 z_2 =0\}$.
% \end{lemma}

% \begin{proof}
% A PROOF OF THIS IS NEEDED. SHALL WE COMPACTIFY FIRST AND USE THE SMALL FOLIATION?

%  Equivalently we find a Hamiltonian isotopy displacing $\{z_1 z_2=0\}$ from $L$. For this, we note that $\{z_1=0\}$ and $\{z_2=0\}$ are represented by finite energy planes asymptotic to the long and short Reeb orbits on $\partial E(1,x)$ respectively. These planes persist under variations of almost complex structure, modulo bubbling in codimension $2$. In particular, generically they persist under the $1$ parameter family of almost complex structures given by stretching along $L$, and, generically again, will converge in the limit to finite energy planes, that is, there will be no degeneration to buildings with negative ends on $L$.

% The deformations of finite energy planes give Hamiltonian isotopies separately displacing $\{z_1=0\}$ and $\{z_2=0\}$ from $L$. However positivity of intersection implies that  throughout the isotopy the planes intersect transversally in a single point. It follows that the isotopies can be realized by a single Hamiltonian flow.
%  \end{proof}

% \bigskip

% \noindent{\bf Setting 1:  $T^*L_1$.}
% The set $P(1,b) \smallsetminus \{ z_1 z_2 =0\}$ can be identified with the  neighborhood 
% $$U(1,b) =\{-1< p_1< 1-r,\, -s<p_2<b-s\}$$ of $T^* L_1$ equipped with the canonical symplectic form $p_1 dq_1 + p_2dq_2$. Viewing $L_1$ here as the zero section and $L$ as a closed Lagrangian torus in $T^*L_1$, the following is a direct consequence of Assumption \ref{assump} and Gromov's intersection theorem for exact Lagrangians in cotangent bundles from Section  $2.3.B_4''$ of \cite{gr}.

% \begin{proposition}\label{nonexact}
% When considered as a Lagrangian torus in $U(1,b) \subset T^* L_1$, $L$ is not exact.    
% \end{proposition}



\bigskip
\noindent{\bf Scene 1.} Let $\omega$ be the standard symplectic form on the sphere $S^2$ with total area equal to $1$. For $\rho>0$, denote the symplectic manifold $(S^2, \rho \omega)$ by $S^2(\rho)$. For our first scene change we compactify the polydisk $P(1,b)$ to the symplectic manifold $$(X_1,\Omega_1) := (S^2 \times S^2, \omega \oplus b \, \omega) =S^2(1) \times S^2(b).$$ 

In this scene, $L_2$ is still disjoint from each of the $L[t]$ and all these tori lie in the complement of the union of symplectic spheres $S_{\infty} \cup T_{\infty}$, where $S_{\infty} = \{\infty\} \times S^2(b)$ and $T_{\infty} = S^2(1) \times \{\infty\}$.
 
 % Moreover, $L_2$ is disjoint from the (closure of) the family of Lagrangian tori  $$L_{1,\tau} = \partial D^1_0(r) \times \partial D^b_\infty(b-s) \subset S^2(1) \times S^2(b),$$ where $D^t_*(\delta)$ is the standard disk in $S^2(t)$ with area $\delta$ and center $*$.
% We use this to prove the following.
% \begin{proposition}\label{homological}
% Considered $L$ as a Lagrangian torus in $T^* L_1$. If $L$  is  homologically nontrivial, i.e.  $[L] \neq 0 \in H_2(U(1,b)) = \ZZ$,  then  $L$ is exact.    
% \end{proposition}

% Together with Proposition\ref{nonexact}, this implies that $L \subset T^*L$  is homologically trivial. This fact constrains the intersection patterns of relative foliations of $L$ and $L \cup L_1$ in $X_2$. In $\S$ \ref{moves}, we use these constraints to obtain from $L$ and $L_1$, two disjoint relatively exact Lagrangian tori in (a new copy of) the cotangent bundle of a torus,$T^*T^2$, in contradictions to Gromov's intersection theorem for exact Lagrangians in cotangent bundles.

\bigskip
\noindent{\bf Scene 2.}
Next, we inflate along the symplectic sphere $S_{\infty}$ by adding a tubular neighborhood of capacity $2r-1$. Denoting the resulting symplectic manifold by $(X_2, \Omega_2)$ we note that $(X_2, \Omega_2)$ is symplectomorphic to $S^2(2r) \times S^2(b)$. 
% and so we can identify $H_2(X_2;\ZZ)$ with $\ZZ \times \ZZ$ in the usual manner. 
Here again, $L_2$ is still disjoint from each of the $L[t]$ and all these tori lie in the complement of $S_{\infty} \cup T_{\infty}$ (this time viewed as symplectic spheres in $(X_2,\Omega_2)$ so that $S_{\infty} = \{\infty\} \times S^2(b)$ and $T_{\infty} = S^2(2r) \times \{\infty\}$).

In Section \ref{foliation} we extend the theory of relative finite energy foliations from \cite{rgi} and the techniques developed in \cite{hiker,boust} to prove the following result.

\begin{proposition}\label{prop:proj} There is an $\Omega_2$-tame (or compatible) almost complex structure $J_2$ on $X_2$  and a continuous (projection) map $p \colon X_2 \to S_{\infty}$ with the following properties.
\begin{enumerate}
    \item The image $p(T_{\infty})$ is a point.\\
    \item For every $z$ in the complement of $p(L_1) \cup p(L_2)$ the set $p^{-1}(z)$ is a $J_2$-holomorphic sphere in the class $(1,0) \in H_2(X_2;\ZZ) = \ZZ \times \ZZ$.\\
    \item $p(L_1)$ and $p(L_2)$ are disjoint embedded closed curves in $S_{\infty}$.\\
    \item The component of $S_\infty \smallsetminus p(L_1)$ that contains $p(T_\infty)$ is contained in the component of $S_\infty \smallsetminus p(L_2)$ that contains $p(T_\infty)$.
\end{enumerate}
\end{proposition}

\begin{remark} The proof of the final assertion is precisely where we use the assumption that $L_2$ is disjoint from $L[t]$ for all $t \in [1,b-s+1)$.
\end{remark}

Let $A \subset S_\infty$ be the component of $S_\infty \smallsetminus p(L_2)$ that does not contain $p(T_\infty)$, and let $C$ be the component of $S_\infty \smallsetminus p(L_1)$ that does. Then $A$ and $C$ are disjoint, and their complement is a cylinder, $B$, as pictured in Figure \ref{fig:abC}.  


\begin{figure}[!h]
    \centering
    \includegraphics[width=0.6\linewidth]{Abc-4.pdf}
    \vspace{-2.5cm}
    \caption{Projections of $L_1$ and $L_2$ onto $S_{\infty}$ under Assumption \ref{assump}.}
    \label{fig:abC}
\end{figure}

\medskip

Choose a point $z_0$ in $A$ and set $T_0 = p^{-1}(z_0).$ By Proposition \ref{prop:proj}, $T_0$ is a regular $J_2$-holomorphic sphere that is disjoint from $L_1 \cup L_2$. In Section \ref{dbuildings}, we study buildings obtained by taking the limit of high degree pseudo-holomorphic curves in $(X_2, \Omega_2)$ with point constraints along $L_1 \cup L_2$ as we stretch-the-neck along these Lagrangians. These buildings are used, together with Proposition \ref{prop:proj}, to prove the following.


\begin{proposition}\label{prop:submanifolds}
For every sufficiently large $d \in \NN$, there exist symplectic embeddings
\begin{align*}
    &F,G \colon S^2 \to X_2 \smallsetminus (L_1 \cup L_2) \\
    &D_A, D^0_2, D^\infty_2 \colon  (D^2,S^1) \to (X_2,L_2) \\
    &D_C, D^0_1, D^\infty_1 \colon  (D^2,S^1) \to (X_2,L_1) \\
    &O_B \colon (S^1 \times [0,1],S^1 \times \{0,1\}) \to (X_2,L_1 \cup L_2)
\end{align*}
such that:\\

\begin{enumerate}

\item  The maps $F$ and $G$ both represent the class $(d,1) \in H_2(X_2; \ZZ)$\\

\item The six disks all have Maslov index $2$.\\



\item The set $\{[\partial D_A], [\partial D^\infty_2]\}$ is an integral basis of $H_1(L_2;\ZZ)$, and $[\partial  D^0_2]=-[\partial  D^\infty_2]$.\\

\item The set $\{[\partial D_C], [ \partial D^\infty_1]\}$ is an integral basis of $H_1(L_1;\ZZ)$, and $[\partial  D^0_1]=-[\partial D^\infty_1]$.\\

\item Denoting the component of the boundary of $O_B$ on $L_i$ by $\partial^i O_B$, the set $\{[\partial^i O_B], [\partial D^\infty_i]\}$ is an integral basis of $H_1(L_i;\ZZ)$\\



\item The images of  $p \circ D_A$, $p \circ O_B$ and  $p \circ D_C$ are the closures of $A$, $B$ and $C$, respectively.\\

\item The symplectic areas of  $D_A$, $O_B$ and  $D_C$ are $s$, $r$ and $b-s$, respectively. The disks $D^\infty_1$ and $D^\infty_2$ both have symplectic area equal to $r$.\\

\item  All nine maps are $J_2$-holomorphic away from arbitrarily small neighborhoods of a collection of Lagrangian tori whose elements are close to, and Lagrangian isotopic to, either $L_1$ or $L_2$.\\

\item We have the following intersection pattern:\\
\begin{enumerate}
    \item $F \bullet G=2d$
    \item $F \bullet D^0_i + F \bullet D^\infty_i = G \bullet D^0_i + G \bullet D^\infty_i =1$ for $i=1,2$
    \item $\{F \bullet D^\infty_1, G \bullet D^\infty_1\} = \{F \bullet D^\infty_2, G \bullet D^\infty_2\} = \{0,1\}$    %$F \bullet D^\infty_1=1$ and $F \bullet D^\infty_2=0$
    \item $F \bullet D_A + G \bullet D_A = d$
    \item $F \bullet D_C + G \bullet D_C = d$
    \item $F \bullet O_B =1$ and $G \bullet O_B=0$
    \item $D^\infty_1 \bullet S_\infty = D^\infty_1 \bullet S_\infty =1$
    \item $F \bullet T_0 = F \bullet T_\infty=G \bullet T_0 = G \bullet T_\infty=1$.\\
\end{enumerate} 



% \item Exactly one of $F$ and $G$ intersect the planes of $\mathfrak{r}_0$ and the other intersects the planes of $\mathfrak{r}_\infty$.\\

% \item Exactly one of $F$ and $G$ intersect the planes of $\mathfrak{s}_0$ and the other intersects the planes of $\mathfrak{s}_\infty$.\\


% \item  $$F \bullet D_A + G \bullet D_A = d$$ $$F \bullet D_C + G \bullet D_C = d$$ $$F \bullet O_B =1$$ $$G \bullet O_B=0$$. 

\item The spheres $F$ and $G$ have exactly $2d$ intersection points. Of these, $d$ are located in $p^{-1}(A)$ and $d$ are located in $p^{-1}(C)$.\\

\end{enumerate} 



\end{proposition}


\bigskip

\noindent{\bf Scene 3.} Let $x_1, \dots, x_{2d}$ be the intersection points of $F$ and $G$. We may assume that $x_1, \dots, x_{d}$ lie in $p^{-1}(A)$ and that $x_{d+1}, \dots, x_{2d}$ lie in $p^{-1}(C)$. Let $H_i$ be the fiber $p^{-1}(p(x_i))$. By the second assertion of Proposition \ref{prop:proj}, the $H_i$ are disjoint symplectic spheres in the class $(1,0)$.



The next scene is obtained by performing a sequence of small symplectic blow-ups followed by a sequence of symplectic blow-downs. First we choose $J_2$ to be $\Omega_2$-compatible, from the start, and perturb it, as in \cite{boust}, to an almost complex structure $J_3$ for which $F$ and $G$ are complex, and which is the standard integrable structure in Darboux balls of capacity $\varepsilon>0$ centered at each of the $x_i \in F \cap G$. 
We may also assume that $F$, $G$ and the $H_i$ coincide with complex planes in the balls.
We then blow up these Darboux balls. The proper transforms of $F$, $G$ and the $H_i$ are well defined smooth symplectic surfaces because of our standard coordinates in the balls. Hence we can blow down the corresponding proper transforms of the $H_i$, symplectically again. The resulting manifold, $(X_3, \Omega_3)$, is diffeomorphic to $S^2 \times S^2$. Given a symplectic submanifold $S$ of $(X_2, \Omega_2)$, we will denote its proper transform in $(X_3, \Omega_3)$ by $\hat{S}$. The symplectic submanifolds $\hat{F}$ and $\hat{G}$ are disjoint spheres in $X_3$ and their self-intersection numbers are both $0$. The same is true of $\hat{T}_0$ and $\hat{T}_{\infty}$.  Moreover, both $\hat{T}_0$ and $\hat{T}_{\infty}$ intersect each of $\hat{F}$ and $\hat{G}$ positively in a single point. In other words, the collection $\hat{F} \cup  \hat{G} \cup \hat{T}_0 \cup \hat{T}_{\infty}$ serves as a set of (topological) axes in the new copy of $S^2\times S^2$ and these axes are disjoint from $L_1 \cup L_2$.
Note, however, that $(X_3, \Omega_3)$ is not monotone. The spheres $\hat{T}_0$ and $\hat{T}_{\infty}$ have area $2r$ while the spheres $\hat{F}$ and $\hat{G}$ have area $2dr + b - 2d \varepsilon$.


\bigskip

\noindent{\bf Scene 4.} Finally, we work away from $\hat{F} \cup  \hat{G} \cup \hat{T}_0 \cup \hat{T}_{\infty}$ and view $L_1$ and $L_2$ as disjoint Lagrangian tori in the open symplectic manifold 
$$X_4 := X_3 \smallsetminus ( \hat{F} \cup  \hat{G} \cup \hat{T}_0 \cup \hat{T}_{\infty} )$$
equipped the restriction of $\Omega_3$ which we denote by $\Omega_4$. By classification results for symplectic spheres, see \cite{mcdrat}, this final scene can be identified with the subset of $T^* T^2 = T^2 \times\RR^2$ defined by restricting the momentum coordinates to the rectangle with side lengths $2r$ and $2dr + b - 2d \varepsilon$.

\begin{lemma}\label{lem:pre}
    The Lagrangian tori $L_1$ and $L_2$ are relatively exact in $(X_4, \Omega_4)$ and both are homologous to the zero section.
\end{lemma}

\begin{proof}

To prove that $L_1$ and $L_2$ are relatively exact, it suffices to find two relative cycles,  $C_1$ and $C_2$, representing classes in $H_2(X_4, L_1 \cup L_2;\ZZ)$, such that 
\begin{equation}
\Omega_4(C_1) = \Omega_4(C_2) =0
\end{equation}
\begin{equation}
\{[\partial C_1 \cap L_1], [\partial C_2 \cap L_1]\}\quad \text{is a basis of} \quad   H_1(L_1;\ZZ),
\end{equation}
and 
\begin{equation}
\{[\partial C_1 \cap L_2], [\partial C_2 \cap L_2]\} \quad \text{is a basis of} \quad H_1(L_2;\ZZ).
\end{equation}

To produce these cycles we will use the transforms of the submanifolds of $X_2$ described in Proposition \ref{prop:submanifolds}.

By Proposition \ref{prop:submanifolds} we can choose disks $D_1$ equal to one of the transforms $\hat{D}^0_1$ or $\hat{D}^\infty_1$, and $D_2$ equal to one of the transforms $\hat{D}^0_2$ or $\hat{D}^\infty_2$, such that both disks intersect $\hat{F}$ exactly once but are disjoint from $\hat{G}$. We denote the intersection points with $\hat{F}$ by $q_1$ and $q_2$ respectively.
%Consider the transforms $\hat{D}^\infty_1$ and $\hat{D}^\infty_2$ and let $q_1$ and $q_2$ be their respective intersection points 
%with $S_{\infty}$. 
Let $\gamma_1$ be a smooth curve in $\hat{F}$ connecting $q_1$ to $q_2$, and $\Gamma_1$ a smooth cylinder in the normal bundle to $\hat{F}$ such that $\Gamma_1$ is disjoint from $\hat{F}$ and projects to $\gamma_1$. We may assume $\Gamma_1$ intersects $D_1$ and $D_2$ in circles around their intersections with $\hat{F}$. The circles bound disks $E_1$ and $E_2$ such that $D_1 \smallsetminus E_1$ and $D_2 \smallsetminus E_2$ are cylinders disjoint from $\hat{F}$, and we can define $C_1$ to be the glued cycle $C_1 = (D_1 \smallsetminus E_1) \cup \Gamma_1 \cup (D_2 \smallsetminus E_2)$. %can then be perturbed to be disjoint from $\hat{F} \cup  \hat{G} \cup \hat{T}_0 \cup \hat{T}_{\infty}$ to produce our first cycle $C_1$. 
It clearly has a zero $\Omega_4$-area. Moreover,  $[\partial C_1 \cap L_1]=-[\partial D_1]$ and $[\partial C_1 \cap L_2]=[\partial D_2]$.

To construct the cycle $C_2$, we first consider the transform $\hat{O}_B$ of $O_B$. It follows from assertion (5) of Proposition \ref{prop:submanifolds} that 
$\{[\partial C_1 \cap L_i], [\partial \hat{O}_B \cap L_i]\}$ is a basis of  $H_1(L_i;\ZZ)$ for $i=1,2$.
However the $\Omega_4$-area of $\hat{O}_B$ is $r$. Moreover, $\hat{O}_B$ intersects the axis $\hat{F}$ positively in precisely one point and so does not represent a class in $H_2(X_4, L_1 \cup L_2;\ZZ)$. Both problems can be fixed simultaneously by choosing a smooth curve $\gamma_2$ connecting $q_1$ to $\hat{O}_B \cap \hat{F}$. We can then form a cylinder $\Gamma_2$ over $\gamma_2$ and use it to glue $D_1$ and $\hat{O}_B$ to produce the required cycle $C_2$ in the complement of $\hat{F}$ as above.
%replacing $\hat{O}_B$ with $(\hat{D}^\infty_1)^- \# \gamma_2 \# \hat{O}_B$ where $\gamma_2$ is a smooth curve in $\hat{F}$ from the point  $\hat{D}^\infty_1 \cap \hat{F}$ to the point $\hat{O}_B \cap \hat{F}$. 
%The resulting pair-of-pants map can then be perturbed off of $\hat{F}$ so that its image lies in $X_4$. This is the required cycle, $C_2$.


It remains to show that $L_1$ (or $L_2$) is homologous to the zero section. It is enough to show that the inclusion $H_1(L_1;\ZZ) \to H_1(X_4;\ZZ)$ is injective or, equivalently, that the boundary map $$\partial \colon H_1(X_4,L_1;\ZZ) \to H_1(L_1;\ZZ)$$ is trivial. If the boundary map is not trivial then there is a relative 2-cycle $C$ such that $[\partial C] = (m,n) \neq (0,0) \in H_1(L_1, \ZZ).$ Note that $C$ also represents a class in $H_2(X_3, L_1;\ZZ)$. Hence, adding suitable multiples of $\hat{D}_C$ and $\hat{D}_1^{\infty}$ to $C$ we get a closed cycle $C_{m,n}$ in $X_3$ whose intersections with the topological axes are determined by the disks added. If $m\neq 0$, then $C_{m,n}$, like $\hat{D}_C$,  will intersect $\hat{T}_0$ but not $\hat{T}_{\infty}$. Otherwise, if $m=0$ and $n\neq0$, then $C_{m,n}$, like $\hat{D}_1^{\infty}$,  will intersect $\hat{F}$ but not $\hat{G}$. In both cases, the resulting intersection pattern is impossible since $[\hat{T}_0]=[\hat{T}_{\infty}] \in H_2(X_3;\ZZ)$ and $[\hat{F}] =[\hat{G}]\in H_2(X_3;\ZZ)$.
\end{proof}

Lemma \ref{lem:pre} immediately yields the desired contradiction. Within $(X_4, \Omega_4) \subset (T^*T^2, d \lambda)$, it follows from Lemma \ref{lem:pre} and Theorem B of \cite{rgi}, that $L_1$ and $L_2$ are both Hamiltonian isotopic to the same constant section of $T^*T^2$. With this, Gromov's Theorem 2.3.$B_4''$ from \cite{gr} implies that $L_1$ and $L_2$ must intersect.  This contradicts Assumption \ref{assump}. To prove Theorem \ref{center} then, it remains for us to prove Proposition \ref{prop:proj} and Proposition \ref{prop:submanifolds}.

\subsection{Proof of Theorem \ref{lagint}}\label{adjust}
 
To extend the proof of Theorem \ref{center} to a proof of Theorem \ref{lagint}, we need to consider a third Lagrangian torus, $L_3=L(r, b-s)$. The (weaker) standing assumption is now as follows.

\begin{assumption}\label{assump2}
There are real numbers $r$, $s$ satisfying 
\[ \frac{1}{2} \leq r <1 \leq s,\]
a number $b>2s$, and a Hamiltonian diffeomorphism $\phi$ of $\RR^4$ that maps $L(r,s)$ into $P(1,b)$, such that  $\phi(L(r,s))$ is disjoint from the family of Lagrangian tori $\bigcup_{t \in [s,b-s]} L(r, t)$ from $L_1$ to $L_3$.
\end{assumption}



Under this assumption, relative foliation theory implies the following analogue of Proposition \ref{prop:proj}.


\begin{proposition}\label{adjustprop} There is an $\Omega_2$-tame (or compatible) almost complex structure $J_2$ on $X_2$  and a continuous (projection) map $p \colon X_2 \to S_{\infty}$ with the following properties.
\begin{enumerate}
    \item The image $p(T_{\infty})$ is a point.\\
    \item For every $z$ in the complement of $p(L_1) \cup p(L_2)  \cup p(L_3)$ the set $p^{-1}(z)$ is a $J_2$-holomorphic sphere in the class $(1,0) \in H_2(X_2;\ZZ) = \ZZ \times \ZZ$.\\
    \item The images $p(L_i)$, for $i= 1,2,3$ are disjoint embedded closed curves in $S_{\infty}$.\\
    \item The component of $S_\infty \smallsetminus p(L_3)$ that contains $p(T_\infty)$ is contained in the component of $S_\infty \smallsetminus p(L_1)$ that contains $p(T_\infty)$. As well, the projection $p(L_2)$ is contained in the component of $S_\infty \smallsetminus p(L_3)$ that contains $p(T_\infty)$ or it is contained in the component of $S_\infty \smallsetminus p(L_1)$ that does not contain $p(T_\infty)$.
\end{enumerate}
\end{proposition}

Let $\mathcal{A}$ to be the component of $S_\infty \smallsetminus p(L_1)$ that does not contain $p(T_\infty)$, and $\mathcal{C}$ be the component of $S_\infty \smallsetminus p(L_3)$ that contains $p(T_\infty)$. Then $\mathcal{A}$ and $\mathcal{C}$ are disjoint disks and their complement is a cylinder, 
$\mathcal{B}$. This is pictured in Figure \ref{fig:abc2}. The three possible positions of $p(L_2)$ are pictured in Figure \ref{fig:abc-l2}.

\begin{figure}[!h]
    \centering
    \includegraphics[width=0.6\linewidth]{Abc-cal-1.pdf}
    \vspace{-2.5cm}
    \caption{Projections of $L_1$ and $L_3$ onto $S_{\infty}$ under Assumption \ref{assump2}.}
    \label{fig:abc2}
\end{figure}

\medskip

\begin{figure}[!h]
    \centering
    \includegraphics[width=0.9\linewidth]{Abc-cal-l2.pdf}
    \vspace{-10.5cm}
    \caption{The three possibilities for the location of $p(L_2)$.} 
    \label{fig:abc-l2}
\end{figure}




% An analogue of Proposition \ref{prop:proj} in this situation implies that
% %By Assumption \ref{assump2},
% $p(L_2)$ is contained in either $\mathcal{A}$ or $\mathcal{C}$. The proof of this follows the same lines as that of Proposition \ref{prop:proj} in section \ref{foliation}, except that instead of considering the family of Lagrangian tori $L[t]$ for $t \in [1, b-s+1-\epsilon]$ we restrict to the family $L[t]$ for $t \in [1, b-2s+1]$. The result is that $p(L_2)$ is disjoint from $\mathcal{B}$, rather than what would be the complement of $\mathcal{A}$ as in  Proposition \ref{prop:proj}.


If we denote the component of $S_\infty \smallsetminus p(L_2)$ that does not contain $p(T_\infty)$ by $A$, then we have the following three cases to consider: 

\begin{enumerate}
    \item $A$ is contained in $\mathcal{A}$,
    \item $A$ is contained in $\mathcal{C}$,
    \item $A$ contains both $\mathcal{A}$ and $\mathcal{B}$.
\end{enumerate}
As summarized in Figure  \ref{fig:abc3}, for each of these three cases the proof of Theorem \ref{lagint} can be reduced to that of Theorem \ref{center}.
% The three possible positions of $p(L_2)$, and hence $A$, are pictured in Figure \ref{fig:abc3}. Also pictured is the way in which the proof of Theorem \ref{lagint} can be reduced to that of Theorem \ref{center} in each case. As described below, the only technical issue involved in this extension is the need to ensure the existence of the analogue of the symplectic embedding $D_A$ in Proposition \ref{prop:submanifolds} in Case 3, where $L_2$ plays a different role in the proof of Theorem \ref{center} is modified, as it is in Case 3. 
% We can find an almost complex structure compatible with all three tori, and the corresponding projection map gives three disjoint embedded circles in $S_{\infty}$, namely the $p(L_i)$ for $i=1,2,3$. We define subsets $E$, $F$, $G$ by setting $E$ to be the component of $S_{\infty} \smallsetminus p(L_1)$ which does not contain $p(T_{\infty})$, and setting $G$ to be the component of $S_{\infty} \smallsetminus p(L_3)$ which does contain $p(T_{\infty})$. As in Proposition \ref{prop:proj} we can see that these sets are disjoint, and let $F$ be their complement, a cylinder.
% Next, as above, we set $A$ to be the component of $S_{\infty} \smallsetminus p(L_2)$ which does not contain $p(T_{\infty})$. Proposition \ref{prop:proj} applies similarly in this situation to show that $p(L_2) \subset E \cup G$, crucially using our hypothesis that $L_2$ is disjoint from $\cup_{t \in [s, b-s]} L(r,t)$. The proof then separates into three cases depending upon the region $A$.
In Case 1 we are in the exact scenario described in statement (4) of Proposition \ref{prop:proj} and the proof of Theorem \ref{center} applies directly. For Case 2, we fix a point $z_0 \in \mathcal{A}$ and set $T_0 = p^{-1}(z_0)$. Then, substituting $T_0$ for $T_{\infty}$ and $L_3$ for $L_1$, we can again follow the proof of Theorem \ref{center}. 
% In particular, the assumption that $L_2$ is Hamiltonian isotopic to $L(r,s)$ inside $\RR^4$ is enough to find a symplectic embedding $D_A$ , projecting to $A$ and of area $s$, as in Proposition \ref{prop:submanifolds}. 
In the end, the argument implies that $L_2$ must intersect $L_3$, which contradicts  Assumption \ref{assump2}.
Finally, in Case 3, the proof of Theorem \ref{center} can be applied directly, but with the roles of $L_1$ and $L_2$ reversed. 

\begin{figure}[!h]
    \centering
    \includegraphics[width=0.9\linewidth]{Abc-reduction.pdf}
    \vspace{-10.5cm}
    \caption{The colorings, when compared to that Figure \ref{fig:abC}, encode the manner in which the proof of Theorem \ref{lagint} can be reduced to that of Theorem \ref{center} in each case.}
    \label{fig:abc3}
\end{figure}

\medskip

\subsection{More Lagrangian Intersections} \label{sec:2lag}

As an alternative to taking the limit $b\to 2s$, Corollary \ref{newlag} could also be obtained by following the argument in Section \ref{adjust} but setting $L_1 = L_3$. In this case, $\mathcal{B} = \emptyset$ and Proposition \ref{adjustprop} (4) is automatic, that is, we do not need to study the $1$-parameter family of tori. Hence, rather than assuming $L_1 = L(r,s)$, we may instead assume that $L_1 = \phi_1(L(r,s))$ for some Hamiltonian diffeomorphism $\phi_1$ of $\RR^4$ that maps $L(r,s)$ into $P(1,b)$. With this we get the following more general intersection result.

\begin{corollary}\label{newlagv2}
Suppose that $1/2\leq r<1$ and $1 \leq s <b \le 2s$. If $\phi_1$ and $\phi_2$ are Hamiltonian diffeomorphisms of $\RR^4$ which each map $L(r,s)$ into $P(1,b)$, then $$\phi_1(L(r,s)) \cap \phi_2(L(r,s)) \neq \emptyset.$$
\end{corollary}

\begin{remark} Corollary \ref{newlagv2} fails if we simply assume that $L_3$ has the Maslov and area class of $L(r,s)$. Indeed, there is an isomorphism from $H_1(L(3/5, 4/5))$ to $H_1(L(4/5,1))$ which preserves both Maslov and area classes (while the product tori are clearly disjoint in $P(1,2)$). The assumption on the Hamiltonian isotopy class is used in Corollary \ref{>r}.
\end{remark}

% As mentioned in section \ref{secgen}, the assumption that $s \ge 1$ is also necessary. Without this the projection map described in Proposition \ref{prop:proj} may not have the required properties, in particular we cannot guarantee the broken planes $\mathfrak{r}_0$, $\mathfrak{r}_{\infty}$, $\mathfrak{s}_0$, $\mathfrak{s}_{\infty}$ described in section \ref{features} have area $r$, nor that the planes $D_A$ and $D_C$ in Proposition \ref{prop:submanifolds} have areas $s$ and $b-s$. However, in the case of monotone Lagrangian tori these area conclusions are automatic, as all Maslov $2$ planes have the same area. Thus Corollary \ref{newlagv2} has a monotone version as follows.

% \begin{corollary}\label{newlagv3}
% Suppose that $1/2\leq r<1$ and $b \le 2r$. Suppose $L_1$ and $L_2$ are monotone Lagrangian tori in $P(1,b)$ with Maslov $2$ disks having area $r$. Then $L_1 \cap L_2 \neq \emptyset$.
% \end{corollary}

% By \cite[Propisition 3.6]{hiker} any two monotone Lagrangian tori in $S^2(1) \times S^2(1)$ can be displaced from the axes, giving monotone tori in $P(1,1)$ with Maslov $2$ disks of area $1/2$. Hence Corollary \ref{newlagv3} implies the following.

% \begin{corollary} Any two monotone Lagrangian tori in $S^2 \times S^2$ intersect.
% \end{corollary}







\section{The proof of Proposition \ref{prop:proj}}\label{foliation}

We first show that the relative foliation theory from \cite{ivriit} and \cite{rgi} can be applied to the pair $L_1 \cup L_2 \subset X_2$, despite the fact that neither Lagrangian torus is monotone. To simplify our notation, we set $(X_2,\Omega_2) = (X,\Omega)$.
%The same holds for the pairs  $L[t]$ and $L_2$ for $t \in (1, b-s+1).$

\subsection{Foliations of $(X, \Omega)$ relative to $L_1 \cup L_2$.}\label{subsec:morefol} 

Let $J$ be an $\Omega$-tame almost complex structure on $X$ such that $J$ is equal to the standard split complex structure near $S_{\infty} \cup T_{\infty}$. 
In fact, by the construction of $X=X_2$, in Scene $2$, we may assume that $J$ is standard in a tubular neighborhood of $S_{\infty}$ of capacity $2r-1$. This implies that $J$-holomorphic curves intersecting $S_{\infty}$ have area at least $2r-1$.

Since $J$ is standard near $T_{\infty}$, it follows from Theorem 0.2.A of \cite{gr} that the $J$-holomorphic sphere $T_{\infty}$ extends to a foliation $\mathcal{F}(J)$ of $X$ comprised of $J$-holomorphic spheres in the class $[S^2 \times \{pt\}]$. %By our choice of $J$,  $T_{\infty}$ is a leaf of the foliation $\mathcal{F}(J)$.
Moreover, there is a smooth projection map from $X$ to $S_{\infty}$ defined by taking a point $p$ in $X_2$ to the unique intersection point of $S_{\infty}$ and the leaf of $\mathcal{F}(J)$ that contains $p$.

To produce a relative version of this picture, we restrict $J$ further. For each Lagrangian $L_i$ we fix a parameterization  $\psi_i \colon \mathbb{T}^2 \to L_i$. This extends to a symplectic embedding  $\Psi_i \colon \mathcal{U}(\epsilon) \to X$ where $$\mathcal{U}(\epsilon) = \{|p_1| <\epsilon,\, |p_2|< \epsilon\} \subset T^* \mathbb{T}^2$$ for some sufficiently small $\epsilon>0.$  Set $$\mathcal{U}_i =\Psi_i(\mathcal{U}(\epsilon)).$$ We assume that $J$ is {\it compatible} with the parameterizations $\psi_i$ in the sense that in each neighborhood $\mathcal{U}_i$ we have 
$$J \frac{\partial}{\partial q_i} = - \frac{\partial}{\partial p_i}.$$
Given this, $J$ can be stretched along $L_1 \cup L_2$, as in \cite{behwz}, to yield  an almost complex structure $J_{\infty}$ on $X \smallsetminus (L_1 \cup L_2)$. 
This new almost complex structure is {\it adapted} to the parameterizations $\psi_i$, i.e. in $\mathcal{U}_i \smallsetminus L_i$ we have $$J_\infty \frac{\partial}{\partial q_i} = -  \sqrt{p^2_1 + p_2^2}\,\frac{\partial}{\partial p_i}.$$ The almost complex structure $J_\infty$ is also standard near $S_{\infty} \cup T_{\infty}$. 
% and, for a generic choice of the original $J$, we may assume that $J_{\infty}$ is regular for all $J_{\infty}$-holomorphic punctured spheres that are somewhere injective and have finite energy.
As described in \cite{behwz}, in stretching along $L_1 \cup L_2$, the curves of the Gromov foliation $\mathcal{F}(J)$ converge to a collection of holomorphic buildings, $\bar{\mathcal{F}}_{L_1,L_2}(J)$. Each building consists of genus zero pseudo-holomorphic curves of three possible level types. The {\em top level} curves  map to $X \smallsetminus (L_1 \cup L_2)$ and are $J_{\infty}$-holomorphic. The {\em middle level} curves map to one of two copies of the symplectization of the unit cotangent bundle of the flat torus that correspond to 
$L_1$ and $L_2$. These middle level curves are all $J_{\mathrm{cyl}}$-holomorphic  where $J_{\mathrm{cyl}}$ is the cylindrical almost-complex structure defined by $$J_{\mathrm{cyl}} \frac{\partial}{\partial q_i} = -  \sqrt{p^2_1 + p_2^2}\,\frac{\partial}{\partial p_i}.$$ The {\em bottom level} curves map to one of two copies of $T^*\mathbb{T}^2$  that again correspond to 
$L_1$ and $L_2$.\footnote{In what follows we will identify the copy of $T^*\mathbb{T}^2$ corresponding to $L_i$ with $T^*L_i$.} These bottom level curves are $J_{\mathrm{std}}$-holomorphic where  $J_{\mathrm{std}}$ is of the form 
$$J_{\mathrm{std}} \frac{\partial}{\partial q_i} = -  \rho \left(\sqrt{p^2_1 + p_2^2}\right)\,\frac{\partial}{\partial p_i}$$ for a fixed positive increasing function $\rho(t)$ such that $\rho(0)=1$ and $\rho(t) =t$ for $t \ge 2.$ 


A generic building of the collection $\bar{\mathcal{F}}_{L_1,L_2}(J)$ is (still) a simple $J$-holomorphic sphere in $X \smallsetminus (L_1 \cup L_2)$. For example, (the standard parameterization of) $T_{\infty}$ is such a building. The other buildings in $\bar{\mathcal{F}}_{L_1,L_2}(J)$ are split along either $L_1$ or $L_2$ and have multiple components. 

% These split buildings  have the following two possible types.

\begin{definition}\label{one}
    A split building in $\bar{\mathcal{F}}_{L_1,L_2}(J)$ is of {\it Type A} if it has three component curves; two $J_{\infty}$-holomorphic planes  in $X \smallsetminus (L_1 \cup L_2)$ of index $1$ and asymptotic to the same Lagrangian torus, say  $L_i$, and a $J_{\rm std}$-holomorphic cylinder of index $2$ with image in $T^*L_i$. Exactly one of the planes intersects $S_{\infty}$.  
\end{definition}


\begin{definition}\label{two}
    A split building in $\bar{\mathcal{F}}_{L_1,L_2}(J)$ is of {\it Type B} if it has five component curves; two  $J_{\infty}$-holomorphic planes of index $1$ in $X \smallsetminus (L_1 \cup L_2)$, %neither of which intersect $S_{\infty}$, 
    a nontrivial $J_{\infty}$-holomorphic cylinder of index 0 in  $X \smallsetminus (L_1 \cup L_2)$, and  
     %This cylinder connects Reeb orbits over disjoint embedded geodesics on $L_1 \cup L_2$. In particular, both ends of the cylinder can be on $L_1$, both can be on $L_2$, or the cylinder can have one end on each. 
pseudo-holomorphic cylinders in $T^*L_1 \cup T^*L_2$ of index $2$. 
\end{definition}

\begin{definition}
    An almost complex structure $J$ on $(X, \Omega)$ that is compatible with the parameterized Lagrangian tori $L_1$ and $L_2$ is of {\it foliation type} relative to the parameterized Lagrangian tori $L_1$ and $L_2$ if $\bar{\mathcal{F}}_{L_1,L_2}(J)$ consists only of $J_{\infty}$-holomorphic spheres and split buildings of Type A.
\end{definition}


\begin{lemma}\label{lem:noB}
Suppose  $J$ is standard near  $S_{\infty} \cup T_{\infty}$, compatible with $L_1$ and $L_2$,  and $J_{\infty}$ is regular for somewhere injective curves. Then $J$ is of foliation type relative to $L_1$ and $L_2$.  
\end{lemma}


\subsubsection{Proof of Lemma \ref{lem:noB}}
It follows as in Proposition 3.5 and Proposition 5.3 of \cite{rgi}, that for $J$ as in the statement, every split building in 
$\bar{\mathcal{F}}_{L_1,L_2}(J)$ is either of Type A or Type B. If $L_1$ and $L_2$ are monotone then, as described in Corollary 5.4 of \cite{rgi}, there are no split buildings of Type B. However, in the current setting, $L_1$ and $L_2$ are nonmonotone. 
To preclude buildings of Type B in our setting, we invoke the following result implied by Lemma 3.7 of \cite{ho}.

\begin{lemma}\label{>r2} Let $L$ be a parameterized Lagrangian torus in a symplectic manifold $(M,\omega)$ of dimension four. For an $\omega$-compatible almost complex structure $J$ on $M$ which is compatible with the parameterization of $L$ set $$a(J) = \min\{  \omega(u) \, | \, \mathrm{index}(u) \ge \mathrm{end}(u) \}$$
where the minimum is taken over $J_{\infty}$-holomorphic curves $u$ of genus zero with ends asymptotic to $L$, and $\mathrm{end}(u)$ denotes the number of ends of $u$.

\begin{enumerate}
\item If $M$ is closed, then $a(J)$ is independent of $J$.\\

\item Suppose that $(M,\omega)$ is an open and  geometrically bounded symplectic $4$-manifold. Let $J_{\mathrm{fix}}$ be an $\omega$-compatible almost complex structure such that the Riemannian metric $\omega(\cdot, J_{\mathrm{fix}} \cdot )$ is complete, has bounded sectional curvature and a positive lower bound for its injectivity radius. For a bounded open neighborhood $U$ of $L$, let $\mathcal{J}_{\mathrm{fix}} (L,U)$ be the set of $\omega$-compatible almost complex structures which are compatible with the parameterization of $L$ and equal to  $J_{\mathrm{fix}}$ in the complement of $U$. Then $a(J)$ is constant on $\mathcal{J}_{\mathrm{fix}} (L,U)$.
\end{enumerate}


% Suppose $(M,\omega)$ is an open,  geometrically bounded symplectic $4$-manifold and let $J_{\mathrm{fix}}$ be an $\omega$-compatible almost complex structure such that the Riemannian metric $\omega(\cdot, J_{\mathrm{fix}} \cdot )$ is complete, has bounded sectional curvature and a positive lower bound for its injectivity radius. Let $L$ be a parameterized Lagrangian torus in $(M,\omega)$ and fix a bounded open neighborhood $U$ of $L$. Let $\mathcal{J}_{\mathrm{fix}} (L,U)$ be the set of $\omega$-compatible almost complex structures which are compatible with the parameterization of $L$ and equal to  $J_{\mathrm{fix}}$ in the complement of $U$. For $J \in \mathcal{J}_{\mathrm{fix}} (L,U)$ set
% $$a(J) = \min\{  \omega(u) \, | \, \mathrm{index}(u) \ge \mathrm{end}(u) \}$$
% where the minimum is taken over $J_{\infty}$-holomorphic curves $u$ of genus zero with ends asymptotic to $L$, and $\mathrm{end}(u)$ denotes the number of ends of $u$. Then, the numbers $a(J)$ are independent of $J \in \mathcal{J}_{\mathrm{fix}} (L,U)$. In particular, given almost complex structures $J_0$ and $J_1$ in $\mathcal{J}_{\mathrm{fix}} (L,U)$, all $J_1$-holomorphic Maslov $2$ planes asymptotic to $L$ have area at least $a(J_0)$. 
% % Moreover, in a generic family $\{J_t\}$ in $\mathcal{J}_{\mathrm{fix}} (L,U)$, the universal moduli spaces of curves satisfying $\mathrm{index}(u) \ge \mathrm{end}(u)$ and $\omega(u) = a(J_0)$ are compact.
\end{lemma}

\begin{proof}

In both cases, the proof follows directly from that of  \cite[Lemma 3.7]{ho} which we briefly recall. Arguing by contradiction, we assume that there are $\omega$-compatible almost complex structures $J_0$ and $J_1$ that are compatible with the parameterization of $L$, and lie %in the second case are 
in the same set $\mathcal{J}_{\mathrm{fix}} (L,U)$, such that there is a $J_1$-holomorphic curve $u_1$ with $\mathrm{index}(u_1) \ge \mathrm{end}(u_1)$ and $\omega(u_1) < a(J_0).$  Let $\{J_t\}_{t \in[0,1]}$ be a smooth family of almost complex structures from $J_0$ to $J_1$ which all have the properties of $J_0$ and  $J_1$ described above. As noted in \cite[Lemma 3.7]{ho}, the results from \cite{we}, \cite{ze} and \cite{ohzh} imply that for a generic family $\{J_t\}$  the corresponding family moduli space, $\widetilde{\mathcal{M}}$, of curves in the class of $u_1$, is a smooth finite dimension manifold with boundary corresponding to $t=0$ and $t=1$. Since all curves in the moduli space have the same area, $\omega(u_1)$, and, by assumption, $\omega(u_1)<a(J_0)$, there must be no curves in $\widetilde{\mathcal{M}}$ corresponding to $t=0$. Fixing a constraint point on each of the asymptotic Reeb orbits of $u_1$ and a suitable number of constraint points in the image of $u_1$, we can cut $\widetilde{\mathcal{M}}$ down to a manifold $\mathcal{M}$ of dimension one. Since $\mathcal{M}$ also has no curves corresponding to $t=0$, it must be noncompact. However, as shown in \cite[Lemma 3.7]{ho}, any degeneration, at say $t_c \in (0,1)$, will result in a building which includes a $J_{t_c}$-holomorphic curve $u_{t_c}$ with $\mathrm{index}(u_{t_c}) \ge \mathrm{end}(u_{t_c})$ and $\omega(u_{t_c}) < a(J_0)$. Repeating the argument on the interval $[0,t_c]$, one eventually arrives at such a curve at $t=0$, and hence a contradiction.

% which implies the same result for $Y^4 = \RR^4$ and a product torus $L(1,x)$. The idea is argue by contradiction and to consider a universal moduli space of curves corresponding to a generic family $\{J_t\}$, and representing the class of a $J_1$-holomorphic curve $C_1$ with $\sigma(C_1) < a(J_0)$. We note that automatic regularity holds for curves satisfying $\mathrm{index}(C) \ge \mathrm{end}(C)$, so these universal moduli spaces are manifolds. The moduli space may be noncompact, but \cite[Lemma 3.7]{ho} shows that any degenerations will result in a $J_t$-holomorphic curve $C_t$, again with $\mathrm{index}(C_t) \ge \mathrm{end}(C_t)$ and $\sigma(C_t) < a(J_0)$. Repeating the argument on the interval $[0,t]$, eventually we arrive at a contradiction to the definition of $a(J_0)$.

% This works generally, the assumption on $Y$ and $L$ in \cite[Lemma 3.7]{ho} was used only to calculate $a(J_0) =1$ when $J_0$ is the standard product structure.
\end{proof}



Applying Lemma \ref{>r2} to product Lagrangian tori in $\RR^4$ we get the following useful Corollary.

\begin{corollary}\label{>r}
Suppose that $0< r\leq s$. Let $J$ be any tame almost complex structure on $\RR^4$ that is compatible with a parametrization of $L(r,s)$, and agrees with the standard complex structure, $J_0$, outside of a bounded domain.
If $J_{\infty}$ is the almost complex structure on $\RR^4 \smallsetminus L(r,s)$ resulting from stretching-the-neck along $L(r,s)$, then every $J_{\infty}$-holomorphic plane in $\RR^4 \smallsetminus L(r,s)$ with  deformation index 1 has symplectic area at least $r$. 
\end{corollary}

\begin{proof} For a product Lagrangian $L(r,s)$ with $r\le s$ it is straightforward to show that $a(J_0) =r$.
    
\end{proof}

%For the sake of completeness, a proof of Corollary \ref{>r} is contained in Appendix \ref{appen}.
With Corollary \ref{>r} in hand, we can now complete the proof of Lemma \ref{lem:noB}.
Arguing by contradiction, we assume that $\bar{\mathcal{F}}_{L_1,L_2}(J)$ includes a building $\mathbf{B}$ of Type B. As $T_{\infty}$ is a simple (unbroken) holomorphic sphere, all the curves of $\mathbf{B}$ must be disjoint from $T_{\infty}$. A single curve of $\mathbf{B}$ intersects $S_{\infty}$, and so by monotonicity has area greater than $2r-1$.

If the two top level planes of $\mathbf{B}$ both avoid $S_{\infty}$, then we can apply Corollary \ref{>r} to conclude they both have area at least $r$. As limiting buildings have the same total area as curves in the Gromov foliation, namely $2r$, this gives a contradiction since the top level cylinder of $\mathbf{B}$ must also have positive area. Hence, we conclude that exactly one of the planes of $\mathbf{B}$ intersects $S_{\infty}$, and so by monotonicity has area greater than $2r-1$. By Corollary \ref{>r}, the other plane has area at least $r$, and fits together with the cylinder to give a Maslov $2$ disk in $\RR^4$ with boundary on either $L_1$ or $L_2$. This class has area strictly greater than $r$, which implies that it has area at least $s$. But then the area of the building $\mathbf{B}$ exceeds $s + (2r-1) \ge 2r$, which is again a contradiction. This completes the proof of Lemma \ref{lem:noB}.


% Now arguing by contradiction let us assume that $\bar{\mathcal{F}}_{L_1,L_2}(J)$ has a building not of Type A. As $T_{\infty}$ is a simple (unbroken) holomorphic sphere, all curves in our building are disjoint from $T_{\infty}$. A single curve intersects $S_{\infty}$, and so by monotonicity has area greater than $2r-1$.

% If the two top level planes both avoid $S_{\infty}$, then we can apply Corollary \ref{>r} to conclude they both have area at least $r$. As limiting buildings have the same total area as curves in the Gromov foliation, namely $2r$, this gives a contradiction as the top level cylinder must also have positive area. Hence we conclude that one of the planes intersects $S_{\infty}$, and so by monotonicity has area greater than $2r-1$. The other plane has area at least $r$, and fits together with the cylinder to give a Maslov $2$ class in $\RR^4$ with boundary on either $L_1$ or $L_2$. This class has area strictly greater than $r$, which implies it has area at least $s$. But then the area of the building exceeds $s + (2r-1) \ge 2r$, which is again a contradiction. This completes the proof of Lemma \ref{lem:noB}.


 
\subsubsection{The features of a foliation of $X$ relative to $L_1 \cup L_2$}\label{features}
Assume that $J$  is of foliation type relative to $L_1$ and $L_2$. Discarding the middle and bottom  level curves of the Type A buildings in $\bar{\mathcal{F}}_{L_1,L_2}(J)$, one obtains a foliation 
$\mathcal{F}_{L_1,L_2}(J)$ of $X\smallsetminus (L_1 \cup L_2)$ with the following properties.

\bigskip

\begin{itemize}
\item[(F1)] The foliation $\mathcal{F}_{L_1,L_2}(J)$ has three kinds of leaves: unbroken ones consisting of a single closed $J_{\infty}$-holomorphic sphere in $X\smallsetminus (L_1 \cup L_2)$ of class $[S^2 \times \{pt\}]$, broken leaves along $L_1$ that consist of a pair of finite energy $J_{\infty}$-holomorphic planes in $X\smallsetminus (L_1 \cup L_2)$ that are asymptotic to $L_1$, and broken leaves along $L_2$ that consist of a pair of finite energy $J_{\infty}$-holomorphic planes in $X\smallsetminus (L_1 \cup L_2)$ that are asymptotic to $L_2$. 
  
\item[(F2)] By positivity of intersection, each leaf of $\mathcal{F}_{L_1,L_2}(J)$ intersects $S_{\infty}$ in exactly one point. For a broken leaf this means that exactly one of its planes intersects $S_{\infty}$. We denote the two families of broken planes asymptotic to $L_1$ by $\mathfrak{r}_0$ and $\mathfrak{r}_{\infty}$, and the two families of planes asymptotic to $L_2$ by  $\mathfrak{s}_0$ and $\mathfrak{s}_{\infty}$. We label these so that the planes in $\mathfrak{r}_{\infty}$ and $\mathfrak{s}_{\infty}$ intersect $S_{\infty}$. The planes in $\mathfrak{r}_0$ and $\mathfrak{s}_0$ are disjoint from $S_{\infty}$.
  
\item[(F3)] The ends of the two planes of a broken leaf are asymptotic to the same embedded geodesic in $L_i$, but with opposite orientations. We denote the homology class of this geodesic, equipped with the orientation determined by the plane that intersects $S_{\infty}$, by $\beta_i\in H_1(L_i; \ZZ)$. The class, $\beta_i$, is the same for all broken leaves of $\mathcal{F}_{L_1,L_2}(J)$ along $L_i$. It will be  referred to as the foliation class of $\mathcal{F}_{L_1,L_2}(J)$ for $L_i$.  
  

\item[(F4)] Each point in $X \smallsetminus (L_1 \cup L_2)$ lies in a unique leaf of $\mathcal{F}_{L_1,L_2}(J)$, and each point of $L_i$ lies on a unique geodesic in the foliation class $\beta_i$. Hence, there is a well-defined map $p \colon X \to S_{\infty}$ which takes $z \in X \smallsetminus (L_1 \cup L_2)$ to the unique intersection of its leaf with $S_\infty$, and takes $z \in L_i$ to the intersection with $S_{\infty}$ of the broken leaf asymptotic to the unique geodesic through $z$ representing the foliation class $\beta_i$. %The map is well defined by positivity of intersection since $S_{\infty}$ is complex.
The images $p(L_1)$ and $p(L_2)$ are closed, embedded, disjoint curves in $S_{\infty}$.

\end{itemize}

This picture can also be refined by {\it straightening} the leaves of $\mathcal{F}_{L_1,L_2}(J)$ near $L_1 \cup L_2$ following 
Proposition 5.16 of \cite{rgi} and Lemma 3.4 of \cite{hiker}.

\begin{lemma} Perturbing $J$ away from  $S_\infty \cup T_\infty \cup\,\mathcal{U}_1 \cup \,\mathcal{U}_2$ we may assume that the unbroken leaves of 
 $\mathcal{F}_{L_1,L_2}(J)$ that intersect $ \mathcal{U}_1 \cup \mathcal{U}_2$ do so along the annuli $$\{p_1=\delta, q_1= \theta, -\epsilon <p_2<\epsilon\}$$ for some $\theta \in S^1$ and $\delta \in (-\epsilon, \epsilon).$  As well, the planes of the broken leaves of $\mathcal{F}_{L_1,L_2}(J)$ through  $S_\infty$ intersect each $\mathcal{U}_i$ along the annulus $$\{p_1 = 0, q_1 =
\theta , 0 < p_2 < \epsilon \},$$ for some $ \theta \in S^1,$  and the planes of the broken leaves of $\mathcal{F}_{L_1,L_2}(J)$ through  $S_0$ intersect each $\mathcal{U}_i$ along the annulus $$\{p_1 = 0, q_1 =
\theta , -\epsilon < p_2 < 0 \},$$ for some $ \theta \in S^1.$
\end{lemma}

\begin{remark}\label{rem:reg}
 Using the regularity assumption on $J_\infty$ we can also assert that each sphere and plane that appears in a leaf of $\mathcal{F}_{L_1,L_2}(J)$ is the image of regular somewhere injective $J_{\infty}$-holomorphic curve.
\end{remark}

\begin{remark}\label{rem:flex}
The assumption that both $L_1$ and $L_2$ are symplectomorphic to $L(r,s)$ is not necessary here. The argument, in particular the use of Corollary \ref{>r}, goes through without change if $L_1$ is replaced by  $L(r, s-1+t)$ for any $t \in [1, b-s+1).$ This fact is crucial to the proof of Proposition \ref{prop:proj} below.
\end{remark}



\subsection{Completion of the proof of Proposition \ref{prop:proj}}
Choose an $\Omega$-tame almost complex structure $J$ on $X$ that is standard in some open neighborhood $U_\infty$ of $S_{\infty} \cup T_{\infty}$, is compatible with the parameterizations of $L_1$ and $L_2$, and whose limit, $J_{\infty}$, is regular for somewhere injective curves. By Lemma \ref{lem:noB}, we have the foliation $\mathcal{F}_{L_1,L_2}(J)$ of $X$ relative to $L_1 \cup L_2$. The properties of its projection map $p \colon X \to  S_{\infty}$, as described above in (F4), imply all but the final assertion of Proposition \ref{prop:proj}.  It remains to verify that the component of $S_\infty \smallsetminus p(L_1)$ containing $p(T_\infty)$,  is contained in the component of $S_\infty \smallsetminus p(L_2)$ containing $p(T_\infty)$. To prove this
we will use foliations of $X$ relative to $L[t]$ and $L[t] \cup L_2$  where $$L[t]= L(r, s-1+t)$$ and $t \in [1, b -s +1).$
The crucial fact to be exploited is that, by Assumption \ref{assump}, $L_2$ is disjoint from all of the $L[t]$.

\bigskip

\subsubsection{Stretching in a family.}
Let $\epsilon>0$ be small enough so that for all $t \ge b-s+1 -\epsilon$ the torus $L[t]$ lies in the neighborhood  $U_\infty$ where $J$ is standard. Setting $c-b-s+1-\epsilon$, let $\{J_t\}$ for $t \in [1,c]$ be a family of $\Omega$-compatible almost complex structures such that 
\begin{itemize}
    \item $J_1=J$,\\
    \item each $J_t$ is equal to $J$ in $U_\infty.$\\
    \item each $J_t$ is compatible with the parameterizations of $L[t]$ and $L_2$.
\end{itemize}

\medskip
In what follows, we will either stretch the $J_t$ along $L[t] \cup L_2$ to get a limiting almost complex structure $J_{t, \infty}$, or we will stretch the $J_t$ along just $L[t]$ to get a limiting almost complex structure $J'_{t, \infty}$. 

\begin{definition}\label{regulart}
 We say that $t \in [1,c]$ is regular if $J_{t, \infty}$ and $J'_{t, \infty}$  are both regular for somewhere injective curves.
\end{definition}



If $t$ is regular then we get a foliation of $X$ relative to $L[t] \cup L_2$, which we denote by $\mathcal{F}(J_t)$, and 
a foliation of $X$ relative to $L[t]$ which we denote by $\mathcal{F}'(J_t)$. We denote the corresponding projection maps by $p_t \colon X \to S_\infty$ and $p'_t \colon X \to S_\infty$, respectively.

\medskip


We may assume that $t=1$ and $t=c$ are regular and that the set of  $t \in (1,c)$ which are not regular is finite.
We now show that on intervals of regular $t$ the families of foliations $\mathcal{F}(J_t)$ and $\mathcal{F}'(J_t)$ vary smoothly.





% In particular,  we cannot guarantee that all $J_t$ are of foliation type. 
% However there are still limiting foliations of $X \smallsetminus (L[t] \cup L_2)$, or $X \smallsetminus L[t]$, which we denote by $\mathcal{F}(J_t)$, or $\mathcal{F}'(J_t)$ respectively. The corresponding projections are denoted $p_t, p'_t : X \to S_{\infty}$, so $p=p_1$ and we set $p'=p'_1$.


\medskip



\begin{lemma}\label{lem:fam}
Suppose that $t \in (1,c)$ is regular. For sufficiently small $\delta>0$,  the points $\tau \in (t - \delta, \, t+\delta)$ are all regular and the corresponding foliations $\mathcal{F}(J_\tau)$ vary smoothly over this interval. 
\end{lemma}

\begin{proof}
We use Fukaya's trick for moving regular curves with Lagrangian boundary as the boundary moves. Given $\delta > 0$, consider a family of diffeomorphisms $f_{\tau}$ of $X$, for $\tau \in [t - \eta, \, t+\eta]$, 
such that 
\medskip

\begin{itemize}
    \item $f_{\tau}(L[\tau]) =L[t]$ \\
    \item $f_{\tau} \circ \psi_{\tau} = \psi_{t}$ \\
    \item $f_{\tau}$ agrees with the identity map in a neighborhood of $S_{\infty}\cup T_{\infty} \cup L_2$ \\
    \item $\|f_{\tau}\|_{C^1}$ is of order $1$ in $\tau -t$.\\
\end{itemize}
Consider the family of almost complex structures on $X$ given by $$\widetilde{J}_{\tau} = (f_{\tau})_*J_{\tau}.$$ By the second and third conditions these are compatible with the parameterizations $\psi_{t}$ on $L[t]$ and $\psi_2$ on $L_2$, and their stretched limits $$\widetilde{J}_{\tau, \infty} = (f_{\tau})_*J_{\tau,\infty} $$  on $X \smallsetminus (L[t] \cup L_2)$ are also adapted to these parameterizations.\footnote{Each $f_{\tau}$ can also be viewed as a diffeomorphism from $X \smallsetminus (L[\tau] \cup L_2)$ to $ X \smallsetminus (L[t] \cup L_2)$.}

By automatic regularity for holomorphic spheres and holomorphic planes, see \cite[Theorem 1]{we}, for sufficiently small $\delta>0$ our $J_{t, \infty}$-holomorphic foliation deforms to give $\widetilde{J}_{\tau, \infty}$-holomorphic relative foliations $\mathcal{F}(\widetilde{J}_\tau)$ of $X$ relative to $L[t] \cup L_2$ for $\tau \in (t-\delta, t+\delta)$. Moreover, (the curves of) these foliations vary smoothly with $\tau$. The map $f^{-1}_{\tau}$ then takes the relative foliations $\mathcal{F}(\widetilde{J}_\tau)$ to the relative foliation $\mathcal{F}(J_\tau)$ of $X$ relative the $L[\tau] \cup L_2$.
\end{proof}

An identical argument yields the following. 

\begin{lemma}\label{lem:fam'}
Suppose that $t \in (1,c)$ is regular. For sufficiently small $\delta>0$,  the points $\tau \in (t - \delta, \, t+\delta)$ are all regular and the corresponding foliations $\mathcal{F}'(J_\tau)$ vary smoothly over this interval. 
\end{lemma}


To manage the behavior of these relative foliations at the  nonregular points in $[1,c]$ it is helpful to shift from the language of foliations to the language of moduli spaces of pseudo-holomorphic curves. Abusing notation, we will denote by $\mathfrak{r}_{0,t}$ the moduli space of $J_{t, \infty}$-holomorphic planes asymptotic to $L[t]$ which are disjoint from $S_{\infty}$, have Maslov class $2$ and area $r$. If $t$ is regular,  then $\mathfrak{r}_{0,t}$ is an $S^1$-family of classes each represented by a plane asymptotic to a geodesic representing negative the foliation class of $L[t]$. Moreover, each class in $\mathfrak{r}_{0,t}$ has a representative plane that is one half of a broken leaf of the foliation $\mathcal{F}(J_t)$. 
%In general, these planes lie in $X \smallsetminus (S_{\infty} \cup T_{\infty})$ and so their relative homology class is determined by the conditions above. 
% Next we define 
% $\mathfrak{r}_{\infty,t}$ to be the moduli space of $J_{t, \infty}$-holomorphic planes asymptotic to $L[t]$ which have area $r$ and Maslov class $2$, intersect $S_{\infty}$ exactly once, and whose boundaries represent the  foliation class of $L[t]$. Again these conditions specify the relative homology class of the planes representing classes in $\mathfrak{r}_{\infty,t}$, and in the case when $J_t$ is of foliation type we get an $S^1$-family of classes each of which is represented by a plane that is one half of a broken leaf of the foliation $\mathcal{F}(J_t)$. 
% 
The moduli space 
$\mathfrak{r}_{\infty,t}$ is defined analogously, as are the moduli spaces $\mathfrak{s}_{0,t}$ and $\mathfrak{s}_{\infty,t}$ of planes asymptotic to $L_2$. If we stretch just along $L[t]$ then we have moduli spaces $\mathfrak{r}'_{0,t}$ and $\mathfrak{r}'_{\infty,t}$. Finally, we can define family moduli spaces $\tilde {\mathfrak{r}}_0$, $\tilde {\mathfrak{r}}_{\infty}$, $\tilde {\mathfrak{s}}_0$, $\tilde {\mathfrak{s}}_{\infty}$, $\tilde {\mathfrak{r}}'_0$, $\tilde {\mathfrak{r}}'_{\infty}$, where
$$\tilde {\mathfrak{r}}_0= \{(u,t) \mid u \in \mathfrak{r}_{0,t},\, t\in [1,c]\},$$
and the other spaces are defined analogously.

\bigskip

%Given these preliminaries, we proceed with the proof of Proposition \ref{prop:proj}.
% We may assume that all but finitely many $J_{t, \infty}$ are regular for somewhere injective curves.
%, and the following lemma says that the corresponding foliation varies smoothly near such $t$.


%To deal with compactness issues we invoke Lemma \ref{>r2}.

% \begin{lemma}\label{>r2}
% Let $L \subset (Y^4, \sigma)$ be a Lagrangian torus in a symplectic $4$-manifold. Given an almost complex structure $J$ on $Y$, compatible with $L$, set
% $$a(J) = \min\{  \sigma(C) \, | \, \mathrm{index}(C) \ge \mathrm{end}(C) \}$$
% where the minimum is taken over $J_{\infty}$-holomorphic curves asymptotic to $L$, and $\mathrm{end}(C)$ denotes the number of ends of $C$.

% Then the numbers $a(J)$ are independent of $J$. In particular, given almost complex structures $J_0$ and $J_1$, all $J_1$-holomorphic Maslov $2$ planes asymptotic to $L$ have area at least $a(J_0)$.

% Moreover, in a generic family $\{J_t\}$, universal moduli spaces of curves satisfying $\mathrm{index}(C) \ge \mathrm{end}(C)$ and $\sigma(C) = a(J_0)$ are compact.
% \end{lemma}

% \begin{proof}
% This exactly follows \cite[Lemma 3.7]{ho}, which implies the same result for $Y^4 = \RR^4$ and a product torus $L(1,x)$. The idea is argue by contradiction and to consider a universal moduli space of curves corresponding to a generic family $\{J_t\}$, and representing the class of a $J_1$-holomorphic curve $C_1$ with $\sigma(C_1) < a(J_0)$. We note that automatic regularity holds for curves satisfying $\mathrm{index}(C) \ge \mathrm{end}(C)$, so these universal moduli spaces are manifolds. The moduli space may be noncompact, but \cite[Lemma 3.7]{ho} shows that any degenerations will result in a $J_t$-holomorphic curve $C_t$, again with $\mathrm{index}(C_t) \ge \mathrm{end}(C_t)$ and $\sigma(C_t) < a(J_0)$. Repeating the argument on the interval $[0,t]$, eventually we arrive at a contradiction to the definition of $a(J_0)$.

% This works generally, the assumption on $Y$ and $L$ in \cite[Lemma 3.7]{ho} was used only to calculate $a(J_0) =1$ when $J_0$ is the standard product structure.
% \end{proof}


% Without assuming the regularity  $J_{t, \infty}$ for a fixed $t$, we can,  in some cases, still understand limits of $J_{t_k, \infty}$ curves, or $J'_{t_k, \infty}$ curves, as $t_k$ converges to $t$.

Lemma \ref{>r2} can be used to show that three of these family moduli spaces are compact.




\begin{lemma}\label{lem:fam2} Let $t_k$ be a regular sequence converging to a $t \in [1,c]$ that is not regular.

\begin{enumerate}

\item %Suppose $J_{t_k, \infty}$ are a sequence of almost complex structures which are stretched along $L_{t_k}$ but smooth elsewhere. 
Suppose that $u_k$ is a sequence of $J'_{t_k, \infty}$-holomorphic planes in $\mathfrak{r}'_{0,t_k}$ ($\mathfrak{r}'_{\infty,t_k}$). Any subsequence of $u_k$ which converges to a holomorphic building must in fact converge to a single  $J'_{t, \infty}$-holomorphic plane in $X \smallsetminus L[t]$. Hence, the family moduli spaces $\tilde {\mathfrak{r}}'_0$ and $\tilde {\mathfrak{r}}'_{\infty}$ are compact.  %In other words, the moduli space of $J_{t, \infty}$-holomorphic planes in the classes $\mathfrak{r}_0$ and $\mathfrak{r}_{\infty}$ are compact.

\item %Suppose $J_{t_k, \infty}$ are a sequence of almost complex structures which are stretched along $L_{t_k} \cup L_2$.
Suppose that $v_k$ is a sequence of $J_{t_k, \infty}$-holomorphic planes in $\mathfrak{s}_{0,t_k}$. Any subsequence of $v_k$ which converges to a holomorphic building must in fact converge to a single  $J_{t, \infty}$-holomorphic plane in $X \smallsetminus L_2$.
Hence, the family moduli space $\tilde {\mathfrak{s}}_0$ is compact.
%In other words, the moduli space of planes in the class $\mathfrak{s}_0$ is compact.
Moreover, for each $t$, there exists a unique plane in $\mathfrak{s}_{0,t}$ asymptotic to each geodesic in $L[t]$ in the negative  of the foliation class.
\end{enumerate}
\end{lemma}

\begin{proof}


It follows from the compactification results of \cite{behwz} that to prove the first assertion  it suffices to show that the only possible top level curve in the limiting building is a plane of $\Omega$-area r.  With this, Lemma \ref{>r2} implies that it suffices to show that for some $\Omega$-compatible almost complex structure $J_0$ compatible with $L[t] =L(r, s-1+t)$, the minimal $\Omega$-area of $J_0$-holomorphic curves $u$ with $\mathrm{index}(u) \ge \mathrm{end}(u)$ is $r$. This is easily checked when $J_0$ is the standard product complex structure.

\medskip

To prove (2), we first note that since $L_2$ is Hamiltonian isotopic to the product torus $L(r,s)$ in $\RR^4$, there exists an almost complex structure on $\RR^4$ for which the minimal area of holomorphic curves is $r$. Then, since $X \setminus (S_{\infty} \cup T_{\infty})$ can be identified with a subset of $\RR^4$, we see from Lemma \ref{>r2} that curves $v_k$ cannot degenerate to nontrivial buildings that are broken only along $L_2$.

Suppose then that a limiting building, $\mathbf{B}$, of the curves $v_k$ is broken along $L[t]$. In other words $\mathbf{B}$ contains curves in $T^* L[t]$. Since the $t_k$ are regular, each $J_{t_k}$ is of foliation type and the curves in $\mathfrak{s}_{0,t_k}$ belong to the limiting foliation $\mathcal{F}(J_{t_k})$. In particular, they are disjoint from the planes in $\mathfrak{r}_{0,t_k}$ and $\mathfrak{r}_{\infty,t_k}$. It then follows, from positivity of intersection, that the curves of $\mathbf{B}$ mapping to $T^* L[t]$ must cover cylinders in the foliation class of $L[t]$. Given such a curve, the building $\mathbf{B}$ can be divided into two nontrivial connected sub-buildings, one of which has the property that all its unmatched ends cover geodesics in multiples of the foliation class in $L[t]$. This sub-building will have area at least $r$. Since the other sub-building is notrivial it will have positive area. As the total area of $\mathbf{B}$ is $r$ we have a contradiction.
\end{proof}

Now we can complete the proof of Proposition \ref{prop:proj}. Arguing by contradiction, we assume that 
\begin{equation}\label{assumpa}
p(L_2) \subset C \smallsetminus p(T_\infty).    
\end{equation}
We note that $J'_{t,\infty}$ agrees with $J_{t,\infty}$ away from a neighborhood of $L_2$. Therefore any curves of $\mathcal{F}(J_t)$ which avoid a neighborhood of $L_2$, and in particular are unbroken along $L_2$, will also form leaves of the foliation $\mathcal{F}'(J_t)$.


%It will be useful to view the projection $p$ as part of the outcome of the following two step process. In the first step, we stretch  $J_t$ along $L[t]$ to get $J'_{t,\infty}$ on $X \smallsetminus L[t]$. For regular $t$, this yields a projection map $p'_t \colon X \to S_\infty$ in terms of the relative foliation $\mathcal{F}'(J_t).$ In the second step, we stretch the $J'_{t,\infty}$ along $L_2$ to get $J_{t,\infty}$ on $X \smallsetminus L[t]\cup L_2$. For regular $t$, this yields the projection maps $p_t \colon X \to S_\infty$ defined in terms of the relative foliations $\mathcal{F}(J_t).$ In these terms, we have $p = p_1$. 

Our choice of $c$ and the condition that $J$ is standard in $U_\infty$ implies that  $p'_t(T_\infty) = p_t(T_\infty) = p_1(T_\infty)$ for all $t \in [1,c],$ and that $p'_c(L[c])$ is the boundary of a small disk $D_c$ in $S_\infty$ containing $p_1(T_\infty)$. Choosing  $c$ sufficiently close to $b-s+1$ we can refine assumption \eqref{assumpa} to 
\begin{equation}\label{assumpb}
p_1(L_2) \subset C \smallsetminus D_c.    
\end{equation}

%In step 2, the almost complex structure $J'_{1,\infty}$ is only changed in a neighborhood of $L_2$. Since this does not affect the broken planes along $L_1$, 
Then, as the broken leaves over $\partial C$ are disjoint from $L_2$, we also have
\begin{equation}\label{bound}
    p_1'(L[1])= p_1(L[1]) = \partial C
\end{equation} and hence 
\begin{equation}\label{bound1}
    (p'_1)^{-1}(C)= (p_1)^{-1}(C).
\end{equation}
Similarly, as our almost complex structures remain standard on $U_{\infty}$ we have  
\begin{equation}\label{bound2}
    (p'_t)^{-1}(D_c)= (p_t)^{-1}(D_c),
\end{equation}
for all $t$.

Formulas \eqref{bound1} and \eqref{bound2}, together with Assumption \eqref{assumpb}, imply that
\begin{equation}\label{assumpc}
p'_1(L_2) \subset C; \, \, p'_c(L_2) \subset S_{\infty} \smallsetminus D_c.    
\end{equation}
In fact we also have
\begin{equation}\label{assumpc2}
p'_1(\mathfrak{s}_{0,1}) \subset C; \, \, p'_c(\mathfrak{s}_{0,c}) \subset S_{\infty} \smallsetminus D_c.    
\end{equation}
as the broken $J_{1,\infty}$ or $J_{c,\infty}$ holomorphic curves asymptotic to $L_2$ are disjoint from those asymptotic to $L[1]$ and $L[c]$ respectively (in the first case by the assumption that $J_1$ is regular, and in the second case as $J_t=J$ is standard on $U_{\infty}$).


By Lemmas \ref{lem:fam} and \ref{lem:fam2} the moduli spaces $\mathfrak{s}_{0,t}$ vary continuously with $t$. Using a fixed geodesic on $L_2$ in the foliation class, we can identify a continuous family of planes $\{w_t\}_{t \in [1,c]}$ with $[w_t] \in  \mathfrak{s}_{0,t}$. Let $\{ x_t \}_{t \in [1,c]}$ be a continuous family of points such that $x_t$ is in the image of $ w_t$. By \eqref{assumpc2} we have $$x_1 \in (p'_1)^{-1}(C), \, \, x_c \in (p'_c)^{-1}(S_{\infty} \smallsetminus D_c).$$


Lemmas \ref{lem:fam'} and \ref{lem:fam2} imply that the moduli spaces $\mathfrak{r}'_{0,t}$ and $\mathfrak{r}'_{\infty,t}$ also vary continuously with $t$. Taking paths of representative planes for each and invoking  \eqref{bound} and \eqref{bound2}, we obtain a homotopy between $(p_1')^{-1}(\partial C)$ and $(p_c')^{-1}(\partial D_{c})$. By \eqref{bound2}, the homotopy avoids $(p_1)^{-1}(D_c)$, and so gives a surjection onto $(p'_1)^{-1}(C \smallsetminus D_c)$. Therefore, for the path $\{ x_t \}_{t \in [1,c]}$ there is an $s \in [1,c]$, and a $J'_{s, \infty}$-holomorphic plane representing either $\mathfrak{r}'_{0, s}$ or $\mathfrak{r}'_{\infty, s}$, that intersects $x_s$.


Let $J^N_{t, \infty}$ be the result of stretching $J'_{t, \infty}$ to length $N$ along $L_2$ so that $J^N_{t, \infty} \to J_{t,\infty}$ as $N \to \infty$. Arguing as above, for a fixed $N$, we can find an $s^N \in [1,c]$, and a $J^N_{s^N, \infty}$-holomorphic plane representing  $\mathfrak{r}'_{0, s^N}$ or $\mathfrak{r}'_{\infty, s^N}$ whose image intersects $x_{s^N}$. Passing to a subsequence, we may assume that the $s^N$ converge to $s^{\infty} \in [1,c]$ and that the limits of the $J^N_{s^N, \infty}$-holomorphic planes representing  $\mathfrak{r}'_{0, s^N}$ and $\mathfrak{r}'_{\infty, s^N}$ both converge to $J_{s^{\infty}, \infty}$-holomorphic buildings. One of these limiting buildings, $\bf{B}$, must intersect $x_{s^{\infty}}$ and thus the $J_{s^{\infty}, \infty}$-holomorphic plane $w_{s^{\infty}}$ that contains $x_{s^{\infty}}$ and represents a class in $\mathfrak{s}_{0,s^\infty}$. Either the building $\bf{B}$ includes a curve that covers the plane $w_{s^{\infty}}$, or it does not.

If $\bf{B}$ does include such a curve, this curve will have area at least $r$. However, the building $\bf{B}$ has a total area $r$ and must also include a curve asymptotically to $L[{s^{\infty}}]$ of positive area. This gives a contradiction.

If $\bf{B}$ does not include a curve that covers the plane $w_{s^{\infty}}$, then, by positivity of intersection, it includes a curve that intersects $w_{s^{\infty}}$ positively. However, this would imply that, for large $N$, the planes $w_{s^N}$ representing $\mathfrak{s}_{0,s^N}$  must intersect the planes in either $\mathfrak{r}_{0, s^N}$ or $\mathfrak{r}_{\infty, s^N}$. This is a contradiction whenever the $J^N_{s^N}$ are of foliation type.






% Then the limits $J^{\infty}_{t, \infty}$ are a family of almost complex structures stretched along $L[t] \cup L_2$. These almost complex structures have corresponding moduli spaces $\mathfrak{s}_{0,t}$, which, by Lemmas \ref{lem:fam} and \ref{lem:fam2}(2) also vary smoothly with $t$. 
% %Assuming our $J^{\infty}_{t, \infty}$ are constant on the complement of $p^{-1}(C)$, this complement remains foliated by holomorphic spheres, and the curves in $\mathfrak{s}_{0,t}$ lie in $p^{-1}(C \smallsetminus D_{\epsilon})$ by positivity of intersection.
% We fix a geodesic on $L_2$. This determines a $J^{\infty}_{t, \infty}$-holomorphic plane $w_t$ in $\mathfrak{s}_{0,t}$ (unique by Lemma \ref{lem:fam2}, part 2) and we can also fix a family of points $x_t$ in the image of $w_t$.
% The points $x_t$ yield $s^N$ as above, such that there exists a $J^N_{s^N, \infty}$-holomorphic plane in $\mathfrak{r}'_{0, s^N}$ or $\mathfrak{r}'_{\infty, s^N}$ with image intersecting $x_{s^N}$. Taking a subsequence, we may assume $s^N \to s^{\infty} \in [1,c]$.

% Finally we take a limit as $N \to \infty$. Limits of the $J^N_{s^N, \infty}$-holomorphic planes in $\mathfrak{r}'_{0, s^N}$ and $\mathfrak{r}'_{\infty, s^N}$ together give a $J^{\infty}_{s^{\infty}, \infty}$-holomorphic building. One of these limiting buildings $\bf{B}$ intersects $x_{s^{\infty}}$ and hence $w_{s^{\infty}}$. There are two cases.

% In the first case a curve of $\bf{B}$ covers $w_{s^{\infty}}$. Then the curve has area at least $r$. But the building $\bf{B}$ has total area $r$, and also includes a curve asymptotic to $L_{s^{\infty}}$. This gives a contradiction.

% In the second case, by positivity of intersection, a curve of $\bf{B}$ intersects $w_{s^{\infty}}$ positively. But then planes $w_{s^N}$ for $N$ large will also intersect the planes in either $\mathfrak{r}_{0, s^N}$ or $\mathfrak{r}_{\infty, s^N}$. This is a contradiction when the $J^N_{s^N}$ are of foliation type.



% %First fix a family of $J_t$ which are compatible with the $L[t]$ for $t \in [1, c]$ where $c = b-s+1 - \epsilon$ for $\epsilon$ small enough for $L_c$ to lie in the neighborhood of $T_{\infty}$ where our almost complex structures are standard. We can stretch these almost complex structures along the $L[t]$ (but not yet along $L_2$) and obtain corresponding foliations 
% First we stretch the along just $L[t]$. By Lemmas \ref{lem:fam} and \ref{lem:fam2}(1), the corresponding moduli spaces $\mathfrak{r}'_{0,t}$ and $\mathfrak{r}'_{\infty,t}$ vary continuously with $t$, and give a homotopy between $(p')^{-1}(\partial C)$ and $(p')^{-1}(\partial D_{\epsilon})$, where $D_{\epsilon}$ is a small disk around $\{ \infty \} \in S_{\infty}$.




% Argue by contradiction and assume $p(L_2) \subset C \setminus D_{\epsilon}$. Then %we may choose $J_1$ such that 
% also $p'(L_2) \subset C$, where $p'$ is defined with respect to $J'_{\infty}$. Indeed, changing our almost complex structure in a neighborhood of $L_2$ does not affect the broken planes along $L_1$, and hence does not change the set $p^{-1}(C)$. Similarly, we may assume $p_1(L_2) \notin D_{\epsilon}$, where $p_1$ is the projection with respect to an almost complex structure stretched along $L_c$.

% Our homotopy above gives a surjection onto $p^{-1}(C \smallsetminus D_{\epsilon})$. Therefore, for any point, or family of points $\{ x_t \} \in X$  with $x_0, x_1 \in p^{-1}(C \smallsetminus D_{\epsilon})$, we can find a corresponding $s \in [1,c]$, and a $J'_{s, \infty}$-holomorphic plane in $\mathfrak{r}'_{0, s}$ or $\mathfrak{r}'_{\infty, s}$, with image intersecting $x_s$.

% Let $J^N_{t, \infty}$ be the result of stretching $J'_{t, \infty}$ to length $N$ along $L_2$.
% Then the limits $J^{\infty}_{t, \infty}$ are a family of almost complex structures stretched along $L[t] \cup L_2$. These almost complex structures have corresponding moduli spaces $\mathfrak{s}_{0,t}$, which, by Lemmas \ref{lem:fam} and \ref{lem:fam2}(2) also vary smoothly with $t$. 
% %Assuming our $J^{\infty}_{t, \infty}$ are constant on the complement of $p^{-1}(C)$, this complement remains foliated by holomorphic spheres, and the curves in $\mathfrak{s}_{0,t}$ lie in $p^{-1}(C \smallsetminus D_{\epsilon})$ by positivity of intersection.
% We fix a geodesic on $L_2$. This determines a $J^{\infty}_{t, \infty}$-holomorphic plane $w_t$ in $\mathfrak{s}_{0,t}$ (unique by Lemma \ref{lem:fam2}, part 2) and we can also fix a family of points $x_t$ in the image of $w_t$.
% The points $x_t$ yield $s^N$ as above, such that there exists a $J^N_{s^N, \infty}$-holomorphic plane in $\mathfrak{r}'_{0, s^N}$ or $\mathfrak{r}'_{\infty, s^N}$ with image intersecting $x_{s^N}$. Taking a subsequence, we may assume $s^N \to s^{\infty} \in [1,c]$.

% Finally we take a limit as $N \to \infty$. Limits of the $J^N_{s^N, \infty}$-holomorphic planes in $\mathfrak{r}'_{0, s^N}$ and $\mathfrak{r}'_{\infty, s^N}$ together give a $J^{\infty}_{s^{\infty}, \infty}$-holomorphic building. One of these limiting buildings $\bf{B}$ intersects $x_{s^{\infty}}$ and hence $w_{s^{\infty}}$. There are two cases.

% In the first case a curve of $\bf{B}$ covers $w_{s^{\infty}}$. Then the curve has area at least $r$. But the building $\bf{B}$ has total area $r$, and also includes a curve asymptotic to $L_{s^{\infty}}$. This gives a contradiction.

% In the second case, by positivity of intersection, a curve of $\bf{B}$ intersects $w_{s^{\infty}}$ positively. But then planes $w_{s^N}$ for $N$ large will also intersect the planes in either $\mathfrak{r}_{0, s^N}$ or $\mathfrak{r}_{\infty, s^N}$. This is a contradiction when the $J^N_{s^N}$ are of foliation type.

\section{The proof of Proposition \ref{prop:submanifolds}}\label{dbuildings}

The symplectic disks $D^1_{\infty}$ and $D^2_{\infty}$ in the statement of Proposition \ref{prop:submanifolds} are obtained by compactifying a plane in $\mathfrak{r}_\infty$ and a plane in $\mathfrak{s}_\infty$, respectively. The remaining submanifolds in the statement of Proposition \ref{prop:submanifolds}  will be obtained using $(d,n_1)$-buildings.  These are defined and analyzed below in the next section. Two existence results for $(d,n_1)$-buildings are established in Section \ref{sub:exist}. The process of obtaining  the desired symplectic submanifolds from these buildings is described in Section \ref{sub:shift}


% Here we prove Propositions \ref{limitF} and \ref{limitG} in a more general context. Let $L_1$ and $L_2$ be Lagrangian tori in $X =S^2(2r) \times S^2(b) \smallsetminus S_{\infty}$ where $S_{\infty} = \{\infty\} \times S^2(b)$. Viewing $X$ as the compactification of $P(2r,b)$ we assume that both $L_1$ and $L_2$ are Hamiltonian isotopic to $L(r,s)$.  As above, the relative foliations $\mathcal{F}_{L_1,L_2}(J)$ of $X \smallsetminus L_1 \cup L_2$ exist, and choosing $J$ to be standard near $S_{\infty}$, $T_0$ and $T_{\infty}$, we may assume that $T_0$ and $T_{\infty}$ are leaves of $\mathcal{F}_{L_1,L_2}(J)$
% and that $L_1, L_2 \subset X \smallsetminus (S_{\infty} \cup T_0 \cup T_{\infty})$.  For the projection $p \colon X \to S_{\infty}$ the images $p(L_1)$ and $p(L_2)$ are circles as in Lemma \ref{compts}.
% To simplify our calculations we now inflate $S_{\infty} \subset X$ by adding a tubular neighborhood of capacity $2r-1$. To simplify notation, we denote the resulting symplectic manifold by $(X, \Omega)$. This manifold can viewed as the compactification of $P(2r,b).$
% Note that $L_1 = L(r,s)$ and $L_2 = \phi(L(r,s))$ are still contained in $X$, the almost complex structure $J$ is $\Omega$-tame and the properties of the foliation $\mathcal{F}_{L_1,L_2}(J)$ are unchanged by the inflation except that each of its leaves has $\Omega$-area  equal to $2r$. In particular, the broken leaves of $\mathcal{F}_{L_1,L_2}(J)$ now consist of two planes, each with $\Omega$-area equal to $r$. Moreover, any $J_\infty$-holomorphic curve that intersect $S_\infty$ has $\Omega$-area at least $2r-1$.

\subsection{On $(d,n_1)$-buildings and their types}\label{sub:d}
Fix a nonnegative integer $d$ and a generic collection of $2d+1$ points on $L_1 \cup L_2$. Let $n_1 \le 2d+1$ be the number of these points on $L_1$. For a generic choice of tame almost complex structure $J$ on $(X,\Omega)$ (compatible with parameterizations of $L_1$ and $L_2$ as before) there exists a smooth $J$-holomorphic sphere $u\colon S^2 \to X$ that represents the class $(d,1)$ and passes through the $2d+1$ points on $L_1 \cup L_2$. This curve is unique up to reparameterization. 
% To see this, note that for the integrable product complex structure such a curve can be written explicitly as the graph of a degree $d$ rational map, and this implies that the Gromov-Witten invariant associated to the homology class and point constraints is $1$. Hence, nodal curves will exist for all tame almost complex structures and away from a  codimension $2$ subset of almost-complex structures we will have smooth curves. The uniqueness in the assertion above follows because spheres in the class $(1,d)$ have self intersection number $2d$, so distinct spheres cannot satisfy the same $2d+1$ point constraints.
As one  stretches along  $L_1 \cup L_2$ to obtain the foliation $\mathcal{F}_{L_1,L_2}(J)$ one gets a sequence of pseudo-holomorphic spheres in the class $(d,1)$. Passing to a convergent subsequence, one obtains a holomorphic building, $\mathbf{F}$, which we will refer to as a $(d,n_1)$-building. As before, this building consists of genus zero holomorphic curves in three levels.
% It consists of genus zero holomorphic curves in three levels; {\em top level} curves of $\mathbf{F}$ mappig to $X \smallsetminus (L_1 \cup L_2)$, {\em middle level} curves mapping to one of two copies of the symplectization of the unit cotangent bundle of the flat torus that correspond to 
% $L_1$ and $L_2$, and {\em bottom level} curves mapping to one of two copies of $T^*\mathbb{T}^2$.  
Each top level curve of $\mathbf{F}$ can be compactified to yield a map from a surface of genus zero with boundary to $(X,L_1 \cup L_2)$. The components of the boundary correspond to the negative punctures of the curve.  They are mapped to the closed geodesics on $L_1$  or $L_2$ underlying the Reeb orbits to which the corresponding puncture is asymptotic. The middle and bottom level curves of $\mathbf{F}$ can be compactified to yield maps to either $L_1$ or $L_2$ with the same type of boundary conditions. These compactified maps can all be glued
together to form a map $\bar{\mathbf{F}}\colon S^2  \to X$ in the class $(d,1)$. %The image of $\bar{\mathbf{F}}$  is an embedded symplectic sphere.

\medskip

The curves of a $(d,n_1)$-building are constrained by its companion foliation $\mathcal{F}_{L_1,L_2}(J)$. To describe these constraints, the following notions will be useful.


% The following notions will be useful here.

\begin{definition}\label{sp}
Given an almost complex structure $J$ on $(X,\Omega)$,
%For a fixed relative foliation $\mathcal{F}_{L_1,L_2}(J)$
a $J_\infty$-holomorphic curve $u$ in $X \smallsetminus (L_1 \cup L_2)$ is said to be {\em essential} if the map $p \circ u$ is injective. Here $p$ is the projection map defined using the foliation $\mathcal{F}_{L_1,L_2}(J)$.
\end{definition}


\begin{definition}\label{ft}
 %For a fixed relative foliation $\mathcal{F}_{L_1,L_2}(J)$, a 
 A puncture of a $J_\infty$-holomorphic curve $u$ in $X \smallsetminus (L_1 \cup L_2)$ is said to be of {\em foliation type with respect to} $L_i$  if it is asymptotic to a closed Reeb orbit which lies on the copy of $S^*\mathbb{T}^2$ that corresponds to $L_i$ and covers a closed geodesic in an integer multiple of the foliation class $\beta_i$. The puncture is of {\em positive (negative) foliation type} if this integer is positive (negative).
\end{definition}


% \begin{definition}\label{fdef:fibersphere}
% A {\em leaf sphere} is a $J_\infty$-holomorphic sphere $\tau \colon S^2 \to X \smallsetminus (L_1 \cup L_2)$ which covers an unbroken leaf of $\mathcal{F}_{L_1,L_2}(J)$. 
% \end{definition}


% \begin{definition}\label{fdef:brokenfibersphere}
% A {\em broken leaf sphere on $L_1$} is a finite collection $\widehat{\tau} =\{\rho_j\}$ of $J_\infty$-holomorphic planes  $\rho_j \colon \CC \to X \smallsetminus (L_1 \cup L_2)$ such that each $\rho_j$ covers one of a pair of planes in $\mathfrak{r}_0 \cup \mathfrak{r}_\infty$ that form a broken leaf of $\mathcal{F}_{L_1,L_2}(J)$, and the compactification of $\widehat{\tau}$ defines a map of the sphere that covers this broken leaf. A {\em broken leaf sphere on $L_2$} is defined analogously.
% \end{definition}


% \begin{definition}\label{fdef:brokenfibersphere}
% A {\em positive (negative) chain of planes along $L_1$} is a finite collection of $J_\infty$-holomorphic planes each of which covers a plane in $\mathfrak{r}_\infty$ ($\mathfrak{r}_0$). A {\em positive (negative) chain of planes on $L_2$} is defined analogously.
% \end{definition}
% \bigskip

\begin{proposition}\label{prop:type} Each  $(d,n_1)$-building  $\mathbf{ F}$ %for the relative foliation $\mathcal{F}_{L_1,L_2}(J)$ 
is of one of the following types:

%\bigskip

% \noindent {\bf Type 0.} $\mathbf{ F}$ has a unique essential curve, $\underline{u}$, which is a $J_\infty$-holomorphic sphere in $X \smallsetminus (L_1 \cup L_2)$ in the class $(d-k,1)$ for some $0 \le k \le d-1$. Any additional top level curves of $\mathbf{ F}$ are either spheres covering unbroken leaves the foliations or collections of planes covering broken leaves of the foliation. Together these additional curves compactify to form a possibly disconnected curve in the class $(k,0)$. 
% Any middle or bottom level curves cover cylinders asymptotic to Reeb orbits in multiples of the relevant foliation class.

% The number of planes in the broken leaf spheres along $L_i$ is at least $d_i +2$.

\bigskip


\noindent {\bf Type 1.} $\mathbf{ F}$ has a unique essential curve, $\underline{u}$, which is a $J_{\infty}$-holomorphic sphere with at least one puncture. The image of the injective map $p\circ \underline{u}$ is $S_{\infty}$ minus finitely many points on $p(L_1) \cup p(L_2)$. Any punctures of $\underline{u}$ asymptotic to $L_1$ are of foliation type and are either all positive or all negative. The same holds for any punctures of $\underline{u}$ along $L_2$. 

% For every positive (negative) puncture of $\underline{u}$ on $L_1$ there is a corresponding negative (positive) chain of planes along $L_1$ in $\mathbf{ F}$. The analogous statement holds for punctures of $\underline{u}$ asymptotic to $L_2$. The remaining top level curves of $\mathbf{ F}$ are either leaf spheres or planes that comprise a set of broken leaf spheres along $L_1$ and $L_2$. 

%These top level holomorphic planes are all either in $\mathfrak{s}_0 \cup \mathfrak{r}_0$ or they are in $\mathfrak{s}_{\infty} \cup \mathfrak{r}_{\infty}$.
%of the same class and are all asymptotic to Reeb orbits covering geodesics in the same class, say $m\beta$.   
%The middle and bottom level curves of $\mathbf{ F}$ are all covers of cylinders of the form $v^{\mathbf{ \theta}, \mathbf{ m}, \mathbf{ p}}$ with $\mathbf{ m}=(0, \pm m)$.



\bigskip


\noindent {\bf Type 2.1.} $\mathbf{ F}$ has exactly two essential  curves, $u_2$ and $\underline{u}$. The closures of the images of the maps $p \circ u_2$ and $p \circ \underline{u}$  are $A$ and $B \cup C$, respectively. Any punctures of $\underline{u}$ on $L_1$ are all of foliation type and are either all positive or all negative. 

% For every positive (negative) puncture of $\underline{u}$ on $L_1$ there is a corresponding negative (positive) chain of planes along $L_1$ in $\mathbf{ F}$. 

%Moreover, $u^0_{L_1}$ and $u^{\infty}_{L_1}$ are both $J$-holomorphic planes asymptotic to closed Reeb orbits over geodesics in classes of the form $(+1, k)$ and $(-1, l)$. 
% The building $\mathbf{ F}$ may also include leaf spheres or planes that comprise a set of broken leaf spheres along $L_1$ and $L_2$. Any remaining top level curves of  $\mathbf{ F}$ cover planes in $\mathfrak{s}_0 \cup \mathfrak{s}_\infty$.




% The other top level curves of $\mathbf{ F}$  cover (broken or unbroken) leaves of $\mathcal{F}_{L_1,L_2}(J)$. Any middle and bottom level curves in the copy of $T^* \mathbb{T}^2$ corresponding to $L_1$ cover cylinders asymptotic to Reeb orbits in multiples of the foliation class.


\bigskip


\noindent {\bf Type 2.2.} $\mathbf{ F}$ has exactly two essential curves, $\underline{u}$ and $u_1$. The closures of the images of the maps $p \circ \underline{u}$ and $p \circ u_1$  are $A \cup B$ and $C$, respectively.  Any punctures of $\underline{u}$ on $L_2$ are all of foliation type and are either all positive or all negative. 

% For every positive (negative) puncture of $\underline{u}$ on $L_2$ there is a corresponding negative (positive) chain of planes along $L_2$ in $\mathbf{ F}$. 
%Moreover, $u^0_{L_1}$ and $u^{\infty}_{L_1}$ are both $J$-holomorphic planes asymptotic to closed Reeb orbits over geodesics in classes of the form $(+1, k)$ and $(-1, l)$. 
% The building $\mathbf{F}$ may also include leaf spheres or planes that comprise a set of broken leaf spheres along $L_1$ and $L_2$. Any remaining top level curves of  $\mathbf{F}$ cover planes in $\mathfrak{r}_0 \cup \mathfrak{r}_\infty$.






% Any punctures of $\underline{u}$ on $L_2$ are all of foliation type and are either all positive or all negative. Corresponding to every such puncture along $L_2$ is a plane of $\mathbf{ F}$ covering a plane of a broken leaf. This plane belongs to $\mathfrak{s}_0$ if the punctures are all positive and $\mathfrak{s}_\infty$  if they are negative.
% %Moreover, $u^0_{L_1}$ and $u^{\infty}_{L_1}$ are both $J$-holomorphic planes asymptotic to closed Reeb orbits over geodesics in classes of the form $(+1, k)$ and $(-1, l)$. 
% The other top level curves of $\mathbf{ F}$  cover (broken or unbroken) leaves of $\mathcal{F}_{L_1,L_2}(J)$. Any middle and bottom level curves in the copy of $T^* \mathbb{T}^2$ corresponding to  $L_2$ cover cylinders asymptotic to Reeb orbits in multiples of the foliation class.


\bigskip


\noindent {\bf Type 3.} $\mathbf{ F}$ has exactly three essential curves, $u_1$, $\underline{u}$,  and $u_2$. The closures of the images of the maps  $p \circ u_1$, $p \circ \underline{u}$,  and $p \circ u_2$ are  $C$, $B$ and $A$, respectively. 
% The building $\mathbf{F}$ may also include leaf spheres or planes that comprise a set of broken leaf spheres along $L_1$ and $L_2$. Any remaining top level curves of  $\mathbf{F}$ cover planes in $\mathfrak{r}_0 \cup \mathfrak{r}_\infty \cup \mathfrak{s}_0 \cup \mathfrak{s}_\infty$.




% The other top level curves of $\mathbf{ F}$  again cover (broken or unbroken) leaves of $\mathcal{F}_{L_1,L_2}(J)$. 

\end{proposition}

\begin{proof} This is a refinement of Proposition 3.17 of \cite{hiker}. In that work, the constraint points were not restricted to lie on the relevant Lagrangian tori. This allows for an additional kind of $(d,n_1)$-buildings, Type 0, which is characterized by having a unique essential curve which is a $J_\infty$-holomorphic sphere, without punctures, that maps to $X \smallsetminus (L_1 \cup L_2)$ and represents the class $(d-k,1) \in H^2(X;\ZZ)$ for some $0 \le k \le d-1$.  To prove Proposition \ref{prop:type} it  suffices to show that if $\mathbf{F}$ is a  $(d,n_1)$-building corresponding to $2d+1$ constraint points on $L_1 \cup L_2$, then it cannot be of Type 0. Assume that $\mathbf{F}$ is of Type 0. Then its nonessential curves can be sorted into: a collection of $J_\infty$-holomorphic spheres that cover unbroken leaves of the foliation,
and a collection of sub-buildings that cover broken leaves of the foliation. Together, these nonessential curves compactify to form a disconnected curve in the class $(k,0)\in H^2(X;\ZZ)$. In particular, the number of sub-buildings covering broken leaves is at most $k \le d$. On the other hand, the middle and bottom level curves of every one of these sub-building cover a unique cylinder, in the foliation class, in one of the two copies of $T^*\mathbb{T}^2$ associated to $L_1$ or $L_2$. Hence, each of the $k \le d$ sub-buildings can pick up, at most, one of the generic point constraints. Since there are $2d+1$ of these constraints  and no other way for $\mathbf{F}$ to hit them, we have arrived at a contradiction and $\mathbf{F}$ cannot be of Type 0.

\end{proof}




\subsubsection{Managing the nonessential curves of $(d,n_1)$-buildings}\label{common}

Here we describe a scheme for managing the nonessential curves of $(d,n_1)$-buildings of Types $1$, $2.1$ and $2.2$. In what follows we will denote the collection of essential curves of a $(d,n_1)$-building $\mathbf{F}$ by $\mathbf{u}$.


\bigskip

\noindent {\bf Type $1$.}
Let $\mathbf{F}$ be a $(d,n_1)$-building of Type $1$. Then $\mathbf{u} = \{\underline{u}\}$, where $\underline{u}$ is the essential curve of $\mathbf{F}$ described in Proposition \ref{prop:type}. The image of $p \circ \mathbf{u}$ is equal to $S_{\infty}$ minus a finite collection of points on $p(L_1) \cup p(L_2)$.
The following dichotomy will be useful: the image of every nonessential curve of $\mathbf{F}$ either projects to a point in the image of $p\circ \mathbf{u}$, or it maps to one of the missing points. 


The curves whose projections lie in the image of $p\circ \mathbf{u}$, are either $J_\infty$-holomorphic spheres that cover unbroken leaves of the foliation $\mathcal{F}_{L_1,L_2}(J)$, or they are $J_\infty$-holomorphic planes that form a collection of sub-buildings  of $\mathbf{F}$ that each cover a broken leaf of the foliation. Let $N$ be the total number of times these spheres and sub-buildings cover the class $(1,0)$. Then $N=N_0+N_1 + N_2$ where $N_0$ is the number of times an unbroken leaf is covered, and $N_i$ is the number of times a broken leaf on $L_i$ is covered.

The remaining curves of $\mathbf{F}$, whose images don't project to the image of $p\circ \mathbf{u}$, each belong to a sub-building defined by one of the punctures of $\mathbf{u}$. Label these punctures as  $$\{x_1, \dots, x_{k_1}, y_1, \dots, y_{k_2}\}.$$ Here, each $x_i$ is asymptotic to a closed geodesic on $L_1$ of the form $\gamma_i^{\ell_i}$, where $\gamma_i$ represents the foliation class $\beta_1$ and all the $\ell_i$ have the same sign. Similarly, each $y_j$ is asymptotic to a closed geodesic on $L_2$ of the form $\sigma_j^{m_j}$ where $[\sigma_j]=\beta_2$ and all the $m_j$ have the same sign. In these terms, the image of $p\circ \mathbf{u}$ is equal to $S_\infty$ minus the set  $$\{p(\gamma_1), \dots, p(\gamma_{k_1})\}\cup\{ p(\sigma_1), \dots, p(\sigma_{k_2})\} \subset p(L_1) \cup p(L_2).$$ 

The relevant features of the sub-building corresponding to the puncture $x_i$ are as follows. It compactifies to a disk and has a single unmatched end on $\gamma_1^{\ell_i}$. Its top level curves all cover one of the planes in a single matching pair of planes in $\mathfrak{r}_0 \cup \mathfrak{r}_\infty$. Let $M^0_1(i)$ be the degree of the map from the sub-building to the plane in $\mathfrak{r}_0$ and let $M^\infty_1(i)$  be the degree of the map from the sub-building to the plane in $\mathfrak{r}_\infty.$ We then have 
\begin{align}
 \ell_i+M^{\infty}_1(i)-M^0_1(i)=0.   
\end{align}
The description of the sub-buildings corresponding to the $y_j$ is completely analogous, and yields nonnegative numbers $M^0_2(j)$ and $M^\infty_2(j)$ satisfying
\begin{align}
 m_j+M^{\infty}_2(j)-M^0_2(j)=0.   
\end{align} 
Set  $$M_1^0 = \sum_{i=1}^{k_1} M^0_1(i),$$ and define the numbers $M_1^\infty$, $M_2^0$, and $M_2^\infty$ analogously. 

With these labels in place we can derive some equations and inequalities  which will be useful later on. First we note that if the punctures of $\mathbf{u}$ along $L_1$ are all positive then each $\ell_i$ is at least one, and we have 
\begin{align}\label{pos1}
M_1^0 = M_1^\infty + \sum_{i=1}^{k_1} \ell_i \geq M_1^\infty + k_1.
\end{align}
If these punctures are all negative we have 
\begin{align}\label{neg1}
M_1^\infty = M_1^0 - \sum_{i=1}^{k_1} \ell_i \geq M_1^0 + k_1.
\end{align}
The fact that there are  $n_1$ point constraints on $L_1$ implies that 
\begin{align}\label{constraint1}
k_1 + N_1 \geq n_1. 
\end{align}
Computing the $\Omega$-area of the curves of  $\mathbf{F}$ we get
\begin{align}\label{area1}
\Omega (\mathbf{u})+ r(M_1^0 +M_1^\infty + M_2^0 + M_2^\infty)+2rN = 2rd+b.  
\end{align}
Finally, computing the intersection number of the curves of $\mathbf{F}$ with $S_{\infty}$ we get the expression
\begin{align}\label{S1}
\mathbf{u} \cdot S_{\infty}+ M_1^\infty + M_2^\infty + N  =  d.
\end{align}


\noindent {\bf Type $2.1$.}
Let $\mathbf{F}$ be a $(d,n_1)$-building of Type $2.1$. In this case  $\mathbf{u} = \{u_2,\underline{u}\}$. The numbers $N$, $N_0$, $N_1$ and $N_2$ can be defined as above. Labeling the punctures of $\underline{u}$  on $L_1$ as  $\{x_1, \dots, x_{k_1}\}$, each $x_i$ is again asymptotic to a closed geodesic on $L_1$ of the form $\gamma_i^{\ell_i}$, where $[\gamma_i] =\beta_1$ and all the $\ell_i$ have the same sign. For each puncture $x_i$ there is a corresponding sub-building of $\mathbf{F}$ that compactifies to a disk whose boundary maps to $\gamma_1^{\ell_i}$. The top level curves of this subuilding yield degrees $M^0_1(i)$ and $M^\infty_1(i)$ as before,  satisfying
\begin{align}\label{eq:2.1match}
 \ell_i+M^{\infty}_1(i)-M^0_1(i)=0,   
\end{align}
and hence inequality  \eqref{pos1} or \eqref{neg1}. To accommodate the point constraints on $L_1$ we again have 
\begin{align}\label{constraint2}
k_1 + N_1 \geq n_1. 
\end{align}


The remaining top level components of $\mathbf{F}$ cover planes in $\mathfrak{s}_0$ and $\mathfrak{s}_\infty$. Let $\widetilde{M}^0_2$ be the total degree of these curves that  cover planes in $\mathfrak{s}_0$ and let  $\widetilde{M}^\infty_2$ be the total degree of these curves covering planes in $\mathfrak{s}_\infty$. The fact that there are $n_2 = 2d+1-n_1$ point constraints on $L_2$ implies that 
\begin{align}\label{constraints2.1}
 \widetilde{M}^0_2+\widetilde{M}^\infty_2   \geq n_2 - 1.
\end{align}
This follows from an index argument. Indeed, the deformation index of curves in $T^*\mathbb{T}^2$ is $2s-2$, where $s$ is the number of positive ends, so we must have $2s-2 \geq 2n_2$. Since two ends can match the essential curves, the rest must match planes in $\mathfrak{s}_0$ or $\mathfrak{s}_\infty$, and the inequality follows.

In these terms, the computation of the $\Omega$-area of the curves of  $\mathbf{F}$ yields
\begin{align}\label{area2}
\Omega (\mathbf{u})+ r(M_1^0 +M_1^\infty + \widetilde{M}_2^0 + \widetilde{M}_2^\infty)+2rN = 2rd+b,  
\end{align}
and the intersection number of the curves of $\mathbf{F}$ with $S_{\infty}$ yields the expression
\begin{align}\label{S2}
\mathbf{u} \cdot S_{\infty}+ M_1^\infty + \widetilde{M}_2^\infty + N  =  d.
\end{align}

\noindent {\bf Type $2.2$.}
Let $\mathbf{F}$ be a $(d,n_1)$-building of Type $2.2$. In this case  $\mathbf{u} = \{\underline{u}, u_1\}$ and the discussion is completely analagous to that for buildings with Type 2.1, with the role of the puncture $x_i$ on $L_1$ replaced by the role of the puncture $y_j$ on $L_2$  and all the other labels swapped. In particular, we have
\begin{align}
 m_j+M^{\infty}_2(j)-M^0_2(j)=0,   
\end{align}
\begin{align}\label{constraint2.2}
k_2 + N_2 \geq n_2 ,
\end{align}
\begin{align}\label{area2.2}
\Omega (\mathbf{u})+ r(\widetilde{M}_1^0 +\widetilde{M}_1^\infty + M_2^0 + M_2^\infty)+2rN = 2rd+b,  
\end{align}
and 
\begin{align}\label{S2.2}
\mathbf{u} \cdot S_{\infty}+ \widetilde{M}_1^\infty + M_2^\infty + N  =  d.
\end{align}

%DO WE NEED THIS?


\subsection{On $(d,n_1)$-buildings of high degree}\label{sub:exist}

In this section we establish two classification results for $(d,n_1)$-buildings for sufficiently large $d$. These are used in the next section to prove the existence of the remaining symplectic submanifolds from Proposition \ref{prop:submanifolds}. 

The following monotonicity lemma, established by Boustany in \cite{boust}, will be used several times to leverage the assumption that $d$ is large. 
\begin{lemma}\label{lem:monotonicity}(\cite{boust}, Lemma 5.6)
Let $S \subset  X \smallsetminus (L_1 \cup L_2)$ be a closed $J_{\infty}$-holomorphic sphere. There exists an $\epsilon = \epsilon(S)>0$ such that for every (possibly) punctured $J_{\infty}$- holomorphic sphere $u$ in $X \smallsetminus (L_1 \cup L_2)$ we have
\begin{align}\label{lem:mon} \Omega (u) \geq \epsilon (u \bullet S).\end{align}
    
\end{lemma}

% We start with the following, which substitutes the assumption in \cite{hiker} that both tori are disjoint from $\{z_1=0\}$.

Our first result concerns the case when all the constraint points are on one of the Lagrangian tori.


\begin{proposition}\label{prop:Fd}
For $d$ sufficiently large, every $(d,2d+1)$-building is of Type 2.2. Each such building, $\mathbf{F}$. consists of two essential planes, $\underline{u}$, and $u_1$, a single lower level curve which maps to $T^*L_1$ and has $2d+1$ (positive) punctures, and $2d-1$ top level planes in $\mathfrak{r}_0 \cup \mathfrak{r}_\infty$. All the curves of $\mathbf{F}$ are somewhere injective and hence regular.
\end{proposition}

\begin{proof}

\noindent The first step of the proof is to use Lemma \ref{lem:monotonicity} and the  labeling described  in Subsection \ref{common} to show that when $d$ is sufficiently large here are no  $(d,2d+1)$-buildings of Type 1 or Type 2.1.  

Arguing by contradiction we assume first that for an arbitrarily large  $d$ there is a $(d,2d+1)$-building $\mathbf{F}$ of Type $1$. Let $\underline{u}$ be the unique essential curve of $\mathbf{F}$.


\medskip

\noindent{\em Case 1.} Assume that the $k_1$ punctures of $\underline{u}$ along $L_1$ are all negative. Then \eqref{neg1} yields $$M_1^\infty \geq M_1^0 + k_1.$$ In the current setting,  \eqref{constraint1} implies that $$k_1 +N_1 \geq 2d+1.$$ Taken together, these two inequalities imply that $M^{\infty}_1+ N \geq 2d+1$ which  contradicts \eqref{S1}.

\medskip

\noindent{\em Case 2.} Assume that the $k_1$ punctures of $\underline{u}$ along $L_1$ are all positive. By \eqref{pos1} we have $$M_1^0 \geq M_1^\infty + k_1.$$ In this case, \eqref{area1} implies that 
\begin{align*}
\Omega(\underline{u})+ r(2M_1^\infty +k_1 + M_2^0 + M_2^\infty)+2r(N_0+N_1+N_2)
\leq 2rd+b  
\end{align*}
Together with the implication from  \eqref{constraint1} that $k_1 +N_1 \geq 2d     +1$ we get
\begin{align*}
\Omega(\underline{u})+ r(2M_1^\infty + M_2^0 + M_2^\infty) +rN_1 +2r(N_0 + N_2)  \leq b-r.  
\end{align*}
This implies that the numbers $\Omega(\underline{u})$,  $M_1^\infty$, $M_2^\infty$, and $N$, which implicitly depend on $d$, are uniformly bounded from above independently of $d$. By \eqref{S1} we must then have  \begin{align*} \lim_{d \to \infty} \underline{u} \cdot S_{\infty} = \infty.\end{align*}  By Lemma \ref{lem:monotonicity}, this contradicts the assertion above that $\Omega(\underline{u})$ is uniformly bounded from above independently of $d$. Cases 1 and 2 imply that for sufficiently large $d$ the building $\mathbf{F}$ cannot be of Type 1.

\medskip

Next, we assume first that for arbitrarily large $d$ there is a $(d,2d+1)$-building $\mathbf{F}$ of Type $2.1$. Let ${\bf{u}}_d=\{ u_2, \underline{u}\}$ be the essential curves of $\mathbf{F}$.

\medskip
 
 \noindent{\em Case 3.} Assume that the punctures of $\underline{u}$, and hence $\mathbf{u}_d$, along $L_1$ are all negative. By inequality \eqref{neg1} we have $M^{\infty}_1 \geq M^1_0 +k_1$. Together with \eqref{constraint2} this implies that $M^\infty_0 + N \geq 2d+1$ which contradicts \eqref{S2}.

\medskip

\noindent{\em Case 4.} Assume that  the punctures of $\underline{u}$ along $L_1$ are all positive. In this case, we have $M_1^0 \geq M_1^\infty + k_1 $ and  $k_1 +N_1 \geq 2d+1$. Together with \eqref{area2}, these imply that
% \begin{align*}
% \mathrm{area}(\underline{u})+ r(2M_1^\infty +n_1 + M_2^0 + M_2^\infty)+2r(N_0+N_1+N_2) 
% \leq 2re+b  
% \end{align*}
% and thus 
\begin{align*}
\Omega(\mathbf{u}_d)+ r(2M_1^\infty + \widetilde{M}_2^0 + \widetilde{M}_2^\infty)+rN_1 +2r(N_0 + N_2)  \leq b-r.  
\end{align*}
From this it follows that $\Omega(\mathbf{u}_d)$,  $M_1^\infty$, $\widetilde{M}_2^\infty$, and $N$ are uniformly bounded from above independently of $d$. With this, equation  \eqref{S2} implies that  \begin{align*}\lim_{d \to \infty} \underline{u} \cdot S_{\infty} = \infty.\end{align*} By Lemma \ref{lem:monotonicity}, this contradicts the assertion above that $\Omega(\bf{u}_d)$ is uniformly bounded from above independently of $d$. Cases 3 and 4 imply that, for sufficiently large $d$, the building $\mathbf{F}$ cannot be of Type 2.1.

\medskip


Finally, we use index computations to show that any  $(d,2d+1)$-building  cannot be of Type $3$. Let $\mathbf{F}$ be such a building.  Let $s_i$ be the total number of positive ends of the curves of $\mathbf{F}$ that map to $T^*L_i$. The total index of the curves mapping to $T^*L_i$ is at most $2s_i -2$, where the upper bound corresponds uniquely to the case of a single curve with $s_i$ positive ends. 

To accommodate the $2d+1$ constraint points on $L_1$ we must have \begin{equation}\label{eq:s1}
    s_1  \geq 2d+2.
\end{equation}
Since two of these $s_1$ ends match with essential curves of $\mathbf{F}$, there must be $s_1 -2$ addition curves of $\mathbf{F}$ corresponding to each of the remaining ends. These curves can be included in planar sub-buildings that cover curves in $\mathfrak{r}_0 \cup \mathfrak{r}_\infty$. Similarly, 
there are also $s_2 -2$ planar sub-buildings of $\mathbf{F}$ that cover curves in $\mathfrak{s}_0 \cup \mathfrak{s}_\infty$.  %Since planes have odd indices

By the index bound above for curves in $T^*L_i$, the total index of the curves of $\mathbf{F}$ (defined to be the sum of the indices of the individual curves minus any matching conditions) is at most \begin{equation}\label{eq:ind3}
\begin{split}
   \mathrm{index}(\mathbf{u}) + (2s_1-2) + (2s_2-2)+ \\
   \sum_{i=1}^{s_1-2}\mathrm{index}(\mathbf{v_i}) + \sum_{j=1}^{s_2-2}\mathrm{index}(\mathbf{w_j}) - (s_1 +s_2). 
\end{split}
\end{equation}
The first term, $\mathrm{index}(\mathbf{u})$, is the total index of the three essential curves, the next two terms are upper bounds for the index of the curves in $T^*L_i$, the two sums represent the total index of the planar sub-buildings mapping into $\mathfrak{r}_0 \cup \mathfrak{r}_\infty$ and $\mathfrak{s}_0 \cup \mathfrak{s}_\infty$ respectively, and the last term %, $-(s_1 +s_2)$, 
comes from the matching constraints at Reeb orbits for the Morse-Bott Reeb flow of the flat metric on $\mathbb{T}^2$. 
As planar sub-buildings, like holomorphic planes, have odd index, we have
\begin{equation}\label{eq:ind35}
   \sum_{i=1}^{s_1-2}\mathrm{index}(\mathbf{v_i}) \geq s_1 -2; \, \, \,  \sum_{j=1}^{s_2-2}\mathrm{index}(\mathbf{w_j}) \geq s_2 -2.
\end{equation}
Since the total index is preserved under limits, the total index of $\mathbf{F}$ is equal to $2(2d+1)$, and it then follows from \eqref{eq:s1}, \eqref{eq:ind3} and  \eqref{eq:ind35} that 
\begin{equation}\label{eq:ind4}
   \mathrm{index}(\mathbf{u}) \leq 2(3-s_2). 
\end{equation}
The three essential curves of $\mathbf{F}$ are $u_1$, $\underline{u}$,  and $u_2$. The curves $u_1$ and $u_2$ are holomorphic planes and hence have odd indices. Since $s_2 \ge 2$ , we must  then have $s_2=2$, $\mathrm{index}(u_1) = \mathrm{index}(u_2) =1$, and 
$\mathrm{index}(\underline{u})=0$. Furthermore, all our inequalities become equalities, and hence $s_1 = 2d+2$. For a generic choice of $J$, %after the constraint points are fixed, this is impossible. There 
this means the building $\mathbf{F}$ has a single curve in $T^*L_1$, and this curve is rigid (that is, the curve has deformation index $0$ when constrained to intersect the $2d+1$ points). In turn, this means the constraint points determine a finite collection of possible positive ends, and these ends will not, generically,  match with the ends of the rigid curve $\underline{u}$.

\medskip

It remains for us to show that the limit $\mathbf{F}$, of Type 2.2, consists entirely of regular curves.  Arguing as above, the total index of the curves of $\mathbf{F}$ is at most \begin{equation}\label{eq:ind5}
\begin{split}
   \mathrm{index}(\mathbf{u}) + (2s_1-2) + 2s_2 + \\
   \sum_{i=1}^{s_1-2}\mathrm{index}(\mathbf{v_i}) + \sum_{j=1}^{s_2}\mathrm{index}(\mathbf{w_j}) - (s_1 + 2s_2). 
\end{split}
\end{equation} 
The third term, $2s_2$, corresponds to curves of $\mathbf{F}$ covering cylinders in $T^*L_2$, and the next two terms represent the total index of the planar sub-buildings as above with ends on $L_1$ and $L_2$ respectively, and the  
%The fifth term, $s_2$, corresponds to planes of $\mathbf{F}$ covering planes in $\mathfrak{s}_0 \cup \mathfrak{s}_\infty$. The 
final term corresponds to matching condition constraints.

The analogue of \eqref{eq:ind35} is
\begin{equation}\label{eq:ind36}
   \sum_{i=1}^{s_1-2}\mathrm{index}(\mathbf{v_i}) \geq s_1 -2; \, \, \,  \sum_{j=1}^{s_2}\mathrm{index}(\mathbf{w_j}) \geq s_2 -2.
\end{equation}

Since the total index of the curves of $\mathbf{F}$ is equal to $2(2d+1)$, it follows from \eqref{eq:s1}, \eqref{eq:ind5} and \eqref{eq:ind36}  that 
\begin{equation}\label{eq:ind7}
   \mathrm{index}(\mathbf{u}) \leq 2-s_2. 
\end{equation}
The two essential curves of $\mathbf{F}$ are $\underline{u}$ and $u_2$. Since  $u_2$ is a holomorphic plane it has odd index and we have \begin{equation}\label{eq:ind6}
   \mathrm{index}(\underline{u}) \leq 1-s_2. 
\end{equation} 
Therefore $s_2  \le 1$, and if $s_2=1$ then the essential curve $\underline{u}$ is rigid. In this case our various inequalities must be equalities, in particular $s_1=2d+2$ and there is a single curve in $T^*L_1$ which is also rigid. However for generic $J$ we do not expect to find rigid curves with matching asymptotics. 
%The curve $\underline{u}$ must have an end asymptotic to $L_1$ and hence must have positive index. 
Thus, $s_2 =0$ and the inequalities \eqref{eq:ind7} and \eqref{eq:ind6}, and hence \eqref{eq:ind36}, are equalities. This means that our planar sub-buildings are in fact single, somewhere injective holomorphic planes, and so all the curves of $\mathbf{F}$ are somewhere injective as required.


% The image of every nonessential curve of $\mathbf{F}$ either projects to a point in the image of $p\circ \underline{u}$ or to one of the missing points. 
 
 
 
%  In this setting $M_1$ be the number of planes in the chains of planes along $L_1$ that correspond to the punctures of $\underline{u}$ along $L_1$. Recall that $M_1=M_1^{0}$ if these punctures are positive and $M_1=M_1^\infty$ if they are negative.
 
%  % Define $N$ and $\widehat{N}$ as in the previous cases. The remaining top level curves of $\mathbf{F}$ are all $J_\infty$-holomorphic planes that cover planes in $\mathfrak{s}_0 \cup \mathfrak{s}_\infty$. Let $M_2^*$ be the number of these planes, counted with multiplicity, that cover planes in $\mathfrak{s}_*$.
% % nd the closures of the images of $p \circ u_2$ and $p \circ \underline{u}$ are $A$ and $B \cup C$, respectively.  
% % Any punctures of $\underline{u}$ on $L_1$ are all of foliation type and are either all positive or all negative.
% %Moreover, $u^0_{L_1}$ and $u^{\infty}_{L_1}$ are both $J$-holomorphic planes asymptotic to closed Reeb orbits over geodesics in classes of the form $(+1, k)$ and $(-1, l)$. 
% % Hence the nonessential top level curves of $\mathbf{ F}_e$ that are asymptotic to $L_1$ are planes that either all cover planes that belong to $\mathfrak{r}_0$ or all cover planes that belong to $\mathfrak{r}_{\infty}$. Let $M_1$ be the total number of these planes counted with multiplicity. Let $M_2$ be the number of top level curves of $\mathbf{F}$ covering planes in  $\mathfrak{s}_0 \cup \mathfrak{s}_{\infty}$, $N_1$ the number of top level curves of $\mathbf{F}$ covering unbroken leaves, and $N_2$ the number of pairs of planes covering broken leaves. Again, these are all counted with multiplicity.

% Computing the area of $\mathbf{F}$ we get
% \begin{align}
%  \mathrm{area}(u_2) + \mathrm{area}(\underline{u})+ r(M_1+M_2+ \widehat{N})+2rN= 2re+b.
% \end{align}

% To accommodate the $2e+1$ point constraints on $L_1$ we must have 
% \begin{align}\label{constraint2}
% M_1 + \widehat{N} \geq 2e+2. 
% \end{align}


% Hence,
% \begin{align}\label{eq:nope2}
%  \mathrm{area}(u_2) + \mathrm{area}(\underline{u}) + rM_2 + 2rN + r\widehat{N} \leq b-r.  
% \end{align}

% Consider the intersections of $\mathbf{F}$ with $S_{\infty}$.
% If all the punctures of $\overline{u}$  asymptotic to $L_1$ are negative, then 
% \begin{align}
% u_2 \cdot S_{\infty}+ \underline{u}\cdot S_{\infty}+ M_1 + M_2^{\infty}+ N + \widehat{N} = e.
% \end{align}
% This contradicts \eqref{constraint2}. Hence all punctures of $\overline{u}$  asymptotic to $L_1$ are positive and we have 
% \begin{align}
% u_2 \cdot S_{\infty}+ \underline{u}\cdot S_{\infty} + M_2^{\infty}+ N + \widehat{N} = e.
% \end{align}
% By inequality \eqref{eq:nope2}, each of $M_2^\infty$, $N$ and $\widehat{N}$ is uniformly bounded from above independently of $e$ and so we must have \begin{align}\label{inf} \lim_{e \to \infty} \left( u_2 \cdot S_{\infty}+ \underline{u}\cdot S_{\infty} \right) = \infty.\end{align} Again Lemma \ref{lem:monotonicity} implies that, for sufficiently large $e$, this will lead to a violation of inequality \eqref{eq:nope2}. Thus, for sufficiently large $e$, the building $\mathbf{F}$ can not be of Type 2.1

\end{proof}

% \begin{proposition}\label{prop:fdreg}
% If $\mathbf{F}$ is a $(d,2d+1)$-building of Type 2.2 or Type 3, then each of its constituent curves is somewhere injective and hence regular.
% \end{proposition}

% \begin{proof}

% \medskip

% \noindent{\em Case 1.} Assume that $\mathbf{F}$ is of Type 2.2. In this case, $\mathbf{F}$ has two essential curves, $\underline{u}$ and $u_1$ such that the image of  $p \circ \underline{u}$ is $A \cup B$ and the image of $p \circ u_1 $ is  $C$.  By definition, these essential curves are injective. To prove that the remaining curves of $\mathbf{F}$ are somewhere injective, we will find a lower bound for the number of these curves and an upper bound for their total area. Together these bounds will imply that none of the nonessential curves of $\mathbf{F}$ can be multiply covered.

% We begin with the area bound. The total area of all the curves of $\mathbf{F}$ is $2dr +b$. We will show that the total area of   
% $\underline{u}$ and $u_1$ is at least $b= s+ (b-s)$ so that the remaining curves have total area at most $2dr$.


% \end{proof}

% TO PROVE THE NEXT COROLLARY WE NEED TO MOVE THE ENDS OF THE TOP LEVEL CURVES OF $\mathbf{F}$ AROUND USING FUKAYA'S TRICK (BEFORE GLUING THINGS TOGETHER). TO DO THIS, WE NEED THESE CURVES TO BE SOMEWHERE INJECTIVE AND HENCE REGULAR (NO?). HOW DO WE KNOW THIS? 

The primary implication of Proposition \ref{prop:Fd} is the following.


\begin{proposition}\label{subaxis} For sufficiently large $e \in \NN$, there exist embedded symplectic spheres $A_0, A_{\infty}, B_0, B_{\infty} \subset X \smallsetminus (L_1 \cup L_2)$ in the class $(e,1)$, with the following properties:
\begin{enumerate}
\item $A_0$ intersects each plane in $\mathfrak{r}_0$ once and is disjoint from the planes in $\mathfrak{r}_{\infty}$;
\item $A_{\infty}$ intersects each plane in $\mathfrak{r}_{\infty}$ once and is disjoint from the planes in $\mathfrak{r}_0$;
\item $B_0$ intersects each plane in $\mathfrak{s}_0$ once and is disjoint from the planes in $\mathfrak{s}_{\infty}$;
\item $B_{\infty}$ intersects each plane in $\mathfrak{s}_{\infty}$ once and is disjoint from the planes in $\mathfrak{s}_0$;
\item either $A_0$ and $A_{\infty}$ both intersect the planes in $\mathfrak{s}_0$ and are disjoint from the planes in $\mathfrak{s}_{\infty}$, or $A_0$ and $A_{\infty}$ both intersect the planes in $\mathfrak{s}_{\infty}$ and are disjoint from the planes in $\mathfrak{s}_0$;
\item either $B_0$ and $B_{\infty}$ both intersect the planes in $\mathfrak{r}_0$ and are disjoint from the planes in $\mathfrak{r}_{\infty}$, or $B_0$ and $B_{\infty}$ both intersect the planes in $\mathfrak{r}_{\infty}$ and are disjoint from the planes in $\mathfrak{r}_0$.
\end{enumerate}
\end{proposition}

\begin{proof} We will establish the existence of the spheres $A_0$ and $A_{\infty}$ . The existence of the spheres $B_0$ and $B_{\infty}$ follows from the same argument by exchanging the roles of $L_1$ and $L_2$.  The starting point for the argument is an $(e,2e+1)$-building, $\mathbf{F}$, of the Type 2.2 as in Proposition \ref{prop:Fd}. We will obtain $A_0$ and $A_{\infty}$ by  deforming $\mathbf{F}$ near $L_1$ and then compactifying and smoothing.  


To deform  $\mathbf{F}$ we will use the procedure developed to prove Proposition 3.24 in \cite{hiker}. Let $\mathbf{v} =(a,b)$ be a vector in $(-\epsilon, \epsilon)^2$ and let $L_1(\mathbf{v})$ be the Lagrangian torus in $\mathcal{U}_1$ defined by $\{p_1=a, p_2=b\}$. Let $u$ be a top level curve of $\mathbf{F}$. By Proposition \ref{prop:Fd}, $u$ is a regular plane asymptotic to $L_1$. With this,  Lemmas 3.27 and 3.29 of \cite{hiker} imply the following.

\begin{lemma}\label{lem:shift1}
 If $\|v\|$ is sufficiently small, and $|a|$ is sufficiently small with respect to $|b|$, then there is an almost complex structure $J_{\mathbf{v}}$ on $X \smallsetminus (L_1 \cup L_1(\mathbf{v})\cup L_2)$ and a $J_{\mathbf{v}}$-holomorphic curve $u(\mathbf{v})$ whose domain is identical to that of $u$ and whose asymptotics are identical to those of $u$ but with the ones on $L_1$ translated to  $L_1(\mathbf{v})$. 
\end{lemma}

For  $\mathbf{v}$ as in Lemma \ref{lem:shift1}, let $\mathbf{F}(\mathbf{v})$ be the building obtained by translating each top level curve $u$ of $\mathbf{F}$ to $u(\mathbf{v})$ and identifying the bottom level curves of of $\mathbf{F}$ with the same curves now formally mapping to a copy of $T^*\TT^2$  that corresponds to $L_1(\mathbf{v})$. 

Compactifying the building $\mathbf{F}(\mathbf{v})$ and smoothing the resulting map near $L_1(\mathbf{v})$, we get a symplectic  embedding $$F_e(\mathbf{v}) \colon S^2 \to X \smallsetminus (L_1 \cup L_2).$$ Arguing as in  Lemma 3.33 of \cite{hiker}, we have the following. 

\begin{lemma}\label{lem:shift2}
   Suppose that  $a \neq 0$. If $b<0$, then $F_e(\mathbf{v})(S^2)$ intersects each plane in $\mathfrak{r}_0$ once and is disjoint from the planes in $\mathfrak{r}_{\infty}$. If $b>0$, then $F_e(\mathbf{v})(S^2)$ is disjoint from the planes in $\mathfrak{r}_0$ and intersects each plane in $\mathfrak{r}_{\infty}$ once.
 \end{lemma}

For $a_1\neq 0$ and  $b_1<0$ as above, set $$A_0 = F_e(\mathbf{v})(S^2).$$
Similarly, choose $a\neq 0$, and  $b>0$, and set $$A_\infty = F_e(\mathbf{v})(S^2).$$ For these choices the symplectic spheres $A_0$ and $A_\infty$ satisfy conditions (1) and (2) of Proposition \ref{subaxis}, respectively. To verify property (5) we note that Proposition \ref{prop:Fd} implies that either $\underline{u} \bullet \mathfrak{s}_0 =1$ and  $\underline{u} \bullet \mathfrak{s}_\infty =0$ or $\underline{u} \bullet \mathfrak{s}_0 =0$ and  $\underline{u} \bullet \mathfrak{s}_\infty =1$. By the construction of $A_0$ and $A_{\infty}$ from $\mathbf{F}(\mathbf{v})$, the same alternative holds for the intersections of $A_0$ and $A_{\infty}$ with  $\mathfrak{s}_0 $ and  $\mathfrak{s}_\infty$.

\end{proof}

We now consider $(d,d)$-buildings. 


\begin{proposition}\label{prop:dd} For $d$ sufficiently large, every $(d,d)$-building is of Type 3. Each such building, $\mathbf{F}$, has three essential curves $u_2$, $\underline{u}$, and $u_1$ with areas $s$, $r$, and $b-s$, respectively. Moreover, the only top level curves of $\mathbf{F}$, other than its three essential curves, are $d-1$ broken planes in $\mathfrak{r}_0 \cup \mathfrak{r}_{\infty}$ and $d$ broken planes in $\mathfrak{s}_0 \cup \mathfrak{s}_{\infty}$. 
All the constituent curves of  $\mathbf{F}$ are somewhere injective and hence regular. 
\end{proposition}

% DON'T YET SEE HOW TO GET THE PRECISE PLANE COUNTS IN THE STATEMENT. I UNDERSTAND THAT THE NUMBER OF PLANES ON $L_1$ IS AT LEAST $d-1$ AND THE NUMBER OF PLANES ON $L_2$ IS AT LEAST $d$. BUT, USING AREA, THE POTENTIAL EXCESS NUMBER OF DISKS SEEMS TO BE BOUNDED FROM ABOVE BY $$(b+r - \Omega(\mathbf{u}) -2rN)/r.$$ THIS SEEMS TO ALLOW FOR REAL EXCESS SINCE WE DON'T BOUND $b$.
% (Here $\mathbf{u}$ is the collection of three essential curves.)


% \begin{proof} We rely on the spheres from Lemma \ref{subaxis}, but do not refer to $S_{\infty}$. Therefore, relabelling our families of broken planes if necessary, we may assume that $A_0$ and $A_{\infty}$ both intersect $\mathfrak{s}_{\infty}$, and that both $B_0$ and $B_{\infty}$ both intersect $\mathfrak{r}_{\infty}$. Also we may assume these spheres are holomorphic, and a monotonicity lemma applies to intersecting holomorphic curves.

% We fix $d$ points on $L_1$ and $d+1$ points on $L_2$ and consider J-holomorphic curves in $X$ in the class $(d,1)$ that pass through these points. Stretching the neck along $L_1 \cup L_2$,  we obtain a limiting building,  $\mathbf{C}$ as well as a relative foliation $\mathcal{F}_{L_1,L_2}(J)$ and its corresponding projection $p_{\infty} \colon X \to S_{\infty}$. 

% It suffices to prove that $\mathbf{C}$ can not be of Type 2 along $L_1$. Assume it is. Then  $\mathbf{C}$ contains exactly one essential plane $C_1$ asymptotic to $L_1$ and precisely $d$ broken planes asymptotic to $L_1$ that are either all in $\mathfrak{r}_0$ or all in  $\mathfrak{r}_{\infty}$.

% \medskip

% \noindent{\bf Case 1.} The broken planes of $\mathbf{C}$ asymptotic to $L_1$ all lie in $\mathfrak{r}_0$.

% \medskip

% \noindent{\it Case $1a$. $\mathbf{C}$ is of Type 2 along $L_2$ and the broken planes of $\mathbf{C}$ asymptotic to $L_2$ all lie in $\mathfrak{s}_0$.}

% In this case, $\mathbf{C}$ has exactly  $d+1$ broken planes asymptotic to $L_2$ and a single essential curve. The essential curve has area $$2dr + b - dr - (d+1)r = b-r$$ and, 
% by the conventions above,  none of the broken planes  of $\mathbf{C}$ asymptotic to $L_s$ intersect  $A_{\infty}$. Hence, the essential curve asymptotic to $L_2$ intersects $A_{\infty}$ $(d,1) \cdot (e,1) = d + e$ times. For sufficiently large $d$ it follows from monotonicity that this will be greater than $b-r$ which is a contradiction.

% \medskip

% \noindent{\it Case $1b$. $\mathbf{C}$ is of Type 2 along $L_2$ and the broken planes of $\mathbf{C}$ asymptotic to $L_2$ all lie in $\mathfrak{s}_{\infty}$.}

% \medskip

% In this case no broken planes of $\mathbf{C}$ asymptotic to $L_2$ intersects $B_0$. Hence, the essential curve  of $\mathbf{C}$ asymptotic to $L_2$ intersects $B_0$ exactly $d+e$ times. Arguing as above, this can not be true for sufficiently large $d$

% \medskip

% \noindent{\it Case $1c$. $\mathbf{C}$ is of Type 3 along $L_2$.}

% \medskip

% Here there are $d$ broken planes asymptotic to $L_2$ lying in $\mathfrak{s}_0 \cup \mathfrak{s}_{\infty}$. Again we divide this into two cases. First suppose there are at least $d/2$ broken planes in $\mathfrak{s}_0$. Then the broken planes have at most $d/2$ intersections with $A_{\infty}$ (coming from the planes in $\mathfrak{s}_{\infty}$). Hence our essential curve intersects $A_{\infty}$ at least $d/2 + e$ times, and this is enough to contradict monotonicity.

% If there are at least $d/2$ broken planes in $\mathfrak{s}_{\infty}$, then we count $d + d/2$ intersections with $A_0$ coming from planes in $\mathfrak{r}_0$ and $\mathfrak{s}_{\infty}$. This contradicts positivity of intersection when $d > 2e$.

% {\bf Case 2.} The broken planes asymptotic to $L_1$ all lie in $\mathfrak{r}_{\infty}$.

% Suppose the limit is also of Type 2 at $L_2$ with $d+1$ broken planes in $\mathfrak{s}_0$. Then no broken planes intersect $A_0$ and this contradicts monotonicity for the essential curve.

% Suppose the limit is also of Type 2 at $L_2$ with $d+1$ broken planes in $\mathfrak{s}_{\infty}$. Then the broken planes intersect $A_{\infty}$ a total of $2d+1$ times, contradicting positivity of intersection.

% Suppose the limit is also of Type 3 at $L_2$ with at least $d/2$ broken planes in $\mathfrak{s}_0$. Then the broken planes intersect $B_0$ at least $3d/2$ times, contradicting positivity of intersection when $d > 2e$. 

% Suppose the limit is also of Type 3 at $L_2$ with at least $d/2$ broken planes in $\mathfrak{s}_{\infty}$. Then the broken planes intersect $A_{\infty}$ at least $3d/2$ times, contradicting positivity of intersection when $d > 2e$. 

% In conclusion, we have established that the limit cannot be Type 2 along $L_1$. A similar argument shows that the limit cannot be Type 2 along $L_2$, and hence we have a Type 3 limit as required.
% \end{proof}




% Fix $e$ and consider holomorphic spheres in the class $(e,1)$ that pass through $2e+1$ points, which we fix to be on $L_1$. We stretch the neck along $L_1 \cup L_2$ we obtain a limiting building $\mathbf{A}$ is disjoint from $L_2$ and has Type 3 on $L_1$, that is, the limiting building contains exactly two essential planes asymptotic to $L_1$ and $2d$ ($2e$?) broken planes lying in $\mathfrak{r}_0 \cup \mathfrak{r}_{\infty}$.

% To see that the limit is disjoint from $L_2$ we can use genericity of the almost complex structure, as breaking along $L_2$ is a codimension 1 condition. (Or alternatively, if there are curves breaking along $L_2$, we can use a perturbation argument to find disjoint curves.)

% Arguing by contradiction, suppose the limit has Type 2. Then, given disjointness from $L_2$, the limiting building has a single essential curve asymptotic to $L_1$ and $2e+1$ broken planes in $\mathfrak{r}_0 \cup \mathfrak{r}_{\infty}$. These planes actually all lie in the same family, and by positivity of intersection (using $(e,1) \cdot [S_{\infty}] = e$) this family must be  $\mathfrak{r}_0$.

% The essential curve has area $$2er + b - (2e+1)r = b - r$$ but has intersection number $e$ with $S_{\infty}$. When $e$ is sufficiently large this contradicts our monotonicity lemma and establishes the claim.

% Next, since the essential curve is disjoint from $L_2$, and intersects leaves of the foliation at most once, it intersects either $\mathfrak{s}_0$ and $\mathfrak{s}_{\infty}$, but not both. Thus perturbing our Type 3 curve to get a smooth sphere gives our result as required.
% \end{proof}

% Now we can prove the main proposition.

% \begin{proposition} Let $d$ be sufficiently large. Then there exists a building in the class $(d,1)$ which is of Type 3 along both $L_1$ and $L_2$. The building has three essential curves, a plane asymptotic to $L_1$, a plane asymptotic to $L_2$, and a cylinder asymptotic to $L_1 \cup L_2$. There are $d-1$ broken planes in $\mathfrak{r}_0 \cup \mathfrak{r}_{\infty}$ and 
% $d$ broken planes in $\mathfrak{s}_0 \cup \mathfrak{s}_{\infty}$.
% \end{proposition}

\subsubsection{The proof of Proposition \ref{prop:dd}} The proof uses the symplectic spheres of Proposition \ref{subaxis} for a fixed  integer, say  $e$. By relabeling, if necessary, we may assume that $A_0$ and $A_{\infty}$ both intersect $\mathfrak{s}_{\infty}$, and that $B_0$ and $B_{\infty}$ both intersect $\mathfrak{r}_{\infty}$. We may also assume that these spheres are $J$-holomorphic. Indeed, the sphere $A_0$ is already holomorphic away from an $L_1(\mathbf{v_0})$, with $\mathbf{v_0} = (a_1,b_1)$ and $b_1<0$, and since the spheres is symplectic we can deform $J$ in a neighborhood of $L_1(\mathbf{v_0})$ to make $A_0$ holomorphic. We may assume this neighborhood is disjoint from $B_0$ and $B_{\infty}$, since they are both deformations of buildings with no ends on $L_1$. We may also assume the neighborhood of $L_1(\mathbf{v_0})$ is disjoint from $A_{\infty}$. This follows as in Lemma 3.35 from \cite{hiker}, using a continuation argument, as $L_1(\mathbf{v_0})$ may be chosen close to $\mathfrak{r}_0$ while $A_{\infty}$ is disjoint from the planes in $\mathfrak{r}_0$. As the deformation of $J$ making $A_0$ holomorphic does not affect $B_0$, $B_{\infty}$ or $A_{\infty}$, or indeed any of the families of planes $\mathfrak{r}_{0}$, $\mathfrak{r}_{\infty}$, $\mathfrak{s}_{0}$, $\mathfrak{s}_{\infty}$, we may repeat the argument for the other spheres.
As a result of this, the spheres $A_0$, $A_{\infty}$, $B_0$ and $B_{\infty}$ can each play the role of $S$ in the applications of Lemma \ref{lem:monotonicity} below.

\bigskip


\noindent{\bf Step 1.} First we show that $\mathbf{F}$ cannot be of Type 1 when $d$ is sufficiently large. Arguing by contradiction we assume that we can find a Type 1 building $\mathbf{F}$ for an arbitrarily large $d$. This building  has a single essential curve, $\underline{u}$,  and there are four cases to consider corresponding to the possibilities that the punctures of $\underline{u}$ along $L_1$ are all negative/positive, and the punctures of $\underline{u}$ along $L_2$ are all negative/positive. In each case, we will arrive at a contradiction when $d$ is large enough.

\bigskip

\noindent{\it{Case 1.}} Assume first that the punctures of $\underline{u}$ along $L_1$ are all negative, and the punctures of $\underline{u}$ along $L_2$ are all negative. This implies the collection of inequalities:
\begin{align*}
    &M_1^\infty \geq M_1^0 + k_1,\\
    &M_2^\infty \geq M_2^0 + k_2,\\
    &k_1 +N_1 \geq d,\\
    &k_2 +N_2 \geq d+1.
\end{align*}
From these it follows that 
$$M_1^\infty + M_2^\infty + N_1 + N_2 \ge 2d+1$$
which contradicts \eqref{S1}.


\bigskip

\noindent{\it{Case 2.}} If the punctures of $\underline{u}$ along $L_1$ are all negative, and the punctures of $\underline{u}$ along $L_2$ are all positive we have:
\begin{align*}
    &M_1^\infty \geq M_1^0 + k_1,\\
    &M_2^0 \geq M_2^\infty + k_2,\\
    &k_1 +N_1 \geq d,\\
    &k_2 +N_2 \geq d+1,
\end{align*}
and hence
$$M_1^\infty + M_2^0 + N_1 + N_2 \ge 2d+1.$$
However, taking the intersection of $\mathbf{F}$ with the sphere $B_0$ yields 
\begin{align*}
    \underline{u} \bullet B_0 + M_1^\infty+M_2^0 +N =d+e.
\end{align*}
Fixing $e$, these two formulas cannot both hold when $d$ is  sufficiently large. 



\bigskip

\noindent{\it{Case 3.}} Arguing similarly, if the punctures of $\underline{u}$ along $L_1$ are all positive, and the punctures of $\underline{u}$ along $L_2$ are all negative we have
$$M_1^0 + M_2^\infty + N_1 + N_2 \ge 2d+1.$$ For fixed $e$ and large $d$ this contradicts the equality
\begin{align*}
  \mathbf{F} \bullet A_0 = \underline{u} \bullet A_0 + M_1^0 +M_2^\infty +N =d+e.
\end{align*}


\bigskip

\noindent{\it{Case 4.}} Finally, if the punctures of $\underline{u}$ along $L_1$ and $L_2$ are all positive then we have 
\begin{align}\label{c41}
M_1^0 + M_2^0 + N_1 + N_2 \ge 2d+1. 
\end{align}
Computing areas we have
\begin{align*}
\Omega(\underline{u}) + r(M_1^0 +M_1^{\infty}+M_2^0 +M_2^{\infty}) + 2rN = 2dr + b
\end{align*}
and using \eqref{c41} this gives
\begin{align}\label{c42}
\Omega(\underline{u}) + r(M_1^{\infty}+ M_2^{\infty})
&\leq  b - r - rN.
\end{align}

From this we see that $N$ is uniformly bounded, independently of $d$, as is the total area of the essential curve $\underline{u}$ and all curves covering $\mathfrak{r}_{\infty}$ and $\mathfrak{s}_{\infty}$.
But, computing intersections with $A_{\infty}$, we have
\begin{align*}
  \mathbf{F} \bullet A_{\infty} = \underline{u} \bullet A_{\infty} + M_1^{\infty} +M_2^{\infty} +N =d+e
\end{align*}
and, by \eqref{c42}, for sufficiently large $d$ this contradicts Lemma \ref{lem:monotonicity}. 




% IS $S_0$ THE CORRECT SPHERE HERE?


% Together with \eqref{area1}, these inequalities imply that 
% \begin{align}\label{bound}
%     \Omega(\underline{u}) + 2r(M_1^\infty+M_2^\infty) +2rN_0 +r(N_1 + N_2) \leq b-r.
% \end{align}
% In particular, the values of $\Omega(\underline{u})$, $M_1^\infty$, $M_2^\infty$, and $N$ are bounded from above, independently of $d$.
% On the other hand, taking the intersection of $\mathbf{F}$ with the sphere $A_\infty$ yields 
% \begin{align}
%     \underline{u} \cdot A_\infty + M_1^\infty+M_2^\infty +N =d+e.
% \end{align}
% Since the last three terms on the left are bounded,, for all $d$ sufficiently large we must have \begin{align}\label{Ainf}\epsilon(A_\infty)\underline{u} \cdot A_\infty > b-r,\end{align} where $\epsilon(A_\infty)$ is the constant from Lemma \ref{lem:monotonicity}. Lemma \ref{lem:monotonicity} then implies that  \eqref{Ainf} contradicts \eqref{bound}.

% \medskip

% The remaining three subcases can be dealt with in an entirely similar manner by replacing the role of $A_{\infty}$ above with the appropriate choice of one of $A_0$, $B_0$ and $B_\infty$.

\bigskip


\noindent{\bf Step 2.}  Next we show that  $\mathbf{F}$ cannot be of Type 2.1 when $d$ is large. Assume that for an arbitrarily large $d$ we can find a limit $\mathbf{F}$ of Type 2.1. This building has two essential curves 
$\mathbf{u} = \{u_2,\underline{u_d}\}$ and the punctures of $\underline{u}$ along $L_1$ are either all negative or positive.

\bigskip

\noindent{\it{Case 1.}} Suppose the punctures of $\underline{u}$ along $L_1$ are negative.  In this case we have:
\begin{align*}
    &M_1^\infty \geq M_1^0 + k_1,\\
    &k_1 +N_1 \geq d,\\
    &\widetilde{M}^0_2+\widetilde{M}^\infty_2   \geq d+2.
\end{align*}
% $$M_1^\infty \geq M_1^0 + k_1.$$ $$n_1 +N_1 \geq d,$$ and, by \eqref{constraints2.1}, we have 
% \begin{align}\label{local2.1}
%  \widetilde{M}^0_2+\widetilde{M}^\infty_2   \geq d+2.
% \end{align}
% Together with \eqref{area2}, these inequalities imply that 
% \begin{align*}
% \Omega (\mathbf{u}_d)+ 2rM_1^0 +rN_1 +2r(N_0+N_2) \leq b-2r.  
% \end{align*}
% Hence, $\Omega(u_2)$, $\Omega(\underline{u})$, $M^0_1$ and $N$ are all bounded from above independently of $d$.
By the last of these inequalities, either $\widetilde{M}^0_2 \ge (d+2)/2$ or $\widetilde{M}^\infty_2 \ge (d+2)/2$. Suppose that  \begin{equation}\label{m2}\widetilde{M}^\infty_2 \ge (d+2)/2. \end{equation} Then, intersecting $\mathbf{F}$ with $A_\infty$ we get
\begin{equation}\label{ainf}d+e =\mathbf{u} \cdot A_\infty + M_1^{\infty} + \widetilde{M}^\infty_2 +N.\end{equation}
By \eqref{m2} and the inequalities for $M^{\infty}_1$ and $N_1$ above, it follows that the right hand side of  \eqref{ainf} is at least $d + (d+1)/2$.  This produces a contradiction whenever  $d> 2e-2$. 

Similarly, if $\widetilde{M}^0_2 \ge (d+2)/2$, then intersecting $\mathbf{F}$ with $B_0$ we get
\begin{align*}
    d+e = \mathbf{u} \cdot B_0 + M_1^\infty + \widetilde{M}^0_2 +N. 
\end{align*}
The right hand side is at least $d + (d+2)/2$ and we again get a contradiction whenever  $d> 2e-2$. 



% and hence 
% \begin{align*}
%    \mathbf{u} \cdot A_0 \geq (d-1)/2 +e - (M_1^0 +N). 
% \end{align*}
% Since $M^0_1$ and $N$ are bounded from above independently of $d$, it follows from Lemma \ref{lem:monotonicity} that this will contradict the fact that $\Omega(u_2)$ and $\Omega(\underline{u})$ are also uniformly bounded when $d$ is sufficiently large.

% \medskip

\bigskip

\noindent{\it{Case 2.}} When the punctures of $\underline{u}$ along $L_1$ are positive, we have:
\begin{align*}
    &M_1^0 \geq M_1^\infty + k_1,\\
    &k_1 +N_1 \geq d,\\
    &\widetilde{M}^0_2+\widetilde{M}^\infty_2   \geq d+2.
 \end{align*}
In particular $M_1^0 + N \ge d$.

Intersecting with $A_0$ we get
\begin{align*}
    d+e = \mathbf{u} \cdot A_0 + M_1^0 + \widetilde{M}^\infty_2 +N,
\end{align*}
and together with the above this gives $\widetilde{M}^\infty_2 \le e$, so $\widetilde{M}^0_2 \ge d+2-e$.

Next we recall the area of our limiting building:
\begin{align*}
\Omega(\underline{u}) + r(M_1^0 + M_1^{\infty}+ \widetilde{M}^0_2+\widetilde{M}^\infty_2) + 2rN.
\end{align*}
Combining with the above inequalities this implies

\begin{align*}
\Omega(\underline{u}) + r(M_1^{\infty}+\widetilde{M}^\infty_2) 
&\leq 2dr + b - r(d - N) -r(d+2-e) - 2rN\\ &= b +r(e-2) - rN.
\end{align*}

This means that for fixed $e$ the count $N$, plus the total area of the essential curves together with all curves in our building lying in $\mathfrak{r}_{\infty}$ and $\mathfrak{s}_{\infty}$, is bounded independently of $d$. Thus all curves intersecting $B_{\infty}$ have uniformly bounded area, but as this intersection number is $d+e$ we contradict Lemma \ref{lem:monotonicity} when $d$ is large.




\bigskip


\noindent{\bf Step 3.}  The proof that $\mathbf{F}$ cannot be of Type 2.2 when $d$ is sufficiently large is nearly identical to the proof in {\bf Step 2}. The only change, after relabeling, is the change of the number of constraint points from $d$ to $d+1$.  This does not affect the argument.

\bigskip


\noindent{\bf Step 4.}  It remains to verify the various properties of the constituent curves of the Type 3 limit $\mathbf{F}$.\\

\noindent{\it Indices.} We begin by considering the indices of the curves of $\mathbf{F}$, starting at bottom level. Let $s_i$ be the total number of positive ends of the curves of $\mathbf{F}$ that map to $T^*L_i$. 


The total index of the curves mapping to $T^*L_i$ is denoted by $I_i$, and is at most $2s_i -2$. This upper bound corresponds to the case of a single curve with $s_i$ asymptotic ends. (As shown below, this is the case that actually occurs.) In any case the $I_i$ are necessarily even. To satisfy the point constraints we must have
\begin{equation}\label{consrtaint}
I_1 \ge 2d; \, \, I_2 \ge 2(d+1).
\end{equation}


Since $\mathbf{F}$ is of Type 3, two of the $s_i$ positive ends match with essential curves of $\mathbf{F}$. So, there must be $s_i -2$ additional curves of $\mathbf{F}$ corresponding to each of the remaining ends. Each of these curves form part of a planar sub-building in $\mathfrak{r}_0 \cup \mathfrak{r}_\infty$ when $i=1$ or a  planar sub-building in $\mathfrak{s}_0 \cup \mathfrak{s}_\infty$ when $i=2$. We denote the indices of these sub-buildings by $J_{k,i}$ for $i=1,2$. These numbers are odd, and equal $1$ exactly when the sub-building consists of a single plane.

The total index of the curves of $\mathbf{F}$ is given by \begin{equation}\label{eq:indg3}
   \mathrm{index}(\mathbf{u}) + I_1 + I_2 + \sum_{k=1}^{s_1-2} (J_{k,1}-1) + \sum_{k=1}^{s_2-2} (J_{k,2}-1) - 4 = 2(2d+1), 
\end{equation}
where the first term is the total index of the three essential curves, we subtract $1$ from the index of the planar sub-buildings for the matching conditions are their end, and the final term is the matching condition for the $4$ ends of the essential curves. We recall this sum is $2(2d+1)$ because the total index of a limiting building is equal to the index of holomorphic curves representing the class $(d,1)$.

Hence, using \eqref{consrtaint},
\begin{equation*}
    \mathrm{index}(\mathbf{u}) = \mathrm{index}(u_2)+\mathrm{index}(\underline{u})+\mathrm{index}(u_1) \leq 4,
\end{equation*}
with equality only if $I_1=2d$, $I_2=2(d+1)$ and all $J_{k,i}=1$.

Suppose $\mathrm{index}(\underline{u})=0$, so the essential cylinder is rigid. Then, in order to have matching ends in our building, neither the lower level curves in $T^* L_1$ nor $T^* L_2$ can be rigid, that is, both $I_1 > 2d$ and $I_2 > 2(d+1)$. But as these terms are even, we then have $$\mathrm{index}(\mathbf{u}) = \mathrm{index}(u_2)+\mathrm{index}(\underline{u})+\mathrm{index}(u_1)  \leq 0,$$ and this is a contradiction as holomorphic planes have odd index, and for regular almost complex structures the index is nonnegative.
Similarly we cannot have $\mathrm{index}(\underline{u}) \ge 4$ and both essential planes having odd index, so we conclude that
\begin{equation*}
    \mathrm{index}(u_2)=1, \, \, \mathrm{index}(\underline{u})=2, \, \, \mathrm{index}(u_1)=1.
\end{equation*}



%and the last term, $-(s_1 +s_2)$, comes from the matching constraints at Reeb orbits for the Morse-Bott Reeb flow of the flat metric on $\mathbb{T}^2$.


%Now, to accommodate the $d$ constraint points on $L_1$ we must have \begin{equation}\label{eq:s13}
    %s_1 \geq d+1,
%\end{equation} 
%and to accommodate the $d+1$ constraint points on $L_2$ we must have \begin{equation}\label{eq:s23}
%    s_2  \geq d+2.
%\end{equation}
%Since the total index of the curves of $\mathbf{F}$ is equal to $2(2d+1)$, it follows from \eqref{eq:indg3}, \eqref{eq:s13} and \eqref{eq:s23} that \begin{equation*}
%    \mathrm{index}(\mathbf{u}) = \mathrm{index}(u_2)+\mathrm{index}(\underline{u})+\mathrm{index}(u_1)\leq 4.
%\end{equation*}
%The genericity condition on $J$ then implies that $\mathrm{index}(u_2)= \mathrm{index}(u_1)=1$, and that $\mathrm{index}(\underline{u})=2$. Hence, the inequalities above must be equalities and there is exactly one curve of $\mathbf{F}$ in $T^*L_1$ and it has $d+1$ ends, and  there is exactly one curve of $\mathbf{F}$ in $T^*L_2$ and it has $d+2$ ends. The remaining curves of $\mathbf{F}$ must consist of $d-1$ planes in $\mathfrak{r}_0 \cup \mathfrak{r}_\infty$ and $d$ planes in $\mathfrak{s}_0 \cup \mathfrak{s}_\infty$.


\bigskip


\noindent{\it Areas.} Next we verify the symplectic areas of the essential curves of $\mathbf{F}$ asserted in Proposition \ref{prop:dd}. %From the discussion above and the fact that the total area of the curves of $\mathbf{F}$ is $2rd+b$, we have 
%\begin{equation*}
%\Omega(u_2) + \Omega(\underline{u}) + \Omega(u_1) = b+r.
%\end{equation*}
%It suffices to show that $\Omega(u_2)=s$ and that $\Omega(u_1)=b-s$.

%The essential curve $u_2$ is
\begin{lemma}\label{lem:s}
The $\Omega$-area of   $u_2$ is equal to $s$.
\end{lemma}

\noindent{\em The Proof of Lemma \ref{lem:s}.} We begin with a sequence of intermediate claims. Let $$v_2 \colon (D,\partial D) \to (X , L_2)$$ be the compactification of $u_2$, where  $D \subset \RR^2$ is the closed unit disk. Setting $\gamma_{2,d} = [v_2(\partial D)]  \in H_1(L_2;\ZZ)$, we have $\gamma_{2,d} \bullet \beta_2 =\pm 1.$

\begin{claim}\label{claim:u}
    There is a  smooth map $$u \colon (D,\partial D) \to (X \smallsetminus S_\infty, L_2)$$ such that 
    \begin{enumerate}
         \item $p \circ u (D)$  is equal to the closure of $A$ in $S_{\infty}$,
         \item $[u(\partial D)] =\gamma_{2,d} + \ell \beta_2 \in H_1(L_2; \ZZ)$ for some $\ell \in \ZZ$,
        \item The Maslov index of $u$ is $2$.
        
    \end{enumerate}
\end{claim}
\begin{proof} 
% WORK ON THIS. IT SHOULD JUST BE THE ESSENTIAL CURVE $u_2$
% Recall that $A$ is the component of $S_{\infty} \smallsetminus p(L)$ that contains the point $p(T_0)$. Consider a symplectic trivialization of the fibration $\bar{p}$ over the closure of $A$.  


Let $u'' \colon D \to  S_{\infty}$ be a smooth embedding such that $u''(D)$ is equal to the closure of $A$. 
Trivializing the bundle $p$ over $A$ we can then lift $u''$ to obtain a smooth map $u' \colon (D,\partial D) \to (X \smallsetminus S_\infty, L_2)$ that satisfies conditions (1) and (2). The Maslov index of $u'$ is even. Gluing a suitable multiple of a compactified disk from $\mathfrak{s}_0$ to $u'$ we get a disk $u$ of Maslov index 2 which still maps into $X \smallsetminus S_\infty$ and still satisfies conditions (1) and (2).

\end{proof}

\begin{claim}
A map $u$ as in Claim \ref{claim:u} has symplectic area equal to either $s$ or $2r-s$.
\end{claim} 

\begin{proof} View $L_2$ as a Lagrangian torus in the exact symplectic manifold $X \smallsetminus (S_\infty \cup T_{\infty})$.
By condition (1) of Lemma \ref{claim:u}, $u$ maps into $X \smallsetminus S_\infty \cup T_{\infty}$. Condition (2) implies that  $\{ \beta_2, [u(\partial D)]\}$ is an integral basis of $H_1(L_2; \ZZ)$. The area class and Maslov class of $L_2$ in $X \smallsetminus S_\infty \cup T_{\infty}$ are therefore determined by the areas and indices of the classes $\beta_2$ and  $[u(\partial D)]$. 
(These are well defined independently of a bounding disk, as we work now in $X \smallsetminus (S_\infty \cup T_{\infty}) \subset \RR^4$.)
The area of $\beta_2$ is $r$ and its Maslov index is $2$. By condition (3) of Lemma \ref{claim:u} the Maslov index of $[u(\partial D)]$ is also $2$. Then since $L_2$ is an image of $L(r,s)$ under a symplectomorphism of $\RR^4$, and $\{ \beta_2, [u(\partial D)]\}$ is a basis of $H_1(L_2; \ZZ)$, the area of $[u(\partial D)]$ must either be $s$ or $2r-s$ (depending on the orientation of our integral basis).
\end{proof}

\begin{claim}
The map $v_2$ has the same $\Omega$-area as the map $u$ in Claim \ref{claim:u}.
\end{claim} 

\begin{proof}
The glued map $v_2-u$ is a cycle in $X$ with boundary on $L_2$. Since its intersection number with both $T_0$ and $T_\infty$ is zero, it is homologous to a cycle that is contained in a (compactified) broken sphere of the foliation and has boundary on the equator. Since the Maslov index of the new cycle is zero, its symplectic area also vanishes. Hence, the area of $v_2-u$ is zero, as desired.
    
\end{proof}


We can now complete the proof of Lemma \ref{lem:s}. By the results above, it suffices to prove that the $\Omega$-area of $u_2$, and hence $v_2$, is not equal to $2r-s.$


\bigskip


\noindent{\it Case 1.} Suppose $u_2 \bullet S_\infty=0$.  Then, by positivity of intersection $u_2$ is a pseudo-holomorphic curve in $X \smallsetminus (S_\infty \cup T_\infty \cup L_2) \subset \RR^4$. It then follows from Corollary \ref{>r} that the symplectic area of $u_2$ must be at least $r$, which is greater than $2r-s$ by hypothesis.



% Let  $v \colon D \to X$  be the compactification of the essential curve $u_2$. Then $v$ will be a curve as in Lemma \ref{claim:u} and will have area equal to either $s$ or $2r-s$. Equivalently, the area of $u_2$ will be either $s$ or $2r-s$. Since $s>r$  we have $2r-s<r$ and it follows from Corollary \ref{>r} and the fact that $u_2$ is pseudo-holomorphic, that the area of $u_2$ must be $s$.

\bigskip


\noindent{\it Case 2.} Suppose that $u_2 \bullet S_\infty=m \ge 1$. Since $S_\infty$ has been inflated  by $2r-1$, the  symplectic area of $u_2$ is at least $m(2r-1)$, which, again by hypothesis, is strictly greater than $2r-s$.






% Let  $v \colon D \to X$  be the compactification of the essential curve $u_2$. Let $u$ be a curve as in Lemma \ref{claim:u}. ?????









% % Assume instead that each $u$ as in Lemma \ref{claim:u} has symplectic area equal to $2r-s$.  Since $v$ also satisfies conditions (1), (2)  and (3) of Lemma \ref{claim:u} it has area equal to $2r-s.$ WHY EXACTLY? 

% % By (1), $v$ misses $T_{\infty}$. Since $2r-s <r$, it follows from Corollary \ref{>r} that the intersection number between $v$ and $S_\infty$ bust be a positive number, say $m$. By our choice to inflate $S_\infty$ by $(2r-1)$



% \bigskip

% *******

% \bigskip


% The resulting Let $e = \omega(D)$. Given our area corrections after the inflation, we see that $e$ is independent of the choice of $D$.


% There exists a diffeomorphism $L(r,s) \to L$ which preserves both the area class and the Maslov class. Here, for Lagrangian tori $T \subset \CC^2$ we have $H_2(\CC^2, T) \cong H_1(T)$ and so we can consider area and Maslov classes as functions on $H_1(T)$. %Here we mean that the action of the diffeomorphism on $H_1$ preserves these two functions. 
% As a consequence, $(r,e) = (r,s)$ or $(r,e) = (r,2r-s)$.
% We need to exclude the second possibility and argue by contradiction.


% %We argue by contradiction and suppose $e < r$. As $c \ge r$, checking the possibilities as in section \ref{stretchL} we see that the only option is $(c,e) = (r, 2r-s)$. 

% Our stretching analyses can produce a holomorphic plane representing a Maslov $2$ class $D$, with $\partial D$ transverse to the foliation class and with $p(D)=A$. Hence $\omega(D)=e$. Let $m = D \bullet S_{\infty}$ and assume our almost complex structure is the standard product structure outside of $D(1) \times S^2(y) \subset S^2(2r) \times S^2(y)$. We can do this as $L \subset Z(1)$.

% As $r<s$ we have $e = 2r-s<r$ and so by  \cite{ho}, Lemma 3.7 the plane $D$ cannot lie in $\CC^2$, that is, $m \ge 1$.

% Then by monotonicity $D$ has area $$\mathrm{area}(D) = e > m(2r-1) \ge 2r-1$$ and as $s \ge 1$ this gives a contradiction as required.



% \section{Completion of Proof}
% \subsection{Homologically nontrivial case.}

% Suppose that one family of broken planes asymptotic to $L$ intersects $S_0$ and the other intersects $S_{\infty}$. (This property is always true for the planes asymptotic to $L_1$.) We consider the complement of the axes $\{z_1=0\}$, $\{z_2=0\}$, $\{ z_1 = \infty\}$, $\{ z_2 = \infty\}$ in  $S^2(2r) \times S^2 ( y)$ as a subset of $T^* T^2$. With this identification, $L$ and $L_1$ become homologically trivial Lagrangian tori, and our choice of $L_1$ implies they are relatively exact, that is, any surface in $T^* T^2$ with boundary on $L \cup L_1$ has symplectic area $0$.

% Up to translation we can set $L_1$ to be the zero section, but a theorem of Gromov says that exact Lagrangians in a cotangent bundle have nontrivial intersection with the zero section. This gives a contradiction.

% In the remainder of our proof then, we are free to assume that $L$ is homologically trivial. Nevertheless our scheme is to find an alternative set of axes with respect to which $L$ is homologically nontrivial, and this will lead to the same contradiction.


% \subsection{High degree curves.}\label{high}

% In this subsection we follow \cite{hiker} and consider the behavior of spheres in the class $\mathbb{A}_d := d[S^2 \times p] + [p \times S^2]$ under the stretching operation. For $d \ge 1$ we know that such curves exist, with images intersecting $2d+1$ generic points. The two cases we consider are first when $d$ points lie on $L$ and $d+1$ on $L_1$, and second when $d+1$ points lie on $L$ and $d$ on $L_1$. In the language of \cite{hiker}, we will see that under neck stretching limits of both families of curves have Type 3.

% Fix $d+1$ points on $L_1$ and $d$ points on $L$, where we take $d >>1$. Then, given an almost complex structure $J$, there exists a unique holomorphic sphere $u$ in the class $\mathbb{A}_d$ intersecting the $2d+1$ points. We take a sequence of $J_N$ stretching along $L \cup L_1$ and let $\mathbf{F}$ be the limiting holomorphic building of the corresponding $\underline{u}$.

% \begin{proposition}\label{limitF}(see \cite{hiker} Proposition 3.20)
% \begin{enumerate}
% \item  $\mathbf{F}$ has exactly three essential curves, one projecting to each of $A$, $B$ and $C$ and denoted $D^A_F, S^B_F$ and $D^C_F$ respectively;\\
% \item $D^A_F$ and $D^C_F$ are disks of Maslov class $2$;\\
% \item  $\mathbf{F}$ has exactly $d$ broken planes asymptotic to $L_1$, say $a_F$ intersect $\{z_1=0\}$ and $b_F$ intersect $S_{\infty}$;\\
% \item  $\mathbf{F}$ has exactly $d-1$ broken planes asymptotic to $L_1$, say $c_F$ avoid $S_{\infty}$ and $d_F$ intersect $S_{\infty}$.\\
% \end{enumerate} 
% \end{proposition}

% Counting areas, as $\mathbf{F}$ represents the homology class $\mathbb{A}_d$ we have $$\mathrm{area}(D^A_F) + \mathrm{area}(S^B_F) + \mathrm{area}(D^C_F) + (2d-1)r = 2rd + y$$
% where the term $(2d-1)r$ comes from the broken planes. By assumption $D^A_F$ has area $s$, and as it is asymptotic to the product torus we see that $D^C_F$ has area $y-s$. This implies that the cylinder $S^B_F$ has area $r$. 

% Similarly we can fix $d$ points on $L_1$ and $d+1$ points on $L$. Then, given an almost complex structure $J$, there exists a unique holomorphic sphere $v$ in the class $\mathbb{A}_d$ intersecting the $2d+1$ points. We take a sequence of $J_N$ stretching along $L \cup L_1$ and let $\mathbf{F}$ be the limiting holomorphic building of the corresponding $v_d$.

% \begin{proposition}\label{limitG}(see \cite{hiker} Proposition 3.20)
% \begin{enumerate}
% \item  $\mathbf{F}$ has exactly three essential curves, one projecting to each of $A$, $B$ and $C$;\\
% \item  $\mathbf{F}$ has exactly $d-1$ broken planes asymptotic to $L_1$, say $a_G$ intersect $\{z_1=0\}$ and $b_G$ intersect $S_{\infty}$;\\
% \item  $\mathbf{F}$ has exactly $d$ broken planes asymptotic to $L_1$, say $c_G$ avoid $S_{\infty}$ and $d_G$ intersect $S_{\infty}$.\\
% \end{enumerate} 
% \end{proposition}

% Next we carry out the perturbation procedure described in \cite{hiker}, Proposition 3.24. Here is the result.

% \begin{proposition}\label{spheres}
% There exist smooth symplectic spheres $F$, $G$ representing the class $\mathbb{A}_d$ and with the following properties.
% \begin{enumerate}
% \item  $F$ and $G$ intersect transversally and positively $d$ times in $p^{-1}(A)$ and $d$ times in $p^{-1}(C)$, and these are the only intersections;\\
% \item  $F$ intersects $S^B_F$ positively, transversally in a single point, while $G$ is disjoint from $S^B_F$.\\
% \end{enumerate} 
% \end{proposition}


% \subsection{Concluding arguments.}\label{end}

% Following \cite{hiker}, we can apply a blowing up and down procedure to produce a new copy of $S^2 \times S^2$ with axes given by the proper transforms of $F$ and $G$ in the class $[p \times S^2]$ and the transforms of the fibers $T_0 = \{z_2=0\}$ and $T_{\infty} = \{ z_2 = \infty\}$ representing the class $[S^2 \times p]$.

% The spheres $F$ and $G$ intersect positively, transversally in $2d$ points, $p_1, \dots, p_{2d}$, where we may assume $p_1, \dots, p_{d}$ lie in $p^{-1}(A)$ and $p_{d+1}, \dots, p_{2d}$ lie in $p^{-1}(C)$. Let $H_1, \dots, H_{2d}$ denote the fibers of $\mathcal{F}$ with $p_i \in H_i$. The $H_i$ are distinct since, by positivity of intersection, both $F$ and $G$ intersect each fiber in a single point.

% Now we blow-up balls of capacity $\varepsilon_1$ centered at each $p_i$. The proper transforms $\hat{F}$ and $\hat{G}$ are disjoint spheres in the blow-up, now of self intersection $0$. If we blow down the proper transforms $\hat{H}_i$ of the $H_i$ then we arrive at a new copy $Y$ of $S^2 \times S^2$, which includes the proper transforms $\hat{T}_0$ and $\hat{T}_{\infty}$ as spheres of self intersection $0$ intersecting each of $\hat{F}$ and $\hat{G}$ positively in a single point. We note that $Y$ is not monotone, indeed the spheres $\hat{T}_0$ and $\hat{T}_{\infty}$ have area $2r$ while the spheres $\hat{F}$ and $\hat{G}$ have area $2dr + y - 2d \varepsilon_1$.
% The Lagrangian tori $L_1$ and $L_2$ naturally transform to give tori in $Y$ disjoint from the axes. In particular, $L_1$ and $L_2$ can be viewed as disjoint Lagrangian tori  in the set $$Z := Y \smallsetminus ( \hat{F} \cup  \hat{G} \cup \hat{T}_0 \cup \hat{T}_{\infty} )$$ which can be identified with a subset of $T^* T^2$where the momentum coordinates are restricted to lie in a rectangle of sides $2r$ and $2dr + y - 2d \varepsilon_1$.

% \begin{lemma}
% As Lagrangian tori in $Z \subset T^*T^2$, $L_1$ and $L_2$ are relatively exact and homologically nontrivial.
% \end{lemma}

% \begin{proof}
% To prove this we make use of the curves in $Y$. First we show that $L_1$ and $L_2$ are relatively exact. For $L_1$ there is a disk $D_1$ with image in $Y$ that has area $r$, intersects  $\hat{F}$ in a single point, is disjoint from $\hat{G} \cup \hat{T}_0 \cup \hat{T}_{\infty}$ and has boundary on $L_1$ that  represents its foliation class, $\beta_1$. There is another such disk, $D_2$, for $L_2$. Let $\sigma_\beta \colon [0,1] \to \hat{F}$ be a path connecting the two intersection points, $D_i \cap \hat{F}$. Concatenating the disk maps with this path, as $-D_1 \# \sigma_\beta \# D_2$, we get a cylinder map, whose area is equal to zero and whose boundary represents the class $\beta_2 - \beta_1 \in H_1(L_1 \cup L_2)$. Perturbing this map away from $\hat{F}$ we then get a cylinder map  $C_{\beta}$ with the same properties, but whose image lies in $Z$. To complete the proof that $L_1$ and $L_2$ are relatively exact in $Z$ it suffices to find another cycle,  $C_{\gamma}$, which has zero area, represents a class in $H_2(Z, L_1 \cup L_2))$, and whose boundary is of the form $\gamma_2 -\gamma_1$  such that $\{\beta_i, \gamma_i\}$ is a basis of $H_1(L_i; \ZZ)$ for $i=1,2$. Let  $\hat{S}^B_F$ be the proper transform of the essential curve  $S^B_F$. It has the desired properties of $C_\gamma$ except that it has area $r$ and intersects $\hat{F}$ at a single point. These shortcomings can be corrected simultaneously by concatenating $\hat{S}^B_F$ with $-D_2$ and an appropriate path in $\hat{F}$ connecting the two intersection points. The resulting pair-of-pants map can then be perturbed off of $\hat{F}$ so that its image in $Z$. This is the required cycle, $C_{\gamma}$.

% It remains to show that $L_1$ and $L_2$ represent nontrivial classes in $H_2(Z;\ZZ)$. The proof is the same for each so we will consider only $L_1$. First we show that the boundary map \begin{align}\label{eq:boundary}
%     \partial \colon H_2(Z, L_1) \to H_1(L_1)
% \end{align}is zero. Suppose not. Then we can find a cycle $\alpha$ in $Z$ with boundary representing a nonzero class $$(m,n) \in H_1(L_1;\ZZ).$$
% The transform $\hat{D}^C_F$ of the essential curve $D^C_F$ has boundary that represents a class of the form $(-1,k) \in H_1(L_1;\ZZ).$
% Let $D_1$ be the foliation disk ion $L_1$ as above, whose boundary represents the class $(0,1)$. Using appropriate multiples of these two disks we can then complete 
% $\alpha$ to a closed cycle $\overline{\alpha}$ in $Y$. The important point is that the intersections of $\overline{\alpha}$ with $T_0$ and $T_1$ are determined by those of $\hat{D}^C_F$, and the intersections of $\overline{\alpha}$ with $\hat{F}$ and $\hat{G}$ are determined by those of $D_1$. If $m \neq 0$,  then, like  $\hat{D}^C_F$, the cycle $\overline{\alpha}$ will intersect $T_{\infty}$ and not $T_0$.
% Similarly, if $m=0$ then $n \neq 0$ and $\overline{\alpha}$ will intersect $\hat{F}$ and not $\hat{G}$. Neither of these intersection patterns is possible in $Y$ since $[T_0]=[T_\infty] \in H_2(Y;\ZZ)$ and 
% $[\hat{F}]=[\hat{G}] \in H_2(Y;\ZZ)$. Hence, the map \eqref{eq:boundary} is trivial, as claimed. From this it follows that the inclusion map $H_1(L_1;\ZZ) \to H_1(Z);\ZZ)$ is an isomorphism which, in turn, implies that $L_1$ is homologically nontrivial.

% \end{proof}


% % The foliation class in both $L_1$ and  $L_2$ bounds a disk of area $r$ which intersects  , but is disjoint from $\hat{G} \cup \hat{T}_0 \cup \hat{T}_{\infty}$. From this, we can find a cylinder $S^{\beta}$ with boundary components representing the foliation classes in $L_1$ and $L_2$ respectively, having area $0$, and disjoint from all axes $\hat{F} \cup \hat{G} \cup \hat{T}_0 \cup \hat{T}_{\infty}$.

% % The transform $\hat{D}^A_F$ of $D^A_F$ intersects $\hat{T}_0$ in a single point and is disjoint from  $\hat{T}_{\infty}$. Meanwhile the  transform $\hat{D}^C_F$ of $D^C_F$ intersects $\hat{T}_{\infty}$ in a single point and is disjoint from  $\hat{T}_0$. It follows that $L_1$ and $L_2$ are homologically nontrivial in $$Z := Y \smallsetminus ( \hat{F} \cup  \hat{G} \cup \hat{T}_0 \cup \hat{T}_{\infty} ).$$

% % Next the transform $\hat{S}^B_F$ is a cylinder with area $r$ and boundary components transverse to the foliation classes in $L_1$ and $L_2$ respectively. The cylinder is disjoint from $\hat{G} \cup \hat{T}_0 \cup \hat{T}_{\infty}$ but, by Proposition \ref{spheres}, intersects $\hat{F}$ in a single point. Adding a copy of $-D$, where $D$ is a broken plane asymptotic to $L_2$ intersecting $\hat{F}$, we find a cylinder (PAIR OF PANTS??) $L_B$ disjoint from all of the axes and of area $0$.

% Since $L_1$ is homologically essential in $T^*T^2$, it follows from Theorem B of  \cite{rgi} that $L_1$ is Hamiltonian isotopic to a constant section  of $T^*T^2$. Since $L_1$ and $L_2$ are relatively exact it follows from Gromov intersection theorem in \cite{gr}, Section $2.3.B''_4$, that $L_1$ and $L_2$ must intersect. With this we have arrived at a contradictions to Assumption \ref{assump}. 

% \subsection{Proof of Lemma \ref{dest}.}\label{prooflem}

% There exists a diffeomorphism $L(r,s) \to L$ which preserves both the area class and the Maslov class. Here, for Lagrangian tori $T \subset \CC^2$ we have $H_2(\CC^2, T) \cong H_1(T)$ and so we can consider area and Maslov classes as functions on $H_1(T)$. %Here we mean that the action of the diffeomorphism on $H_1$ preserves these two functions. 
% As a consequence, $(r,e) = (r,s)$ or $(r,e) = (r,2r-s)$.
% We need to exclude the second possibility and argue by contradiction.


% %We argue by contradiction and suppose $e < r$. As $c \ge r$, checking the possibilities as in section \ref{stretchL} we see that the only option is $(c,e) = (r, 2r-s)$. 

% Our stretching analyses can produce a holomorphic plane representing a Maslov $2$ class $D$, with $\partial D$ transverse to the foliation class and with $p(D)=A$. Hence $\omega(D)=e$. Let $m = D \bullet S_{\infty}$ and assume our almost complex structure is the standard product structure outside of $D(1) \times S^2(y) \subset S^2(2r) \times S^2(y)$. We can do this as $L \subset Z(1)$.

% As $r<s$ we have $e = 2r-s<r$ and so by  \cite{ho}, Lemma 3.7 the plane $D$ cannot lie in $\CC^2$, that is, $m \ge 1$.

% Then by monotonicity $D$ has area $$\mathrm{area}(D) = e > m(2r-1) \ge 2r-1$$ and as $s \ge 1$ this gives a contradiction as required.

\medskip

\noindent This completes the proof of Lemma \ref{lem:s}. 

\medskip

As $L_1$ is the standard product torus, and $u_1$ can be compactified to give a Maslov $2$ disk, we also have the following.
\begin{lemma}\label{lem:b-s}
The $\Omega$-area of $u_1$ is equal to $b-s$.
\end{lemma}

Finally, we add the areas of all curves in our limiting building. Writing $v_{k,1}$ and $v_{k,2}$ for the planar sub-buildings asymptotic to $L_1$ and $L_2$ respectively this gives
\begin{align*}
2dr + b &= \Omega(u_1) + \Omega(u_2) + \Omega(\underline{u}) + \sum_{k=1}^{s_1-2} \Omega(v_{k,1}-1) + \sum_{k=1}^{s_2-2} \Omega(v_{k,2}) \\
{}&\geq (b-s) + s + \Omega(\underline{u}) + r(s_1 -2) + r(s_2-2) \\
{}&\geq b + \Omega(\underline{u}) + r(d-1) + rd\\
{}&= \Omega(\underline{u}) + b + r(2d-1).
\end{align*}
Here, the first inequality uses Lemmas \ref{lem:s} and \ref{lem:b-s}, then just says our planar sub-buildings have area at least $r$. The second inequality applies \eqref{consrtaint}.

It follows that $\Omega(\underline{u}) \le r$, with equality only if $s_1 = d+1$ and $s_2 = d+2$. The essential curves $\underline{u}$ and $u_2$, the planes $v_{k,2}$ and lower level curves in $T^* L_2$ can be glued to give a disk with boundary on $L_1$, lying in $X \smallsetminus T_\infty$. Such a disk has even Maslov class, say $2n$, and then area $s + (n-1)r$. As $\Omega(u_2)=s$ and the $v_{k,2}$ have area a multiple of $r$, we see that $\underline{u}$ must also have area a multiple of $r$. As a holomorphic curve it has positive area, and so $\Omega(\underline{u}) = r$, and consequently $s_1 = d+1$ and $s_2 = d+2$.

With this, the proof of Proposition \ref{prop:dd} is complete. 

\bigskip

An entirely similar argument yields the following.

\begin{proposition}\label{prop:dd+1} For $d$ sufficiently large, every $(d,d+1)$-building is of Type 3. Each such building, $\mathbf{G}$, has three essential curves $u_2$, $\underline{u}$, and $u_1$ with areas $s$, $r$, and $b-s$, respectively. Moreover, the only top level curves of $\mathbf{G}$, other than its three essential curves, are $d$ broken planes in $\mathfrak{r}_0 \cup \mathfrak{r}_{\infty}$ and $d-1$ broken planes in $\mathfrak{s}_0 \cup \mathfrak{s}_{\infty}$. 
All the constituent curves of $\mathbf{G}$ are somewhere injective and hence regular. 
\end{proposition}

\subsection{Completion of the proof of Proposition \ref{prop:submanifolds}}
\label{sub:shift}

Let $D_A$, $O_B$ and $D_C$ be compactifications of the three essential curves of $\mathbf{F}$. Conditions (2)-(8) of Proposition \ref{prop:submanifolds} follow directly from this choice and Proposition \ref{prop:dd}. As established in Proposition 3.24 of \cite{hiker},  the desired embedded symplectic spheres $F$ and $G$ can be obtained by deforming the curves of the buildings $\mathbf{F}$ and $\mathbf{G}$ and then compactifying the new buildings to symplectic spheres. In particular, deformations can be prescribed to achieve the intersection properties (9) in Proposition \ref{prop:submanifolds}. The relevant relabeling needed to invoke Proposition 3.4 of \cite{hiker} directly, is as follows: $$\mathbb{L} \to L_2$$ $$L_{1,1} \to L_1$$ $$\mathbb{E} \to D_A$$ $$E_{1,1} \to D_C$$ $$A_0 \to A$$ $$A_\infty \to C.$$  
Proposition 3.24 in \cite{hiker} unfortunately does not state the intersections with $O_B$ as required for (9)(f), however this is implicit in the proof. Indeed, the sphere $G$ intersects the building $\mathbf{F}$ at least $d$ times in $p^{-1}(A)$ and $d$ times in $p^{-1}(C)$, and is disjoint from $L_1$ and $L_2$. Further we may assume $G$ is holomorphic near any intersections. As both $G$ and $\mathbf{F}$ represent the class $(d,1)$, and $(d,1) \cdot (d,1) = 2d$, this implies that there are actually precisely $2d$ intersections in $p^{-1}(A \cup C)$ and $G$ is disjoint from $O_B$. 

Regarding the deformation $F$ of $\mathbf{F}$, \cite[Lemma 3.33]{hiker} implies that $F$ intersects $\mathbf{F}$ precisely $d$ times in the closure of $p^{-1}(C)$ (we include the closure here to include intersections with planes in $\mathfrak{s}_0 \cup \mathfrak{s}_{\infty}$) and precisely $d-1$ times in $p^{-1}(A)$; in other words, each plane in $\mathfrak{r}_0 \cup \mathfrak{r}_{\infty} \cup \mathfrak{s}_0 \cup \mathfrak{s}_{\infty}$ leads to exactly one intersection. As $F$ is disjoint from $L_1 \cup L_2$ this leaves one additional intersection with $\mathbf{F}$, which occurs on $O_B$.

The remaining property, (1), is a simple consequence of the standard compactification procedure. $\Box$

%FOR COMPLETENESS, WE MIGHT NEED TO VERIFY THE INTERSECTIONS WITH $O_B$ SINCE THESE DO' NOT SEEM TO BE DEALT WITH IN \cite{hiker}.




\section{Proof of Theorem \ref{thm:hm}}

It suffices to identify a Lagrangian torus $\mathbf{L} \subset P(2r,b)$ that is Hamiltonian isotopic to $L(r,s)$, and then to construct a Hamiltonian diffeomorphism of $P(2r,b)$ that  displaces each of the $L_j$ from $\mathbf{L}$.

\medskip

\noindent{\bf Product Hamiltonian flows.} Arguing as in \cite{ho}, see also \cite{hicksmak}, we will construct the desired map using Hamiltonian diffeomorphisms of $\CC^2$ generated by functions of the form $$F(z_1,z_2) = G(z_1)H(z_2).$$ Denoting the time $t$ Hamiltonian flow of  $S \in C^{\infty}(M\times \RR /\ZZ)$ by $\phi^t_S$, we have  
\begin{align}\label{move}
    \phi^1_F(z_1,z_2) = \left(\phi^{H(z_2)}_G (z_1),\, \phi^{G(z_1)}_H (z_2) \right).
\end{align} We now recall a model setting in which maps of this type can be used to kill intersection points efficiently, in terms of the size of their support. 

Let $A$ and $B$ be subsets of $\CC^2$ and let $\mathrm{pr}_i$ be projection to the $z_i$-plane. Suppose that $\mathrm{pr}_1(A) = [-2a,\,a ] \times \{0\}$, for some $a>0$, and  $\mathrm{pr}_1(B) = \{0 \} \times [-1,\, 1]$. 
Suppose also that $$\mathrm{pr}_1(A \cap B)=(0,0).$$ 





Let $H(z_2)$ be a Hamiltonian, with $\min(H)=0$, whose time one flow disjoins $$\mathrm{pr}_2(\mathrm{pr}_1^{-1}(0,0) \cap A)$$ from $$\mathrm{pr}_2(\mathrm{pr}_1^{-1}(0,0) \cap B).$$
Define $G(z_1)$, near $\mathrm{pr}_1(A)$, to be a smoothing of the piecewise linear function  $$\chi(x_1) = \begin{cases} 0, &\text{for } x_1<-a+\varepsilon , \\
\frac{x+a-\varepsilon}{a-2\varepsilon}, &\text{for } -a+\varepsilon \leq x_1 \leq -\varepsilon , \\1, &\text{for } x_1>-\varepsilon,
\end{cases}$$ for some $0<\varepsilon<a.$ 
On $A$, the map $\phi^1_F$ is given (approximately) by the formula $$((x_1, y_1- H(x_2,y_2)\chi'(x_1)),\phi^{\chi(x_1)}_H(x_2,y_2)).$$
From this, and our choice of $H$, it follows easily that $\phi^1_F(A)$ is disjoint from $B$. Moreover, the projection $\mathrm{pr}_1(\Phi^1_F(A))$ is (approximately) equal to 
$$\mathrm{pr}_1(A) \cup ([-1, \, 0] \times [-\max(H)/a,\, 0]).$$ In particular, in the $z_1$-projection, the action of $\phi^1_F$ can be viewed as being local, and the rectangle $[-1, \, 0] \times [-\max(H)/a,\,0]$ can be viewed as the {\it cost} of disjoining $A$ from $B$. This is illustrated in Figure \ref{fig:model}. In more general situations, like those to come, this rectangle must be disjoint from $\mathrm{pr}_1(B)$ in order to avoid the creation of new intersections.

\begin{figure}[!h]
    \centering
    \includegraphics[width=1\linewidth]{Model-3.pdf}
    \caption{Displacing local models.}
    \label{fig:model}
\end{figure}

\medskip

With this model in hand, we begin the proof of Theorem \ref{thm:hm}. Before specifying $\mathbf{L}$ we reposition the $L_j = L(r) \times \Lambda_j$. We begin by repositioning the closed curves $\Lambda_j$ in $D(b)$. Equipping $D(b)$ with the rescaled polar coordinates $(\rho, \theta) \in [0,b) \times [0,1]$, fix an $a \in \left(\frac{s}{b},\frac{1}{k}\right)$ and move each $\Lambda_j$ so that is a smoothing of the boundary of $$\left[\epsilon, \frac{s}{a} +\epsilon\right] \times \left((j-1)(a+\delta), ja +(j-1)\delta\right)$$ where $\epsilon>0$ is arbitrarily small and $\delta$ is defined by
$ak + (k-1)\delta=1.$
See Figure \ref{fig:wedge}. 

\begin{figure}[!h]
    \centering
    \vspace{-0cm}
    \includegraphics[width=5.5in, height=7in]{Db.pdf}
     \vspace{-8cm}
    \caption{The repositioned circles $\Lambda_j$.}
    \label{fig:wedge}
\end{figure}

Identifying $D(2r)$with a square of side length $\sqrt{2r}$, we fix two smooth curves, $\gamma_1$ and $\gamma_2$, of area $r$ that are smoothings of the piecewise linear curves pictured in Figure \ref{fig:gamma}. The green rectangles appearing in Figure \ref{fig:gamma} all have area $\mathfrak{o}/k$ for the overlap $\mathfrak{o}$ defined as in \eqref{eq:g}. 
\begin{figure}[!h]
    \centering
    % \hspace{-3.91cm}
    \includegraphics[width=\linewidth]{Gamma12-4.pdf}
      \vspace{-3.5cm}
    \caption{Regions in the disk $D(2r)$.}
    \label{fig:gamma}
\end{figure}

The repositioned $L_j$ are then defined by setting  $L_j = \gamma_1 \times \Lambda_j$ if $j$ is odd and $L_j = \gamma_2 \times \Lambda_j$ if $j$ is even. 
The Lagrangian $\mathbf{L}$ is chosen to be of the form $\Gamma \times \Sigma$ where $\Gamma$ is a curve in $D(2r)$ and $\Sigma$ is a closed curve in $D(b)$. We choose $\Gamma$ to be a smoothing of the piecewise linear curve in $D(2r)$, bounding area $r$ pictured in Figure \ref{fig:gamma123}.
\begin{figure}[!h]
        \centering
    % \hspace{-3.91cm}
    \includegraphics[width=\linewidth]{Gamma12g.pdf}
      \vspace{-3.5cm}
    \caption{The circle $\Gamma$.}
    \label{fig:gamma123}
\end{figure}
The closed curve $\Sigma$ in $D(b)$ of area $s$ is chosen to that it intersects each $\Lambda_j$ as pictured in Figure \ref{fig:Sigma1}. Here, again, the green rectangles are assumed to have area equal to $\mathfrak{o}/k$.
\begin{figure}[!h]
    \centering 
    % \hspace{-3.91cm}
    \includegraphics[width=\linewidth]{Sigma1-2.pdf}
      \vspace{-8cm}
    \caption{The circle $\Sigma$.}
    \label{fig:Sigma1}
\end{figure}


\begin{remark}The configuration of curves of area $r$ and boxes of area $\mathfrak{o}/k$ pictures in Figure \ref{fig:gamma123} is only possible under the assumption that $r>6\mathfrak{o}/k$.\end{remark}




We now displace the repositioned $L_j$ from $\mathbf{L}$ via a Hamiltonian diffeomorphism of $P(2r,b)$. This will be done in stages, for consecutive pairs of the $L_j$. Consider first $L_1$ and $L_2$. Let $H_1$ be a Hamiltonian on $D(b)$ that displaces $\Lambda_1$ from $\Sigma$ such that $\phi^1_{H_1}(\Lambda_1)$ intersects only $\Lambda_1$ and $\Lambda_2$ as pictured in Figure \ref{fig:H1}.  
\begin{figure}[!h]
    \centering 
    % \hspace{-3.91cm}
    \includegraphics[width=\linewidth]{H1.pdf}
      \vspace{-8cm}
    \caption{Moving $\Lambda_1$ away from $\Sigma$.}
    \label{fig:H1}
\end{figure}
Similarly, let $H_2$ be a Hamiltonian that displaces $\Lambda_2$ from $\Sigma$ such that $\phi^1_{H_2}(\Lambda_2)$ intersects only $\Lambda_1$ and $\Lambda_2$ as depicted in Figure \ref{fig:H2}. Both $H_1$ and $H_2$ can be chosen so that their minimum value is zero and their maximum value, $\bar{H}$, is the same. This shared maximum value $\bar{H}$ is greater than $\mathfrak{o}/k$ and can be chosen to be arbitrarily close to $\mathfrak{o}/k$.  Note also that both $H_1$ and $H_2$ can be chosen to have support contained in a region of the form $\{\eta < \theta < 2a + \delta-\eta\}$ for some $\eta>0$.

\begin{figure}[!h]
    \centering 
    % \hspace{-3.91cm}
    \includegraphics[width=\linewidth]{H2.pdf}
      \vspace{-8cm}
    \caption{Moving $\Lambda_2$ away from $\Sigma$.}
    \label{fig:H2}
\end{figure}

Let $G$ be a function on $D(2r)$ such that in the region $g$ pictured in Figure \ref{fig:g}, we have $G(z_1) = \chi(x_1)$ where $\chi$ is a smoothing of a piecewise linear function, similar to the one in the  model, whose constant values along a horizontal strip in $g$ are labeled in Figure \ref{fig:g}. Assume also that $G=0$ outside an arbitrarily small neighborhood of the region $g$. 

\begin{figure}[!h]
    \centering 
    % \hspace{-3.91cm}
    \includegraphics[width=\linewidth]{G.pdf}
      \vspace{-3.5cm}
    \caption{The function $G$.}
    \label{fig:g}
\end{figure}



\begin{figure}[!h]
    \centering 
    % \hspace{-3.91cm}
    \includegraphics[width=\linewidth]{Pr1.pdf}
      \vspace{-3.5cm}
    \caption{Projections of the tori $\tilde{L}_i$.}
    \label{fig:proj}
\end{figure}


Set $F_1(z_1,z_2) =G(z_1)H_1(z_2)$ and $F_2(z_1,z_2) =G(z_1)H_2(z_2)$ and consider the Lagrangian tori $\tilde{L}_1 = \phi^1_{F_1}(L_1)$ and $\tilde{L}_2 = \phi^1_{F_2}(L_2)$. The images of $\tilde{L}_1$ and $\tilde{L}_2$ under $\mathrm{pr}_1$ are pictured in Figure \ref{fig:proj}. Because the additional areas can be arranged to lie in our green regions, these projections intersect in the same two points, $\alpha$ and $\beta$, as do $\gamma_1 =\mathrm{pr}_1(L_1)$ and $\gamma_2=\mathrm{pr}_1(L_2)$. 
Similarly, the images of $\tilde{L}_1$ and $\tilde{L}_2$ under $\mathrm{pr}_1$ also intersect $\Gamma$ in the same points as $\gamma_1 $ and $\gamma_2$, respectively.

\begin{lemma}\label{lem:off}
    For the construction above
    \begin{enumerate}
        \item $\tilde{L}_1$, $\tilde{L}_2$  and $\mathbf{L}$ are disjoint from one another.\\
        \item $\mathrm{pr}_2(\tilde{L}_1)$ and $\mathrm{pr}_2(\tilde{L}_2)$ are contained in a subset of $D(b)$ of the form $\{\eta \leq \theta \leq 2a + \delta-\eta\}$ for some $\eta>0$.
    \end{enumerate}
\end{lemma}

\begin{proof}

It follows from the discussion above that any point of intersection of $\tilde{L}_1$ and $\tilde{L}_2$ must lie in either $(\mathrm{pr}_1)^{-1}(\alpha)$ or $(\mathrm{pr}_1)^{-1}(\beta)$. But the intersections of $\tilde{L}_1$ and $\tilde{L}_2$ with these two fibers are the same as those of $L_1$ and $L_2$ since $G(\alpha)=G(\beta) =0.$ In particular,  $\mathrm{pr}_1^{-1}(\alpha) \cup \tilde{L}_i =\Lambda_i=\mathrm{pr}_1^{-1}(\beta) \cup \tilde{L}_i$ for $i =1,2$. Hence, $\tilde{L}_1$ and $\tilde{L}_2$ are disjoint.

Similarly, any point of intersection of $\tilde{L}_1$ and $\mathbf{L}$ must lie in the fiber of $\mathrm{pr}_1$ over one of the two intersection points of $\mathrm{pr}_1(\tilde{L}_1)$ and $\mathrm{pr}_1(\mathbf{L}) =\Gamma.$ Labeling these two points $\alpha_1$ and $\beta_1$, the definition of $G$ implies that it is identically equal to one near both. Hence, 
$$\mathrm{pr}_1^{-1}(\alpha_1) \cup \tilde{L}_1 = \mathrm{pr}_1^{-1}(\beta_1) \cup \tilde{L}_1 =\phi^1_{H_1}(\Lambda_1)$$ while 
$$\mathrm{pr}_1^{-1}(\alpha_1) \cup \mathbf{L} = \mathrm{pr}_1^{-1}(\beta_1) \cup \mathbf{L} =\Sigma.$$ By the definition of $H_1$, these are disjoint and hence so are $\tilde{L}_1$ and  $\mathbf{L}$. A similar argument yields $\tilde{L}_2 \cap\mathbf{L} = \emptyset.$

The last assertion follows immediately from the fact, mentioned above, that both $H_1$ and $H_2$ can be chosen to have support contained in a region of the desired form.
\end{proof}

Invoking the last assertion of Lemma \ref{lem:off}, and defining $H_i$ analogously to $H_1$ or $H_2$ depending on the parity of $i$, the remaining $\tilde{L}_j$ can be repositioned within $P(2r,b)$, in pairs,  so that in the end they are disjoint from each other and from $\mathbf{L}$, as desired.



% We show here that this rigidity phenomenon vanishes completely if one allows for Hamiltonian diffeomorphisms of $\RR^4$ that are not supported in $P(2r,b)$.


% \begin{proposition}\label{flex} For any $r>0$, there exists a Hamiltonian diffeomorphism $\phi$ of $\RR^4$ that maps  $L(r,s)$ into $P(2r, b)$ such that $\phi(L(r,s))$ is disjoint from the collection $\{L_j\}_{j=1, \dots k}.$ When $2r<b-s$, there exists a Hamiltonian diffeomorphism $\phi$ of $\RR^4$ such that $\phi(\overline{P(r,s)}) \subset P(2r, b)$ and $\phi(\overline{P(r,s)})$ is disjoint from the collection $\{L_j\}_{j=1, \dots k}.$
% \end{proposition}

% THE CURRENT CONSTRUCTION HAS AN ERROR.

% \begin{proof}


% First we note that since $b>ks$ and $k \ge 2$, it follows from  Theorem \ref{thm:hm} that it suffices to consider the case  \begin{equation}\label{eq:cut} 2r<b-s.\end{equation} 
% In this case we will construct a Hamiltonian diffeomorphism $\Psi$ of $\RR^4$ and a Hamiltonian $\Phi$ of (supported in) $P(2r,b)$
% such that $\Psi(\overline{P(r,s)})$ is contained in $P(2r,b)$ and is disjoint from all of the $\Phi(L_j)$.



% % All collections of $k$ circles in $D(y)$ are symplectomorphic, therefore it suffices to find a single family of $\Lambda_n$ such that the corresponding $L_n$ are disjoint from the image of $P(r,s)$. Also, we can replace $L_n$ by $E_n \times \Lambda_n$, where the $E_n$ are circles in $D(2r)$ bounding area $r$, and the Hamiltonian isotopy type of the union of the $L_n$ remains the same.

% \begin{remark}
%     In constructing the Hamiltonian diffeomorphism $\Psi$ of $\RR^4$, our primary concern will be the image of $\overline{P(r,s)}$. We will allow ourselves to consider the effects, on $\overline{P(r,s)}$, of Hamiltonian flows that are complete but do not have compact support. These flows can be cut-off, away from the path of $\overline{P(r,s)}$, in the standard manner. This cut-off procedure will not be mentioned. 
    
%     We will also invoke the Extension after Restriction principle and, at times, speak only of symplectic embeddings. Since $\overline{P(r,s)}$ is star-shaped, any symplectic embedding of $\overline{P(r,s)}$ into $\RR^4$, that extends to a neighborhood of $\overline{P(r,s)}$, can be realized by a Hamiltonian diffeomorphism, \cite{schl-er}.
% \end{remark}

% \subsection{$\Psi$} The map $\Psi$ will be defined as a composition of four  Hamiltonian diffeomorphisms, $$\Psi = \Psi_4 \circ \Psi_3 \circ \Psi_2 \circ \Psi_1.$$ These will be constructed in  sequence.

% \noindent{\it The map $\Psi_1$.}  This first map will act on the factors of $$\overline{P(r,s)} = \overline{D(r)} \times \overline{D(s)},$$ separately. We will view $\RR^4$ as $\RR^2 \times \RR^2$ and denote the coordinates on the two factors by $w=(u,v)$ and $z=(x,y)$, respectively. 



% For any $\varepsilon_1>0$, there is a Hamiltonian diffeomorphism $F_1$ of $\RR^2$ such that $F_1(\overline{D(r)})$ is contained in the rectangle
% $$D_1 = [0,\,1] \times [0, r(1+\varepsilon_1)].$$ This will be the first factor of $\Psi_1$.

  
% Let $m$ be the unique positive integer such that  \begin{equation}\label{eq:em}
%  mr < s \leq (m+1)r.   
% \end{equation} For an arbitrarily small $\varepsilon_2>0$ choose a number $c$ in $(0, (1-\varepsilon_2)r)$  and define $\ell$ by \begin{equation}\label{eq:ell}
%     (2\ell +m)c = s.
% \end{equation}
%  Let $T_{-1} = [-1-\ell, -1] \times [0,c]$, $T_m = [2m, 2m+\ell] \times [0,c]$ and $T_i = [2i, 2i+1] \times [0, c]$ for $0 \le i \le m-1$. Fix an $\varepsilon_3>0$ and  set 
% $$Z =\left( \bigcup_{i=-1}^{m} T_i\right) \cup \left([-1 - \ell, 2m + \ell] \times [0, \varepsilon_3]\right).$$
% The set $Z$, pictured in Figure 1, is connected and its area is $$s+(m+1)\varepsilon_3.$$ Hence, there is a Hamiltonian diffeomorphism $F_2$ of $\RR^2$
% such that $F_2(\overline{D(s)})$ is contained in $Z$. 
% \begin{figure}
%     \centering
%     \vspace{-0cm}
%     \includegraphics[width=5in, height=7in]{Figure Z0-2.pdf}
%      \vspace{-12cm}
%     \caption{$Z$}
%     \label{fig:z0}
% \end{figure}
% \medskip

% Defining the Hamiltonian diffeomorphism $\Psi_1 \colon  \RR^4 \to \RR^4 $ by $$\Psi_1(u,v,x,y) = (F_1(u,v), F_2(x,y)),$$  we have $\Psi_1(\overline{P(r,s)}) \subset D_1 \times Z.$


% \medskip

% \noindent{\it The map $\Psi_2$.} The second  map $\Psi_2$ is constructed  by coupling simple Hamiltonian flows in the plane. The first of these is the flow of the Hamiltonian
% function $g(u,v) = (1+ \varepsilon_2)v$ given by, $$\phi^t_{g}(u,v) = (u+ t(1+\varepsilon_2),v).$$
% Note that for $j \in \NN$, we have 
% $$\phi^j_g(D_1) = [(1 + \varepsilon_2)(j-1) , (1 + \varepsilon_2)j - \varepsilon_2] \times [0, r(1+\varepsilon_1)].$$
% Setting $D_j =\phi^j_g(D_1)$, we note that the rectangles $D_j$ are pairwise disjoint.
 

% The second type of simple planar Hamiltonian flow we consider is determined by a bump function. Fix a constant  $\varepsilon_4>0$ and let $\chi : \RR \to [0,1]$ be a smooth function such that $\chi(x) =0$ for $x \le \varepsilon_4/4$, $\chi(x) =1$ for $x \ge 1- \varepsilon_4/4$,  and $0 \le \chi'(x) \le 1+\varepsilon_4$. For each integer $j$ the Hamiltonian flow of the 
% function $$f_j(x,y)= -\chi(x-2j+3)$$ is given by  $$\phi^t_{f_j}(x,y) = (x,y + t\chi'(x-2j+3)).$$


% % Given an $\varepsilon_3>0$ arbitrarily small we fix a smooth function $\chi : \RR \to [0,1]$ such that $\chi(x) =0$ for $x \le \frac{\varepsilon_3}{2}$, $\chi(x) =1$ for $x \ge 1- \frac{\varepsilon_3}{2}$ and $0 \le \chi'(x) \le 1+2\varepsilon_3$. 
% % \begin{figure}
% %     \centering
% %     \vspace{-4cm}
% %     \includegraphics[width=4in, height=4in]{Figure chi-2.pdf}
% %     \vspace{-3cm}
% %     \caption{The function $\chi$}
% %     \label{fig:chi}
% % \end{figure}
% % The Hamiltonian flow of the 
% % function $f_j(x,y)= \chi(x-2j+3)$ is $$\phi^t_{f_j}(x,y) = (x,y + t\chi'(x-2j+3)).$$
 
% Now consider the product Hamiltonian functions on $\RR^4$ defined by $$H_j(u,v,x,y)= g(v) f_j(x).$$ The flow of $H_j$ is
% $$\phi^t_{H_j}(u,v,x,y) = \left(\phi^{tf_j(x)}_g(u,v), \phi^{tg(v)}_{f_j}(x,y)\right)$$ and it's time one flow is 
% \begin{align}\label{flowj}
%      \phi^1_{H_j}(u,v,x,y) &= (u+f_j(x)(1+\varepsilon_2), v ,x , y+(1+\varepsilon_2)v f_j'(x) ).
% \end{align}
% % In other words
% % \begin{align*}
% %      \phi^1_{H_j}(u,v,x,y) &= \begin{cases} (u, v ,x , y),&  \text{if } x \le 2j-1\\
% % (u+1+\varepsilon_2, v ,x , y ),&  \text{if }2j\le x\\
% % (u+f_j(x)(1+\varepsilon_2), v ,x , y+(1+\varepsilon_2)v f_j'(x) ),&  \text{otherwise}.
% %      \end{cases}
% % \end{align*}




% The second map defining  $\Psi$ is $$\Psi_2 = \phi^1_{H_{m+1}} \circ \phi^1_{H_{m}} \circ \cdots \circ \phi^1_{H_{1}}.$$


% Set $P_0 = D_1 \times Z$. Since $\Psi_1$ embeds $\overline{P(r,s)}$ into $P_0$, in order to understand $\Psi_2 \circ \Psi_1(\overline{P(r,s)})$ it suffices to understand the sets $P_{j} = \phi^1_{H_{j}} (P_{j-1})$ for $j=1$ to $m+1$. 

% %In particular, we have $$\Psi_2 \circ \Psi_1(\overline{P(r,s)}) \subset P_{m+1}.$$ 

% To analyze  $P_j$ it will be useful to consider its images under the projections $\pi_w \colon \RR^4 \to \RR^2$ to the $w$-plane, and $\pi_z \colon \RR^4 \to \RR^2$ to the $z$-plane.
% For $P_1$, equation \eqref{flowj} implies that  $$\pi_w(P_1) = [0, 2+\varepsilon_2] \times [0, r(1+\varepsilon_1)].$$ It also implies that  $\pi_z(P_1)$ is contained in $Z \cup S_0$ where 
% \begin{align*}
% S_0=\left[-1 + \varepsilon_4/4,-\varepsilon_4/4 \right] \times [0,(1+\varepsilon_1)(1+\varepsilon_2)(1 + \varepsilon_4)r +\varepsilon_3].
% \end{align*}
% Finally, it follows from equation \eqref{flowj} that for each  $(x,y) \in  \pi_z(P_1)$, we have 
% \begin{align*}
% \pi_z^{-1}(x,y)=
% \begin{cases}
%  D_1 \times \{(x,y)\},& \,\, \text{if } (x,y) \in x \leq -1 + \varepsilon_4/4\\
%  \phi_g^{ \chi(x +1)} (D_{1}) \times \{(x,y)\},& \,\, \text{if } (x,y) \in S_0\\
%  D_2 \times \{(x,y)\},& \,\, \text{otherwise}.
% \end{cases}
% \end{align*}


% Proceeding in this way, we see that the $w$-projections of the $P_j$ become larger and larger rectangles, while we need to add more rectangles to $Z$ to contain the $z$-projections of the $P_j$. In the end, we obtain the following description of $P_{m+1}$ in terms of its two projections. 

% The $w$-projection, $\pi_w(P_{m+1})$, is equal to the rectangle $$W=[0, m+2 + (m+1) \varepsilon_2] \times [0, r(1+\varepsilon_1)].$$ As pictured in Figure \ref{fig:W}, $W$ contains the disjoint rectangles $D_1,\dots, D_{m+2}$. 
% \begin{figure}
%     \centering
%     \includegraphics[width=4in, height=6in]{Figure W2.pdf}
%     \vspace{-9cm}
%     \caption{$\pi_w(P_{m+1}) =W$}
%     \label{fig:W}
% \end{figure}
% The area of $W$ is arbitrarily close to $(m+2)r$ which, by \eqref{eq:cut} and \eqref{eq:em}, is less than $b$. Hence we may assume that the area of $W$ is less than $b$.

% The $z$ projection, $\pi_z(P_{m+1})$, is contained in the set $$V= Z \cup \bigcup_{i=0}^m S_i.$$ Here, $S_i$ is the rectangle obtained by shifting $S_0$ to the right by $2i.$ 
% Let $\bar{T}_j$ denote the component of $V \smallsetminus \left(\bigcup_{i=0}^m S_i \right)$ that contains $T_j$, see Figure \ref{fig:z0si}.
% \begin{figure}
%     \centering
%     \vspace{-0cm}
%     \includegraphics[width=5in, height=7in]{Figure Vgap-5.pdf}
%     \vspace{-12cm}
%     \caption{The projection $\pi_z(P_{m+1})$ is contained in the set $V = Z \cup \left( \bigcup_{i=0}^m S_i\right)$.}
%     \label{fig:z0si}
% \end{figure}

% For all $(x,y) \in  \pi_z(P_{m+1})$ we have 
% \begin{align}\label{eq:shift}
% \pi_z^{-1}(x,y)=
% \begin{cases}
%  D_{j+2} \times \{(x,y)\},& \,\, \text{if } (x,y)\in \bar{T}_j\\
%  \phi_g^{ \chi(x -2i +1)} (D_{i}) \times \{(x,y)\},& \,\, \text{if } (x,y) \in S_i.
% \end{cases}
% \end{align}




% \medskip

% \noindent{\it The map $\Psi_3$.}  The Hamiltonian diffeomorphism $\Psi_3$ will embed $P_{m+1}$ into $D(2r) \times W$.  To define it, we first construct a symplectic immersion $\psi_{3}$ of $\pi_z(P_{m+1})$ into $D(2r)$. It suffices to define $\psi_3$ on the region $V$. 

% The area of each rectangle $S_i \subset V$ is $$\mathcal{S} = (1-\varepsilon_4/2)((1+\varepsilon_2)(1+\varepsilon_2)(1 + \varepsilon_4)r +\varepsilon_3).$$ The area of each $\bar{T}_j \subset V$ for $j=0, \dots m-1$ is $\mathcal{T}= c + \varepsilon_3 \varepsilon_4/2$,  and the areas of $\bar{T}_{-1}$ and $\bar{T}_{m+1}$ are less than $c$. For a fixed $c<r$, coming from the definition of $Z$ above,   we can choose $\varepsilon_1$, $\varepsilon_2$, $\varepsilon_3$ and $\varepsilon_4$ small enough so that  $\mathcal{S}<2r-c$ and $\mathcal{S} +\mathcal{T}<2r$. This will allow us to construct the immersion $\psi_3$ so that the images of each $S_i$ and $\bar{T}_j$ are distinct. 

% %In fact, $\psi_3$ will be defined so that the $S_i$ have identical images and the $T_j$, for $j=1, \dots, m-1$ have identical images. 

% Consider polar coordinates $R = \pi |z|^2$ and $\theta \in \RR / \ZZ$ on the $z$-plane. On $S_i$, the immersion $\psi_3$ will be given by the explicit formula
% \begin{align}\label{eq:psiS}
%     \psi_3 (x,y) &=\left(R=\frac{2y}{\mathcal{E}(1+\varepsilon_5)} +\varepsilon_6,\, \theta =\frac{ \mathcal{E}(1 +\varepsilon_5)}{2}(x+1-2i)\right),
% \end{align}
% where $\varepsilon_5$ and $\varepsilon_6$ are again small positive constants. We assume here that $\varepsilon_5$ is larege enough with respect to  $\varepsilon_1$, $\varepsilon_2$, $\varepsilon_3$ and $\varepsilon_1$
% so that $R< 2r$ for all $(x,y) \in S_i$. We also assume that $\varepsilon_6$ is small enough for this to remain true. As illustrated in Figure \ref{fig:si}, the image $\psi_3(S_i)$ is a region of the form
% $$\{R_0 \le R \le R_1, \theta_0 \le \theta \le \theta_1\},$$ where $0<R_0 < R_1 <2r$ and $0< \theta_0 <\frac 12 < \theta_1 <1$.  
%  \begin{figure}
%     \centering
%     \includegraphics[width=4in, height=5.5in]{Wgreen-4.pdf}
%     \vspace{-4.5cm}
%     \caption{The images $\psi_3(S_i) \subset D(2r)$ are equal and  pictured in green.}
%     \label{fig:si}
% \end{figure}


% The immersion $\psi_3$ is chosen so that it restricts to each $\bar{T}_j$ as an embedding with image in the complimentary region $$\{0 \le R <2r, \theta_1 \le\theta \le 1\} \cup \{0 \le R <2r, 0 \le \theta \le \theta_0\}.$$  This is possible since the area of this region is $$r\left( 1- \varepsilon_5 + \varepsilon_4/2 + \varepsilon_4 \varepsilon_5/4\right)$$ which is greater than $c$ when the $\varepsilon_j$ are all sufficiently small. For $j=1, \dots, m-1$, the images $\psi_3(\bar{T}_j)$ will be chosen to be identical, as pictured in Figure \ref{fig:tj}.
% \begin{figure}
%     \centering
%     \includegraphics[width=4in, height=5.5in]{Tred-2.pdf}
%     \vspace{-4.5cm}
%     \caption{The images $\psi_3(\bar{T}_j) \subset D(2r)$, for $j=0, \dots m-1$, are equal and pictured in red.}
%     \label{fig:tj}
% \end{figure}
% The images of $T_{-1}$ and $T_{m}$ are assumed to be nearly identical. Overall, the immersion $\psi_3$ can be viewed as a clockwise wrapping of $V$ around the center of $D(2r)$. The first three stages and final three stages of this wrapping are pictured in Figure \ref{fig:wrap}.

% % This is possible since $c < 2r - (1+ \varepsilon_2)r = (1 - \varepsilon_2)r$.
% % The immersion $\psi$ is illustrated in Figure 3. Note we are assuming $\psi(T_{-1})$ and $\psi(T_m)$ have roughly the same image, which has area $lc$.


% % Using polar coordinates $(\rho, \theta)$ in the $uv$-plane we can then find $\rho_0< \sqrt{\frac{2r}{\pi}}$ and $\theta_0 >\pi$ such that the area of the region $$U_0=\{\rho < \rho_0, 0\le\theta\le \theta_0\}$$ has area greater than $(1+\varepsilon_2)(1+\varepsilon_1)(1 + \varepsilon_3)r$ and the region $$V_0=\{\rho <\rho_0, \theta_0\le \theta_0 \le 2\pi\}$$ has area greater than $c$.

% % There exists a symplectic embedding $\alpha_1$ of $T_{-1}$ into $V_0$ such that the image of the boundary segment $\{u=-1, v\in [0,c]\}$ contains a segment on $\theta=\theta_0$. This can be extended to a symplectic embedding $\alpha_2$ of $T_{-1} \cup S_1$  into $D(2r)$ such that $\alpha_2(T_{-1})= \alpha_1(T_{-1})$ and $\alpha_2(S_{1})$ is contained $U_0$. In turn, $\alpha_2$ can be extended to a symplectic immersion $\alpha_3$ of $T_{-1} \cup S_1 \cup T_1$  into $D(2r)$ such that $\alpha_3(T_1)$ is mapped into $V_0$ and the image of the boundary segment $\{u=1, v\in [0,c]\}$ contains a segment on $\theta=\theta_0$. Continuing in this manner we get our desired  symplectic immersion $\psi_3 =\alpha_{2m+1}$, where, in the end, $\psi_3$ maps $T_m$ into $V_0$, see Figure \ref{fig:stages}. We note that $\psi_3$ acts identically on each of the $S_j$. In particular, if $(x_1,y_2) \in S_{j_1}$ and $(x_2,y_2) \in S_{j_2}$ have the same image under $\psi_3$, then $x_1 -2j_1 = x_2 -2j_2$ and $y_1=y_2$.


% % To do this, when $(x,y) \in S_i$ we set $$\psi_{3}(x,y) = \left(R = \frac{2y}{(1+\varepsilon_2)(1+4\varepsilon_3)}, \theta = \frac{(1+\varepsilon_2)(1+4\varepsilon_3) (x+3 - 2i)}{2} \right)$$
% % and map the remainder of $\pi(P)$ to $D(2r) \smallsetminus \{ \theta \in [0, (1 + \varepsilon_2)/2] \}$ such that we have a symplectic embedding on each square $T_i$. Since $c< (1-\varepsilon_2)r$, this is possible for all sufficiently small $\varepsilon_1>0$ and $\varepsilon_3>0$.
% % The immersion $\psi_3$ is illustrated in Figure 3. Note we are assuming $\psi(T_{-1})$ and $\psi(T_m)$ have roughly the same image, which has area $lc$.

% % 
% % 
% \begin{figure}
%     \begin{subfigure}[t]{0.3\linewidth}
%     \includegraphics[width=\linewidth, height=\textheight, keepaspectratio]{FigureT-1.pdf}
%     \vspace{-2cm}
%     \caption{$\psi_3(T_{-1})$}
%     \end{subfigure}
%     \begin{subfigure}[t]{0.3\linewidth}
%     \includegraphics[width=\linewidth, height=\textheight, keepaspectratio]{Wgreen.pdf}
%     \vspace{-2cm}
%     \caption{$\psi_3(S_1)$}
%     \end{subfigure}
%     \begin{subfigure}[t]{0.3\linewidth}
%     \includegraphics[width=\linewidth, height=\textheight, keepaspectratio]{Tred-2.pdf}
%     \vspace{-2cm}
%     \caption{$\psi_3(T_{0})$}
%     \end{subfigure}
%     \begin{subfigure}[t]{0.3\linewidth}
%     \includegraphics[width=\linewidth, height=\textheight, keepaspectratio]{Tred-2.pdf}
%     \vspace{-2cm}
%     \caption{$\psi_3(T_{m-1})$}
%     \end{subfigure}
%     \begin{subfigure}[t]{0.3\linewidth}
%     \includegraphics[width=\linewidth, height=\textheight, keepaspectratio]{Wgreen.pdf}
%     \vspace{-2cm}
%     \caption{$\psi_3(S_m)$}
%     \end{subfigure}
%     \begin{subfigure}[t]{0.3\linewidth}
%     \includegraphics[width=\linewidth, height=\textheight, keepaspectratio]{FigureTm.pdf}
%     \vspace{-2cm}
%     \caption{$\psi_3(T_{m})$}
%     \end{subfigure}
%     \vspace{.5cm}
% \caption{The clockwise  wrapping immersion $\psi_3$, in stages.} 
% \label{fig:wrap}
% \end{figure}

% We now observe that our choice of $\Psi_2$ implies that the lift of $$\psi_3 \colon \pi_z(P_{m+1}) \to D(2r)$$  to $P_{m+1}$ is a symplectic embedding. 


% \begin{lemma}
%     The  map $\Psi_3 (u,v,x,y) = (u,v,\psi_3(x,y)) $ is a symplectic embedding of $P_{m+1}$
%     into $W \times D(2r)$.
% \end{lemma}

% \begin{proof} It suffices to show that $\Psi_3$ is injective on $P_{m+1}$. Arguing by contradiction, assume that $(u,v,x_1,y_1)$ and $(u,v,x_2,y_2)$ are distinct points in $P_{m+1}$ whose images under $\Psi_3$ are the same. That is, \begin{align*}
%     \psi_3(x_1,y_1) = \psi_3(x_2,y_2).
% \end{align*} Since the restriction of  $\psi_3$ to each $S_i$ and $\bar{T}_j$ is an embedding, either we have $(x_1, y_1) \in \bar{T}_{j_1}$ and $(x_2, y_2) \in \bar{T}_{j_2}$ for some $-1\le j_1 \neq j_2 \leq m$, or else $(x_1, y_1) \in S_{i_1}$ and $(x_2, y_2) \in S_{i_2}$ for some $0\le i_1\neq i_2 \leq m$.
% In the first case, it follows from equation \eqref{eq:shift} that $(u,v)$ must be in both $D_{j_1+2}$ and  $D_{j_2+2}$. This is impossible since these rectangles are disjoint for $j_1 \neq j_2$. In the second case, equation \eqref{eq:shift} implies that $(u,v)$ must be in both $ \phi_g^{\chi(x_1 - 2i_1 +1)} (D_{i_1})$ and  $\phi_g^{\chi(x_2 - 2i_2 +1)} (D_{i_2})$. On the other hand, it follows from equation \eqref{eq:psiS} that if  $\psi_3(x_1,y_1)=\psi_3(x_2,y_2)$ then $x_1 -2j_1 = x_2 -2j_2$. Hence,  the sets $\phi_g^{\chi(x_1 - 2j_1 +1)} (D_{j_1 +1})$ and  $\phi_g^{\chi(x_2 - 2j_2 +1)} (D_{j_2 +1})$ are disjoint and we again arrive at a contradiction.
% \end{proof} 

% \medskip

% \noindent{\it The map $\Psi_4$.}  Since the area of $W$ may be assumed to be less than $b$, there is a Hamiltonian diffeomorphism $\psi_4$ of $\RR^2$ that maps $W$ into $D(b)$. Setting $\Psi_4(u,v,x,y) = (x,y,\psi_4(u,v)))$, we have $\Psi_4(\Psi_3(P_{m+1})) \subset P(2r,b).$


% \subsection{$\Phi$} To complete the proof of Proposition \ref{flex}, it remains to find a Hamiltonian diffeomorphism $\Phi$ of $P(2r,b)$
% such that the $\Phi(L_j)$ are all disjoint from $\Psi(\overline{P(r,s)})$. The map $\Phi$ will  be constructed as a composition of two Hamiltonian diffeomorphisms, $$\Phi = \Phi_2 \circ \Phi_1.$$

% \noindent{\it The map $\Phi_1$.}   The first map is of the form $\Phi_1= \mathrm{Id}_{\RR^2} \times \phi_1$ where $\phi_1$ is a Hamiltonian diffeomorphism of $D(b)$. The only important feature of $\phi_1$ will be the images of $\Lambda_1$ and $\Lambda_2$, with respect to $\psi_4(W)$. To simplify the notation, in the domain of $\phi_1$ we will identify each $\Lambda_j$  with it image under $\phi_1$ and we will identify $\psi_4(W)$ with $W$. 

% The image of $\Lambda_1$ is chosen so that it bounds a region containing the rectangle $$W \smallsetminus (D_1 \cup D_{m+2})$$ and intersects $W$ only at (a smoothing of) the union of the line segments $$\{(u,\eta) \mid 0 \le u\le 1\} \text{  and  } \{(1,v) \mid 0 \le v \le \eta\}$$ 
% and again at (a smoothing of) the union of 
% $$\{(u,\eta) \mid (m+1)(1+\varepsilon_2) \le u\le  (m+1)(1+\varepsilon_2)+1\}$$ and $$\{((m+1)(1+\varepsilon_2),v) \mid 0 \le v \le \eta\}.$$ 
% The region of $W$ bounded by $\Lambda_1$ has area arbitrarily close to $$mr + 2(r-\eta).$$
% We fix $\eta \in (r/2,r)$ so that this area is less than $s$.
% Hence $\Lambda_1$ can be chosen, as pictured in Figure \ref{fig:lambda1}, and still bound a total area equal to $s$.

% \begin{figure}
%     \centering
%     \vspace{-0cm}
%     \includegraphics[width=5.5in, height=7in]{Figure W1.pdf}
%     \vspace{-12cm}
%     \caption{The new position of $\Lambda_1$.}
%     \label{fig:lambda1}
% \end{figure}


% % It remains to find $E_n$ and $\Lambda_n$ such that the tori $L_n$ are disjoint from $P'$. Henceforth, to clarify the estimates we will ignore all terms of order $\varepsilon_2$.

% The only constraint placed on the image  of $\Lambda_2$ under $\phi_1$ is that it bounds almost all of the area of $W$ that is not bounded by $\Lambda_1$, as illustrated in Figure \ref{fig:lambda123}. 
% \begin{figure}
%     \centering
%     \vspace{-0cm}
%     \includegraphics[width=5.5in, height=7in]{Figure W12-2.pdf}
%     \vspace{-10.5cm}
%     \caption{The new position of $\Lambda_1$ and $\Lambda_2$.}
%     \label{fig:lambda123}
% \end{figure}
% With this, we may assume that the images of the $\Lambda_j$, with $j>2$, are all disjoint from  $W$.
% Hence, the repositioned Lagrangian tori $\Phi_1(L_3), \dots, \Phi_1(L_m)$ are all disjoint from $\Psi(\overline{P(r,s)})$. 



 
% \medskip

% \noindent{\it The map $\Phi_2$.} The second map, $\Phi_2$, will displace $\Phi_1(L_1)$ and $\Phi_1(L_2)$ from $\Psi(\overline{P(r,s)})$ while leaving the other repositioned Lagrangian tori fixed. Let $f(x,y)$ be a smooth nonnegative function on $\RR^2$ that equals one on $\Lambda_1 \cup \Lambda_2$ and vanishes outside an arbitrarily small neighborhood of $\Lambda_1 \cup \Lambda_2$. The map $\Phi_2$ will be the time one  Hamiltonian flow of a function of the form $$h(u,v,t)f(x,y).$$ By the properties of $f$, the only effect that  $\Phi_2$ can have on the Lagrangian tori $\Phi_1(L_j)$ is on the first factors of $\Phi_1(L_1) =S(r) \times \Lambda_1$ and  $\Phi_2(L_2) =S(r) \times \Lambda_2$ which will both get mapped to $\phi^1_h(S(r))$. It suffices to identify a closed curve $\Gamma =\phi^1_h(S(r))$ in $D(2r)$ such that $\Gamma$ bounds area $r$ and both  $\Gamma \times \Lambda_1$ and $\Gamma \times \Lambda_1$ are disjoint from $\Psi(\overline{P(r,s)})$. (The precise Hamiltonian $h$ whose flow maps  $S(r)$ to this $\Gamma$ will be irrelevant.)

% Consider the subset 
% $$X= \{2\eta <R, 0 \le\theta\le \theta_0\} \cup \left(\left\{\theta_0 \leq \theta \leq 1\right\} \smallsetminus \psi_3(T_{-1} \cup T_m)\right)$$
% of $D(2r)$ pictured in Figure \ref{fig:X}. In the $\varepsilon_j \to 0$ limit,  the area of $X$ approaches $$(r -\eta) + (r- \ell r ) = 2r -\eta- \ell c.$$ 
% Our choice of $\eta$ requires that $mr + 2(r- \eta)<s$. Hence 
% $$-\eta< \frac{s}{2} - \frac{mr}{2} -r.$$ We also have $$\ell c = \frac{s}{2} - \frac{mc}{2}.
% $$ Putting these together, the area of $X$ approaches $$r- \frac{m}{2}(r-c)$$ in the $\epsilon_j \to 0$ limit.


% SINCE THIS IS LESS THAN $r$ THE ORIGINAL ARGUMENT CAN NOT BE COMPLETED.




% \color{red}By \eqref{eq:delta} this limiting area is equal to $\frac{3r}{2} + \delta_1 - \delta_2.$ Hence, when the $\varepsilon_j$ are all sufficiently small we may assume that the area of $X$ is greater that $r$. Let $\Gamma$ be any smooth embedded circle in $X$ that bounds area $r$. It remains to show that, when the $\varepsilon_j$ are all sufficiently small, the Lagrangian tori $\Gamma \times \Lambda_1$ and $ \Gamma \times \Lambda_2$ are both disjoint from $\Psi_4 (\Psi_3 (\Psi_2(P_0)))$ and hence from $ \Psi(\overline{P(r,s)})$. 



% \begin{figure}
%     \centering
%     \vspace{-0cm}
%     \includegraphics[width=5.5in, height=7in]{Figure x.pdf}
%     \vspace{-6cm}
%     \caption{The region $X \subset D(2r)$ is pictured in white. The gray box with arms corresponds to $\psi_3(T_{-1}) \cup \psi_3(T_m)$.}
%     \label{fig:X}
% \end{figure}

% Arguing by contradiction,  assume that for $p = (u,v,x,y)$ in $P_0$ we have 
% \begin{equation}\label{eq:assump}\Psi_3 (\Psi_2(p))\in ( \Lambda_1 \times \Gamma \cup  \Lambda_2\times \Gamma).\end{equation}
% The number $x$ takes values in the interval $[-1-\ell, 2m+l]$ and either belongs to a subinterval of the form $$[-1+2k + \varepsilon_4/4,2k-\varepsilon_4/2],$$ for some $0\leq k \leq m $,  or it does not. 

% In the latter case, $\pi_z(\Psi_2(p))$ lies in some $\bar{T}_j$. Hence, $\psi_3 ( \pi_z(\Psi_2(p)))$ is contained in $\{0<R<2r, \theta_1 \leq  \theta \leq 1\}.$ In fact, it follows from the intersection assumption \eqref{eq:assump} and the choice of $\Gamma$ that $\psi_3 ( \pi_z(\Psi_2(p)))$ is  contained in $\psi_3(T_{-1} \cup T_m)$. 



% $\pi_w(\Psi_2(p))$ lies in $W$. By the construction of $\Lambda_1$ and $\Lambda_1$


% The crucial observation is the following

% \begin{lemma}
%     Let $p=(u,v,x,y) \in P_0 =D_1 \times Z$ be a point such that $v  < \lambda \le r(1+\varepsilon_1)$, and $x \in [-1+2k + \varepsilon_4/4,2k-\varepsilon_4/2]$ for some $0\leq k \leq m$. Then for any $\delta> 0$ we have 
% \begin{equation}
%     |\psi_3(\pi_z(\Psi_2(u,v,x,y)))| \leq 2
%     \lambda +\delta
% \end{equation}
% whenever  $\varepsilon_1, \dots, \varepsilon_6$ are sufficiently small.
% \end{lemma}

% \begin{proof}
%     Since $x$ is in $[-1+2k + \varepsilon_4/4,2k-\varepsilon_4/2]$, we have $y \in [0,\varepsilon_3]$ and $$\pi_z(\Psi_2(u,v,x,y))) = (x,y + (1+\varepsilon_2)v \chi'(x+1 -2k)) \in S_k.$$ The explicit formula for $\psi_3$ on $S_k$ then yields
%     $$
%     |\psi_3(\pi_z(\Psi_2(u,v,x,y)))| = \frac{2}{\mathcal{E}(1+\varepsilon_5)} (y + (1+\varepsilon_2)v \chi'(x+1 -2k))
%     $$
% and hence 
%  $$
%     |\psi_3(\pi_z(\Psi_2(u,v,x,y)))| < 2 (\varepsilon_3 + v(1+\varepsilon_2)(1+\varepsilon_4)).
%     $$
% The result follows easily from this estimate.
% \end{proof}













% The $\Lambda_n$ are disjoint and so the only requirement on the $E_n$ to define disjoint tori $L_n = E_n \times \Lambda_n$ is that they bound disks of area $r$. Also, for $n \ge 3$, the $L_n$ will be disjoint from $P'$ independent of the choice of $E_n$.

% The torus $L_1$ will be disjoint from $P'$ provided $E_1$ lies in the set
% $$U = \{ \theta \in [0, \frac 12], \, R > 2t \} \cup  \left( \{ \theta \in [\frac 12, 1] \} \smallsetminus \psi(T_{-1} \cup T_m) \right).$$
% This is because points in $P$ with $2i-1 < x < 2i$ and $v < t$ have $y < \chi'(x - 2i +1)t \le t$. Also, as $\Lambda_1 \cap W \subset D_1 \cup D_{m+1}$, it intersects the fibers in $P$ over $T_i$ only if $i =-1$ or $i=m$. We compute the area of $U$ as $(r-t) + (r - lr) = 3r/2 > r$. 


% Therefore we can find $E_1$ as required, see Figure 3 again.

% Similarly, $L_2$ will be disjoint from $K'$ provided $E_2$ avoids $U$. Therefore we can simply set $E_2 = E_1$. The remainder of the $E_n$ can be chosen arbitrarily and we have constructed our $L_n$ as required.
% \end{proof}
 

% % \begin{proof}

% % All collections of $k$ circles in $D(y)$ are symplectomorphic, therefore it suffices to find a single family of $\Lambda_n$ such that the corresponding $L_n$ are disjoint from the image of $P(r,s)$. Also, we can replace $L_n$ by $E_n \times \Lambda_n$, where the $E_n$ are circles in $D(2r)$ bounding area $r$, and the Hamiltonian isotopy type of the union of the $L_n$ remains the same.

% % \medskip

% % {\it 1. Preparation.}

% % Let $m \in \NN$ with $mr < s < (m+1)r$. Then we choose $0<c<r$ and $\varepsilon_2 >0$ with $\varepsilon_2$ and $r-c$ small. Let $2lc = s - mc$.

% % First, in the $z = x + iy$ plane, we identify $D(s)$ with a small neighborhood $Z$ of the set
% % $$([-1 - l, 2m + l] \times \{ 0 \}) \bigcup_{i=0}^{m-1} [2i, 2i+1] \times [0, c] \bigcup ([-1-l, -1] \times [0,c]) \bigcup ([2m, 2m+l] \times [0,c]).$$
% % It will be convenient to set $T_{-1} = [-1-l, -1] \times [0,c]$ and $T_m = [2m, 2m+l] \times [0,c]$ and $T_i = [2i, 2i+1] \times [0, c]$ for $0 \le i \le m-1$, see Figure 1.

% % In the $w = u + iv$ plane we define disjoint disks $D_j$ of area $r$ for $1 \le j \le m+1$ by
% % $$D_j = [(1 + \varepsilon_2)(j-1) , (1 + \varepsilon_2)j - \varepsilon_2] \times [0, r].$$ Let $W = [0, (1+ \varepsilon_2)(m+1) - \varepsilon_2] \times [0,r]$, see Figure 2.

% % We note that the Hamiltonian function $G(v) = (1+ \varepsilon_2) v$ generates a Hamiltonian flow $\phi_G^t$ with $\phi_G^1(D_j) = D_{j+1}$. Further $\phi_G^t(D_j) \subset W$ for all $j \le m$ and $0 \le t \le 1$.

% % Next we fix a function $\chi : \RR \to [0,1]$ with $\chi(x) =0$ when $x \le 0$ and $\chi(x) =1$ when $x \ge 1$ such that $0 \le \chi'(x) \le 1+\varepsilon_1$ for all $x$.

% % Let $P_0 = Z \times D_1$, which is Hamiltonian diffeomorphic to $P(r,s)$.

% % \medskip

% % {\it 2. The polydisk embedding $\Psi$.}

% % We define $P_1 = \phi_1(P_0)$, where $\phi_1$ is generated by the Hamiltonian function $$H_1 = \chi(x+1) G(w).$$

% % Then $$X_{H_1} = \chi'(x+1)G(w) \partial_y + \chi(x+1)(1+\varepsilon_2) \partial_u$$ and letting $\pi : \CC^2 \to \CC$ denote the projection onto the $z$ factor we have
% % $$\pi(P_1) \subset Z \bigcup [-1,0] \times [ 0,  (1+\varepsilon_2)r ].$$
% % The fibers satisfy $\pi^{-1}(x,y) = D_1$ when $x \le -1$ and $\pi^{-1}(x,y) = D_2$ when $x \ge 0$, and in general  $\pi^{-1}(x,y) \subset \phi_G^{\chi(x+1)} D_1$.

% % Carrying on in this way we define $P_2, \dots, P_{m+1}$ where $P_j = \phi_j( P_{j-1})$ and $\phi_j$ is generated by a Hamiltonian function $$H_j = \chi(x - 2j + 3) G(w).$$

% % Let $P = P_{m+1}$. We also denote by $S_i$ the square $[2i-1, 2i] \times  [ 0, (1+\varepsilon_2)r]$ for $0 \le i \le m$. Then $$\pi(P) \subset Z \bigcup_{i=0}^m S_i,$$ see again Figure 1.
% % For $2i \le x \le 2i+1$ we have $\pi^{-1}(x,y) = D_{i+2}$ and when $2i-1 \le x \le 2i$ we have $\pi^{-1}(x,y) \subset \phi_G^{ \chi(x - 2i +1)} D_{i+1}$.

% % We will define an immersion $\psi : \pi(P) \to D(2r)$ using polar coordinates $R = \pi |z|^2$, $\theta \in \RR / \ZZ$ on the disk. To do this, when $(x,y) \in S_i$ we set $$\phi(x,y) = (R = \frac{2y}{1+\varepsilon_2}, \theta = \frac{(1+\varepsilon_2) (x+1 - 2i)}{2} )$$
% % and map the remainder of $\pi(P)$ to $D(2r) \smallsetminus \{ \theta \in [0, (1 + \varepsilon_2)/2] \}$ such that we have a symplectic embedding on each square $T_i$. This is possible since $c < 2r - (1+ \varepsilon_2)r = (1 - \varepsilon_2)r$.
% % The immersion $\psi$ is illustrated in Figure 3. Note we are assuming $\psi(T_{-1})$ and $\psi(T_m)$ have roughly the same image, which has area $lc$.

% % {\bf Claim.} The map $\Psi = \psi \times \mathrm{I}_W$ is an embedding of $P$.

% % {\it Proof of Claim.} Suppose $\psi(x_1, y_1) = \psi(x_2, y_2)$. Then either $(x_1, y_1) \in T_{i_1}$ and $(x_2, y_2) \in T_{i_2}$ or else $(x_1, y_1) \in S_{i_1}$ and $(x_2, y_2) \in S_{i_2}$.
% % Also, if $(x_1, y_1)$ and $(x_2, y_2)$ are different we must have $i_1 \neq i_2$.
% % In the first case the fibers (inside $P$) are given by $\pi^{-1}(x_1, y_1) = D_{i_1+2}$ and  $\pi^{-1}(x_2, y_2) = D_{i_2+2}$. Hence there are no corresponding double points of $\Psi$. In the second case, by the formula for $\psi|_{S_i}$ we see that $x_1 - 2i_1 = x_2 - 2i_2$. But then the fibers   $\pi^{-1}(x_1, y_1) \subset \phi_G^{\chi(x_1 - 2i_1 +1)} D_{i_1 +1}$ and  $\pi^{-1}(x_2, y_2) \subset \phi_G^{\chi(x_2 - 2i_2 +1)} D_{i_2 +1}$ are disjoint. $\Box$

% % After a symplectomorphism of the $w$ plane we may assume the image of $\Psi$ lies in $P(2r, y )$, then $P' = \Psi(P)$ is our image of $P(r,s)$.

% % \medskip

% % {\it 3. Lagrangian tori disjoint from $P'$.}

% % It remains to find $E_n$ and $\Lambda_n$ such that the tori $L_n$ are disjoint from $P'$. Henceforth, to clarify the estimates we will ignore all terms of order $\varepsilon_2$.

% % First we define the $\Lambda_n$. Since $(m-1)r + (1+2l)r = s$ (taking $r - c$ of order $\varepsilon_2$, so it's negligible) we can choose $\Lambda_1$ such that it bounds a disk containing $\cup_{j=2}^m D_j$.  Furthermore, we may assume $\Lambda_1$ intersects $W$ only at (a smoothing of) $$\{ u \in [0,1] \cup [m, m+1], \, v = t\} \cup \{ u \in \{1, m\}, \,  v \in [0,t] \}.$$
% % Figure 4 is an illustration of $\Lambda_1$.
% % An area calculation gives $2(r-t) + (m-1)r = s$, or $r-t = r(l + 1/2)$.

% % The remaining $\Lambda_n$ can be chosen so that $\Lambda_2$ intersects $W$ only in the region $\{ u \in [0,1] \cup [m, m+1], \, v < t\}$ and the $\Lambda_n$ for $n \ge 3$ are disjoint from $W$, see again Figure 4.


% % The $\Lambda_n$ are disjoint and so the only requirement on the $E_n$ to define disjoint tori $L_n = E_n \times \Lambda_n$ is that they bound disks of area $r$. Also, for $n \ge 3$, the $L_n$ will be disjoint from $P'$ independent of the choice of $E_n$.

% % The torus $L_1$ will be disjoint from $P'$ provided $E_1$ lies in the set
% % $$U = \{ \theta \in [0, \frac 12], \, R > 2t \} \cup  \left( \{ \theta \in [\frac 12, 1] \} \smallsetminus \psi(T_{-1} \cup T_m) \right).$$
% % This is because points in $P$ with $2i-1 < x < 2i$ and $v < t$ have $y < \chi'(x - 2i +1)t \le t$. Also, as $\Lambda_1 \cap W \subset D_1 \cup D_{m+1}$, it intersects the fibers in $P$ over $T_i$ only if $i =-1$ or $i=m$. We compute the area of $U$ as $(r-t) + (r - lr) = 3r/2 > r$. Therefore we can find $E_1$ as required, see Figure 3 again.

% % Similarly, $L_2$ will be disjoint from $K'$ provided $E_2$ avoids $U$. Therefore we can simply set $E_2 = E_1$. The remainder of the $E_n$ can be chosen arbitrarily and we have constructed our $L_n$ as required.
% % \end{proof}
% \begin{section}{Appendix: Proof of Corollary \ref{>r}}\label{appen}
    
\clearpage

\begin{thebibliography} {99}

\bibitem{atallah} M. Atallah, J-P. Chass\'{e}, R. Leclercq and E. Shelukhin, Weinstein exactness of nearby Lagrangians and related questions, arXiv:2410.04158.

%\bibitem{bo}F. Bourgeois,
%A Morse-Bott approach to contact homology, PhD Thesis, Stanford University, 2002.

%\bibitem{B}
%F. Bourgeois, A Morse-Bott approach to contact homology. Symplectic and contact topology: interactions and perspectives (Toronto, ON/Montreal, QC, 2001), 55--77, Fields Inst. Commun., 35, Amer. Math. Soc., Providence, RI, 2003.

\bibitem{behwz}
F. Bourgeois, Y. Eliashberg, H. Hofer, K. Wysocki, and E. Zehnder,
 Compactness results in symplectic field theory,
{\em Geom. Topol.}, 7 (2003), 799--888.

%\bibitem{bormon} F. Bourgeois and K. Mohnke, Coherent orientations in symplectic field theory, {\em Math. Z.}, 248 (2004), no. 1, 123--146.

\bibitem{boust} K. Boustany, Packing integral tori in Del Pezzo surfaces, {\em J. Symplectic Geometry}, 
Volume 22 (2024) Number 6, 1267--1348.


\bibitem{brsh} F. Bro\'ci\'c and E. Shelukhin, A counterexample to Lagrangian Poincar\'e recurrence, arXiv:2409.14225.

\bibitem{brki} J. Brendel and J. Kim, Lagrangian split tori in $S^2 \times S^2$ and billiards,  arXiv:2502.03324, to appear in {\em Selecta Math.}.



% \bibitem{bu} O. Buse, negative inflation and stability in symplectomorphism groups of ruled surfaces, {\em J. Symplectic Geometry}, {\bf 9} (2011), 147--160.

%\bibitem{busehind} O. Buse and R. Hind, Symplectic embedding of ellipsoids in dimension greater than four, {\em Geom. Top.}, 15 (2011), 2091--2110.

% \bibitem{ch} Y. Chekanov Hofer?s Symplectic Energy and Lagrangian intersections,
% in Contact and Symplectic Geometry, Cambridge, 1996, 296--306.

\bibitem{chek} Y. Chekanov, Lagrangian tori in a symplectic vector space and global symplectomorphisms, 
{\em Math. Z.} 223 (1996), 547--559.

\bibitem{chsc} Y. Chekanov and F. Schlenk, Lagrangian product tori in tame symplectic manifolds, {\em Comment. Math. Helv.} 91 (2016), 445--475.


% \bibitem{cm} K. Cieliebak and K. Mohnke, Punctured holomorphic curves and Lagrangian embeddings, {\em Inventiones Mathematicae}, 212 (2018), 213--295.

% \bibitem{cisc} K. Cieliebak and M. Schwingenheuer, Hamiltonian unknottedness of certain monotone Lagrangian tori in $S^2\times S^2$, {\em Pacific Journal of Mathematics}, {\bf 299} (2019 ), 427--468. 



%\bibitem{choi} K. Choi, D. Cristofaro-Gardiner, D. Frenkel, M. Hutchings and V. G. B. Ramos, Symplectic embeddings into four-dimensional concave toric domains, {J. Topol.}, to appear.

%\bibitem{dr} D. Dragnev,
%Fredholm theory and transversality for noncompact pseudoholomorphic maps in symplectizations, \emph{Comm. Pure Appl. Math.}, \textbf{57} (2004), 726--763.

%\bibitem{eh1} I. Ekeland and H. Hofer, Symplectic topology and
%Hamiltonian dynamics, {\it Math. Z.}, \textbf{200} (1989), 355--378.

%\bibitem{ekehof} I. Ekeland and H. Hofer, Symplectic topology and
%Hamiltonian dynamics II, {\it Math. Z.}, 203 (1990), 553--567.

%\bibitem{evle} J. Evans and Y. Lekili,  Generating the Fukaya categories of Hamiltonian $ G$-manifolds, {\em J. Amer. Math. Soc.} \textbf{32} (2019), 119-162 .



%\bibitem{EGH} Y. Eliashberg, A. Givental and H. Hofer, Introduction to symplectic field theory, GAFA 2000 (Tel Aviv, 1999), {\em Geom. Funct. Anal.}, 2000, Special Volume, Part II, 560--673.

% \bibitem{ep} M. Entov, L. Polterovich, Quasi-states and symplectic intersections, {\em Comm. Math. Helv.} \textbf{81:1} (2006), 75-99.

%\bibitem{flho} A. Floer, H. Hofer and K. Wysocki, Applications of symplectic homology. I. {\em Math. Z.}, 217 (1994), 577--606.

% \bibitem{fooo} K. Fukaya, Y-G. Oh, H. Ohta and K. Ono, Lagrangian Floer theory on compact toric manifolds. I. {\it Duke Math. J.}, \textbf{151} (2010),  23-174.

%\bibitem{frmu} D. Frenkel and D. M\"{u}ller, Symplectic embeddings of 4-dimensional ellipsoids into cubes, arXiv:1210.2266.

\bibitem{rgi} G. Dimitroglou-Rizell, E. Goodman and A. Ivrii, Lagrangian isotopy of tori in $S^2 \times S^2$ and $\CC P^2$, {\it Geometric and Functional Analysis} 26 (2016), 1297--1358.

\bibitem{gr} M. Gromov, Pseudo-holomorphic curves in symplectic manifolds, {\em Inventiones Mathematicae}, 82 (1985), 307--347.

\bibitem{hicksmak} J. Hicks and C. Y. Mak, Some cute applications of Lagrangian cobordisms towards examples in quantitative symplectic geometry, arXiv:2208.14498.

%\bibitem{guth} L. Guth, Symplectic embeddings of polydisks, {\it Inv.  Math}, \textbf{172} (2008), 477--489.


%\bibitem{harris} J. Harris and I. Morrison, Moduli of Curves,
%Springer 1998.

%\bibitem{hind} R. Hind, Symplectic folding and nonisotopic polydisks, {\it Algebr. Geom. Topol.}, 13 (2013), 2171--
%2192.

%\bibitem{hindivrii} R. Hind and A. Ivrii, Ruled 4-manifolds and isotopies of symplectic surfaces, {\it Math. Z.},
 %265, (2010), 639--652.

%\bibitem{hindker} R. Hind and E. Kerman, New obstructions to symplectic embeddings, {\it Inv. Math.}, 196 (2014), 383--452.

% \bibitem{hl} R. Hind and S. Lisi, Symplectic embeddings of polydisks, {\it Selecta Math.}, 21 (2015), 1099--1120.

%\bibitem{hl-err} R. Hind and S. Lisi, Erratum to: Symplectic embeddings of polydisks, {\it Selecta Math.}, 23 (2017), 813--815.

\bibitem{hiker} R. Hind and E. Kerman, Packing Lagrangian tori, {\it Geom. and Topol.} \textbf{28}:5 (2024), 2207-2257.

\bibitem{ho} R. Hind and E. Opshtein, Squeezing Lagrangian tori in dimension $4$, {\it Comment. Math. Helv.}, 95 (2020), 535--567.

\bibitem{hizh} R. Hind and J. Zhang, The shape invariant of symplectic ellipsoids, arXiv:2010.02185.

\bibitem{hizh2} R. Hind and J. Zhang, Hamiltonian knottedness and lifting paths from the shape invariant, {\it Compos. Math.}, 159 (2023), 2416--2457.

%\bibitem{hofer} H. Hofer, Symplectic capacities. In: Geometry of Low-dimensional Manifolds 2
%(Durham, 1989). Lond. Math. Soc. Lect. Note Ser., vol. 151, pp. 15–34. Cambridge
%Univ. Press, Cambridge (1990).

%\bibitem{hls} H. Hofer, V. Lizan and J.-C. Sikorav, On genericity of holomorphic curves in $4$-dimensional almost-complex manifolds, {\it J. Geom. Anal.}, \textbf{7} (1997), 149--159.

%\bibitem{three} H. Hofer, K. Wysocki and E. Zehnder, A characterisation of the tight three-sphere, {\it Duke Math. J.}, 81 (1995), 159--226.

% \bibitem{hofa} H. Hofer, K. Wysocki and E. Zehnder, Properties of pseudoholomorphic curves in symplectisations I: Asymptotics, {\it Ann. Inst. H. Poincar\'{e} Anal. Non Lineaire}, 13 (1996) , 337--379.

%\bibitem{hofi} H. Hofer, K. Wysocki and E. Zehnder, Properties of pseudoholomorphic curves in symplectisations II: Embedding controls and algebraic invariants, {\it Geom. Funct. Anal.}, 5 (1995),  337--379.

%\bibitem{hoff} H. Hofer, K. Wysocki and E. Zehnder, Properties of pseudoholomorphic curves in symplectisations III: Fredholm theory, {\it Topics in nonlinear analysis}, 381-475, Prog. Nonlinear Differential Equations Appl., 35, Birkh\"{a}user, Basel, 1999.

%\bibitem{hoferzehnder} H. Hofer and E. Zehnder, Symplectic invariants and Hamiltonian dynamics,  Birkh\"{a}user, Basel, 1994.

%\bibitem{hwzfol} H. Hofer, K. Wysocki and E. Zehnder, Finite energy foliations of tight three-spheres and Hamiltonian
 % dynamics, {\it Ann. of Math. (2)}, 157 (2003), 125--255.

\bibitem{hutch-talk} M. Hutchings, talk {\em Classification of Some Open Toric Domains},  Joint IAS/Princeton/Montreal/Paris/Tel-Aviv Symplectic Geometry Zoominar, November 1 2024, https://www.youtube.com/watch?v=JGF20oz2CsQ.

%\bibitem{hutchings1} M. Hutchings, An index inequality for embedded pseudoholomorphic curves in symplectizations, {\it J. Eur. Math. Soc. (JEMS)},  4  (2002),  no. 4, 313--361.

%\bibitem{hutchings2} M. Hutchings, The embedded contact homology index revisited, {\it New perspectives and challenges in symplectic field theory},  263--297, CRM Proc. Lecture Notes, 49, Amer. Math. Soc., Providence, RI, 2009.

%\bibitem{hutchings} M. Hutchings, Quantitative embedded contact homology, {\it J. Diffl. Geom.} 88 (2011), 231--266.

%\bibitem{hutchings2} M. Hutchings, Recent progress on symplectic embedding problems in four dimensions, {\it Proc. Natl. Acad. Sci. USA}, 108 (2011), 8093--8099.

%\bibitem{hutchingsl} M. Hutchings, Lecture notes on embedded contact homology.

%\bibitem{hutsul} M. Hutchings and M. Sullivan, Rounding corners of polygons and the embedded contact homology of $T^3$, {\it Geom. Topol.}, \textbf{10} (2006), 169–-266.

\bibitem{ivriit} A. Ivrii, {\it Lagrangian unknottedness of tori in certain symplectic 4-manifolds}, PhD thesis, Stanford University, 2003.

% \bibitem{lm96} F. Lalonde and D. McDuff, J-curves and the classification of rational and ruled symplectic 4-manifolds, {\it Contact and Symplectic Geometry}, Cambridge 1994, 3--42. 

%\bibitem{lalmcd} F. Lalonde and D. McDuff, The geometry of symplectic energy, {\it Ann. of Math.}, 141 (1995), 349--371.

% \bibitem{maksmi} C. Y. Mak and I. Smith, Non-displaceable Lagrangian links in four manifolds, {\it Geometric and Functional Analysis}, 31 (2021), 438--481.

%\bibitem{mcduffblow} D. McDuff, Blowing up and symplectic embeddings in dimension $4$, {\it Topology}, \textbf{30} (1991), 409--421.

\bibitem{mcdrat} D. McDuff, The structure of rational and ruled symplectic 4-manifolds,
{\it J. Amer. Math. Soc.}, \textbf{3} (1990), 679--712.

% \bibitem{mcd01} D. McDuff, Symplectomorphism groups and almost-complex structures, {\it Enseignement Math}, (2001), 1--30.

% \bibitem{mcp} D. McDuff and L. Polterovich, Symplectic packings and algebraic geometry, {\it Inventiones Math.}, \textbf{115} (1994), 405--429.

% \bibitem{mcs} D. McDuff and D. Salamon, $J$-holomorphic curves and
% symplectic topology. American Mathematical Society Colloquium
% Publications, 52. American Mathematical Society, Providence, RI,
% 2004.

%\bibitem{mcdsch} D. McDuff and F. Schlenk, The embedding capacity of $4$-dimensional symplectic ellipsoids, {\it Ann. of Math.}, 175 (2012), 1191--1282.

%\bibitem{mcduff} D. McDuff, Symplectic embeddings of $4$-dimensional ellipsoids, {\it J. Topol.}, 2 (2009), 1--22.

%\bibitem{mcd2} D. McDuff,
%The Hofer conjecture on embedding symplectic ellipsoids, {\it J. Diffl. Geom.}, 88 (2011), 519-–532.

%\bibitem{mcdops} D. McDuff and E. Opshtein, Nongeneric $J$-holomorphic curves and singular inflation, preprint, arXiv:1309.6425.

\bibitem{ohzh} Y.-G. Oh and K. Zhu,  Embedding property of somewhere injective
J-holomorphic curves in Calabi-Yau manifolds for generic $J$, {\it Asian J. Math.}, 13
(2009), 323--340.

%\bibitem{opshtein} E. Opshtein, Maximal symplectic packings of $\PP^2$, {\it Compos. Math.}, 143(2007), 1558--1575.

\bibitem{ps} L. Polterovich and E. Shelukhin, Lagrangian configurations and Hamiltonian maps, {\em Compositio Math.} \textbf{159} (2023), 2483?2520.

% \bibitem{oh} Y.-G. Oh, Addendum to Floer Cohomology of Lagrangian Intersections and Pseudo-Holomorphic Discs, I, {\it Comm. Pure Appl. Math.}, {\bf48} (1995), 1299--1302. 

%\bibitem{pelngo1} A. Pelayo and S. V. Ng\d{o}c, The Hofer question on intermediate symplectic capacities, preprint, arXiv:1210.1537.

%\bibitem{pelngo2} A. Pelayo and S. V. Ng\d{o}c, Sharp symplectic embeddings of cylinders, preprint, arXiv:1304.5250.



% \bibitem{rs} A.F. Ritter, I. Smith, The monotone wrapped Fukaya category and the open-closed string map,{\it Sel. Math. New Ser.} {\bf 23}  (2017), 533--642. 

%\bibitem{rizzel} G. Rizzel, Talk at joint AMS-EMS conference, Porto, June 2015.


%\bibitem{rs} J. Robbin and D. Salamon, The Maslov index for paths, {\it Topology}, 32 (1993), 827--844.


%\bibitem{schl} F. Schlenk, Embedding problems in symplectic geometry
%De Gruyter Expositions in Mathematics 40. Walter de Gruyter Verlag, Berlin. 2005.

% \bibitem{schl-er} F. Schlenk, An extension theorem in symplectic geometry, {\it Manuscripta Math.} {\bf 109}, 329?348 (2002). 


%\bibitem{schwarz} M. Schwarz, Cohomology operations from $S^1$-cobordisms in Floer homology, PhD thesis, ETH Z\"{u}rich, 1995.

\bibitem{schmitz} J. Schmitz, Nodal tangles, arXiv:2506.23754.


%\bibitem{shev} V. Shevchishin, Pseudoholomorphic curves and the
%symplectic isotopy problem, preprint math.SG/0010262.

%\bibitem{diam} J.-C. Sikorav, Some properties of holomorphic curves in almost-complex manifolds, in the book {\it Holomorphic curves in symplectic geometry}, 165--189, Progr. Math., 117, Birkh\"{a}user, Basel, 1994.

%\bibitem{tr} L. Traynor, Symplectic packing constructions, \emph{J. Diff. Geom.}, \textbf{42} (1995), 411--429.

% \bibitem{vianna} R. Vianna, Infinitely many monotone Lagrangian tori in del Pezzo surfaces, {\em Selecta Mathematica}, (2016).

%\bibitem{vit} C. Viterbo, A new obstruction to embedding Lagrangian tori,
%{\em Inventiones mathematicae}, {\bf 100} (1990) 301--320.

\bibitem{we} C. Wendl, Automatic transversality and orbifolds
of punctured holomorphic curves in dimension four,  {\it Comment. Math. Helv.}, 85 (2010), 347--407.



\bibitem{ze} K. Zehmisch, Holomorphic jets in symplectic manifolds, {\it J. Fixed Point
Theory Appl.}, 17 (2015), 379--402.

\end{thebibliography}

\end{document}









